% 本文件由 LaTeX 排版 Agent 根据有效 translation.json 自动生成。
% author=Jerome; work_id=2710; title=驳鲁菲诺护教篇

\chapter{\WorkTitle{驳鲁菲诺护教篇}}

% structure: p0001 source=h1 target=omitted reason=page-title-already-rendered-by-parent

卷一\newline{}卷二\newline{}卷三\par

\SourceRecord{Translated by W.H. Fremantle.；Nicene and Post-Nicene Fathers, Second Series；Vol. 3.；Edited by Philip Schaff and Henry Wace.；Buffalo, NY: Christian Literature Publishing Co.,；1892.；<http://www.newadvent.org/fathers/2710.htm>.}

% structure: p0001 source=h1 target=section reason=page-title

\section{驳斥鲁菲努辩护书（第一卷）}

\Emph{致\Person{帕玛基}和\Person{玛尔策拉}，自\Place{白冷}，公元402年。}\par

\EditorialNote{热罗尼莫在撰写其辩护书时所掌握的文献包括：（1）鲁菲努翻译的\WorkTitle{庞斐禄辩护书}及其所附前言，以及关于篡改奥力振著作的书籍；（2）\OriginalTerm{论原理}{\GreekText{Περὶ} \GreekText{᾿Αρχῶν}}的翻译和鲁菲努的前言（见第430页）；（3）鲁菲努致阿纳斯塔西的辩护书（见第430页）；（4）阿纳斯塔西致耶路撒冷若望的书信（见第432页《辩护书》卷二第14章、卷三第20章）。他还拥有阿纳斯塔西的其他书信，例如致米兰主教的那一封（热罗尼莫《书信》第95封，另见《辩护书》卷三第21章）。但他未获得鲁菲努辩护书全文（第4章、第15章）。402年初春，他收到帕玛基和玛尔策拉的书信，当时鲁菲努于400年底在阿奎莱亚写成的辩护书已在其友人中流传了一段时间。他们未能获得完整副本，但寄来了主要论点，并强烈敦促热罗尼莫答复。与此同时，他的兄弟保利尼亚诺已在西方居留约三年，经由罗马返回巴勒斯坦，在那里听到并看到鲁菲努辩护书的部分内容，将其铭记于心（《辩护书》卷一第21章、第28章），并在白冷复述。热罗尼莫对这些文献作了答复。}\par

% structure: p0004 source=h2 target=subsection reason=relative-heading-rank; base=h2

\subsection{难哉！昔日之友，既已和好，竟以秘密传阅之书攻击我。}

我不仅从你的来信，也从许多人的来信中得知，在\Person{提辣诺}的学派里，人们用\Quote{我的狗从敌人自己的舌头}攻击我，因我将《\OriginalTerm{论原理}{\GreekText{Περὶ} \GreekText{᾿Αρχῶν}}》一书译成了拉丁文。这是何等空前的无耻！他们指责医生发现了毒药，目的却是保护贩卖毒品的人，不是得着无罪的赏报，而是与罪犯同伙；仿佛犯罪人数多能减轻罪行，或者指控取决于我们的个人情感而非事实。有人撰写小册子攻击我，强加于所有人的注意，却不敢公开出版，以便扰乱单纯之人的心，不给我答辩的机会。这是一种伤害人的新方法：提出你害怕传播的指控，写下你不得不隐藏的东西。如果他写的属实，为何惧怕公众？倘若虚假，为何又写下？我们幼时读过\Person{西塞罗}的话：\Quote{我认为，写任何你想隐藏的东西，都是缺乏自制。}我问，他们在抱怨什么？这热狂与愤怒从何而来？是否因为我拒绝了虚假的赞美？因为我拒绝了用不诚实的言语献上的颂辞？因为我以朋友之名察觉了敌人的陷阱？在这篇序言中，我被称作兄弟与同僚，然而我的所谓罪行却被公开宣扬，说我撰文支持奥力振，用赞美将他捧上了天。作者声称他这样做是出于好意。那么，为何他现在像公开的敌人一样，将从前作为朋友赞美的事抛在我的脸上？他宣称，他本想效法我这位先行者的翻译，并从我的一些拙作中为其作品借用权威。若果真如此，他只需声明我曾写作就够了。他何须重复同一件事，通过反复强加于人注意，再再翻来覆去地说同样的话，仿佛没人会相信他的赞美？简单真诚的赞美不会如此焦虑地追求读者的信任。为何他害怕，除非他拿出我的原话作证，否则没人相信他赞美我？你看，我们完全理解他的伎俩；他显然上过戏校，通过不断练习学会了嘲讽式的赞美者的角色。假装单纯是无用的，因为阴谋家的恶意已被识破。犯一次错，或放宽说两次，也许是不幸的偶然；但他为何明知故犯，重复错误，并将这错误编织进所有著作，使我无法否认他赞美我的事？一个真正明白事理的朋友，在我们先前的误解与和解之后，本应避免一切可疑的举止，并留意不要因疏忽做出看似蓄意的事。图利乌斯在为其加里纽斯辩护的书中说：\Quote{我始终认为，保守我所建立的一切友谊是一种最高的宗教责任，尤其是那些经过分歧之后善意得以恢复的友谊。对于从未动摇过的友谊，若有所疏忽，健忘的借口，或最坏也是疏忽，容易被接受；但重归于好之后，若有得罪之举，则不被视为疏忽，而是不友善的意图，这不再是轻率，而是背信。}同样，\Person{贺拉斯}在致\Person{弗洛鲁斯}的信中写道：\par

\Quote{善意，若未紧密结合，便不会粘连而会飞散。}\par

% structure: p0007 source=h2 target=subsection reason=relative-heading-rank; base=h2

\subsection{别人也翻译了奥力振，为何偏偏指责我？}

2. 他发誓说他是出于单纯而犯错，这对我有什么好处呢？如你所知，他的称赞被当面数落，而这位最单纯的朋友的赞誉（但这赞誉中并没有多少单纯或真诚）却被归咎于我，视为罪行。如果他在为自己所做的事寻求权威的基础，并希望表明谁在他之前走过这条路，他手边有宣信者\Person{Hilary}，他翻译了\Person{Origen}关于\WorkTitle{约伯传}和\WorkTitle{圣咏集}的著作，共计四万行。他还有\Person{Ambrose}，其著作几乎全部充满\Person{Origen}所写的内容；还有殉道者\Person{Victorinus}，他真正行事\Quote{单纯}，不设陷阱陷害他人。对于所有这些，他保持沉默；他不理会那些如同教会柱石的人；而我，不过像一只跳蚤，一个微不足道的人，他却从各个角落搜寻出来。也许正是同样的单纯，使他意识不到自己在攻击朋友，也将使他发誓他完全不知道这些作者。但谁会相信他不认识这些人呢？他们尚在人们记忆中，且是拉丁人，而他本人又是如此博学，对古代作家，尤其是希腊作家有如此丰富的知识，以致于他对异邦学问的热忱几乎使他丧失了自己的语言？事实上，与其说我被他赞扬，不如说那些作者未被攻击。但无论他所写的是称赞（如他试图让头脑简单的人相信的那样）还是攻击（正如我从他的创伤带来的痛苦中所感受到的），他都确保我没有同时代人因同受称赞而给我带来荣誉，也没有因同受辱骂而给我带来安慰。\par

% structure: p0009 source=h2 target=subsection reason=relative-heading-rank; base=h2

\subsection{他在\OriginalTerm{论原理}{\GreekText{Περὶ} \GreekText{᾿Αρχῶν}}的前言中给予我虚假的称赞。现在，由于我为自己辩护，他写了三卷书攻击我如敌人。}

3. 我手中有你的信，你在信中告诉我，我受到了指控，并期望我答复指控者，以免沉默被视为承认他的指控。我承认我送去了答复；但尽管我受到伤害，我仍遵守友谊的法则，为自己辩护而没有指控我的指控者。我把它表现得好像一位朋友在罗马提出的异议正被世界各地的许多敌人四处传播，以至于每个人都认为我是在答复这些指控，而不是针对那个人。你是否要告诉我，我另有选择，我受友谊法则的约束，在受到指控时应保持沉默，而且，尽管我感到自己的脸（可以这么说）沾满了污秽，被异端的脏物所溅污，甚至连用清水洗一下都不行，以免不义之举被归咎于他？这样的要求不是任何人应该提出的，也不是任何人应该接受的。你公开攻击你的朋友，以崇拜者的面具提出对他的指控；他甚至不被允许证明自己是公教徒，或答复说这种赞誉所依据的所谓异端并非源于与异端的一致，而是出于对伟大天才的钦佩。他认为把这卷书译成拉丁文是可取的；或者，如他更希望人们认为的那样，他虽然不愿意，却是被迫这样做的。但是，当我在隐退中，与他相隔万水千山时，他有什么必要把我牵扯进来呢？他为何要让我遭受众人的敌意，用他的称赞对我造成更大的伤害，而把他自己推出来作为榜样所得到的好处却更小？现在，既然我已拒绝他的称赞，并删除了他所写的内容，表明我并非我的朋友所宣称的那样，我听说他勃然大怒，写了三卷书攻击我，充满优雅的阿提卡式嘲讽，将他先前称赞的那些东西作为攻击的对象，并把\Person{Origen}的不敬神学说变成指控我的理由；尽管在那一封他如此赞扬我的序言中，他这样谈论我：\Quote{我将遵循我的前辈们所订立的翻译规则，尤其是我刚才提到的这位作家所遵循的那些规则。他将\Person{Origen}七十多篇讲道集以及少数关于宗徒的注释译成了拉丁文；在这些翻译中，凡希腊原文有障碍之处，他都予以处理，并对所用措辞加以修订，使拉丁读者找不到与我们信仰相悖的词语。因此，我打算追随他的脚步，即使不能展现同样有力的辞采，至少遵守同样的规则。}\par

% structure: p0011 source=h2 target=subsection reason=relative-heading-rank; base=h2

\subsection{他谈到我与他在信仰上合一；但他的信仰是什么？为何他的著作秘而不宣？我能应对任何攻击。}

4. 这些话是他自己的，他无法否认。那极为优雅的文体和刻意雕琢的措辞，以及超越这一切、在此显出的基督徒般的\Quote{纯朴}，揭示了作者的性格。但还有另一个不同的层面：看来，欧瑟比对这书作了篡改；而现在，这位指责奥力振的朋友，这位如此关心我名声的人，宣称欧瑟比和我都一同犯了错，然后又声称我们持有相同的正确观点，而且在同一部著作中。但他现在不能一面视我为敌，称我为异端者，而片刻之前却说他自己的信仰与我的信仰并非不谐。那么，我必须问他，他这种平衡而含糊的说话方式意味着什么：\Quote{拉丁读者}，他说，\Quote{在这里不会发现任何与我们信仰不协的东西。} 他称之为\Quote{他的}信仰是什么？是罗马教会所持守的信仰吗？还是奥力振著作中包含的信仰？如果他回答\Quote{罗马教会的信仰}，那么我们是天主教徒，因为我们在翻译中没有采纳奥力振的任何错误。但如果奥力振的亵渎便是他的信仰，那么，虽然他试图指责我自相矛盾，他自己却证明是异端者。如果称赞我的人是正统的，那么按他自己的承认，他便将我看作与他同享正统的人。如果他是异端的，那么他表明，他在我解释之前称赞我，是因为他认为我与他共担错误。不过，要回击他的这些在角落里窃窃私语、暗中诽谤的著作，等它们出版、从阴暗处来到光明中，并因朋友的热情或敌人的轻率而落到我手中时，再作答也不迟。我们不必过于害怕那些其作者不敢出版、只让同党阅读的攻击。到那时，而不是在此之前，我将要么承认他的指控的正当性，要么驳斥它们，要么将指控回击到控告者身上：并表明我的沉默不是由于良心不安，而是出于忍耐。\par

5. 与此同时，我渴望在读者默默的判断中摆脱嫌疑，并在朋友面前驳斥最严重的指控。我不愿给人留下先动手的印象，因为我即使受了伤，也未向攻击者还击，只是设法医治自己的伤口。我恳请读者，不要偏袒任何人，让责任落在先动手之人身上。他不满足于动手；而且，仿佛他是在对付一个已经哑口无言、永远不会再开口的人，他写了三本精心制作的著作，并从我的著作中摘录了一份堪与玛尔基翁的\Quote{矛盾}相媲美的清单。我们全都迫不及待地想知道他的道理是什么，以及这种我们未曾预料的疯狂是什么。或许他已经学会了（尽管时间很短）足以使他成为我老师的一切，一股突然的雄辩将揭示出无人曾料到的他所拥有的知识。\par

父啊，允准；大能的耶稣，允准。\newline{}让他开始面对面交战。\par

尽管他可能挥舞控告的长矛，全力向我们投掷，我们信赖主、我们的救主，祂的真理必如盾牌般围绕我们，我们必能同咏唱着说：\BibleQuote{他们的打击成了小儿的箭矢，} 以及 \BibleQuote{虽有大军向我安营，我的心毫不畏惧；虽有战争向我兴起，我依然满怀信赖。} 但这些容后再谈。现在让我们回到我们开始的地方。\par

% structure: p0016 source=h2 target=subsection reason=relative-heading-rank; base=h2

\subsection{我翻译了\OriginalTerm{论首要原理}{\GreekText{Περὶ} \GreekText{᾿Αρχῶν}}，因为你们要求如此，也因为他的翻译抹杀了奥力振的异端。}

6. 他的追随者们向我反对（而且\par

厌倦了工作，\newline{}他们挥动着刻瑞斯的臂膀，\par

我将奥力振的著作\OriginalTerm{论首要原理}{\GreekText{Περὶ} \GreekText{᾿Αρχῶν}}译为拉丁文，这是有害且与教会信仰相悖的。我对他们的回答简短而扼要：我的兄弟潘玛丘，是您和您的朋友们的信函迫使我这样做。你们宣称，这些书已被另一个人错误地翻译，并且不少地方被篡改、添加或更改。而且，唯恐你们的信函无法令人信服，你们寄来了这个译本的一份副本，以及其中称赞我的序言。我刚看过这些文件，便立即注意到，奥力振关于圣父、圣子和圣神的亵渎性教义——这是罗马人的耳朵所不能忍受的——已被译者篡改，使之具有更正统的含义。他的其他教义，关于天使的堕落、人类灵魂的沉沦、关于复活的错误言论、关于世界的观念，或更确切地说，关于伊壁鸠鲁的中间空间、关于万物恢复到平等状态，以及其他比这些更糟糕的教义——若要一一叙述，花费的时间太长——我发现他要么按希腊文原样翻译，要么以更强硬、夸张的措辞从狄迪慕斯——他是奥力振最公开的维护者——的著作中采用。这一切的效果是，读者发现该书表达了关于天主圣三的大公道理，便会在毫无警觉的情况下接纳这些异端观点。\par

% structure: p0020 source=h2 target=subsection reason=relative-heading-rank; base=h2

\subsection{我的翻译消除了歧义，并揭示了这书以及前一译本的本来面目。}

7. 一个不是他朋友的人或许会对他说：\Quote{要么改变一切坏的，要么公开一切你认为完全好的。如果你为了纯朴基督徒的缘故，删去一切有害的，并且不愿将那些你说是异端者添加的东西译成外语，那就告诉我们一切有害的吧。但是，如果你打算做一个真实且忠实的翻译，为何改变一些而留下另一些不变？你在序言中公开声明你已修订坏的，并留下了一切最好的；因此，如果作品中任何地方被证明是异端的，你就不能享受翻译者的许可，而必须承担作者的权威；并且你会被公开定罪，犯有以蜜涂抹毒杯的犯罪意图，以便感官所遇到的甜味掩盖致命的毒液。}这些，以及比这些更严厉的话，敌人会说；他会将你带到教会的法庭前，不是作为一部坏作品的翻译者，而是作为赞同其教义的人。但我只满足于为自己辩护。我将我在希腊文\WorkTitle{\OriginalTerm{论原理}{\GreekText{Περὶ} \GreekText{᾿Αρχῶν}}}的书卷中所见的，忠实地以拉丁文表达出来，不是希望读者相信我的译文，而是希望他不相信你的译文。我期望我的工作带来双重益处：首先揭露作者的异端，其次证明翻译者的不可靠。为使无人以为我赞同我所翻译的教义，我在序言中声明我是如何被迫作此翻译，并指出读者不应相信什么。第一个翻译为作者带来荣耀，第二个带来羞辱。前者召唤读者相信其教义，后者促使他拒绝相信。在第一个译本中，我被违背己意地声称是赞美作者；在第二个译本中，我非但不赞美他，反而被迫控告那赞美他的人。同样的工作由各人完成，但意图不同；同样的旅程产生了两种不同的结果。我们的朋友删去了存在的词句，声称书卷已被异端者败坏；又添加了不存在的词句，声称这些断言是作者在其他地方说的；但他永远不能说服我们，除非他能指出他所说的实际出处。我的努力是不改变任何实际存在的内容；因为我翻译此书的目的，是揭露我所翻译的谬误教义。你视我为仅仅一个翻译者吗？我更多。我成了告发者。我告发了一个异端者，以清除教会的异端。我先前在某些方面赞美\Person{奥力振}的原因，已在此书前附的论文中阐明。导致我翻译的唯一原因，现在读者已经明了。无人有权指控我认同作者的不虔诚，因为我以虔诚的意图为之，即揭露那曾被以虔诚之名向各教会推荐的不虔诚。\par

% structure: p0022 source=h2 target=subsection reason=relative-heading-rank; base=h2

\subsection{我对\Person{奥力振}的\WorkTitle{注释}的翻译未引起任何骚动；他的第一个译本，\Person{庞斐禄}的\WorkTitle{辩护}，却激起了全\Place{罗马}的愤慨。}

8. 我虽然多年翻译了\Person{奥力振}的许多作品，但从未引起丑闻；而你，虽以前默默无闻，却凭你的第一也是唯一一部作品，因你的鲁莽举动而臭名昭著。你的序言告诉我们，你还翻译了殉道者\Person{庞斐禄}为\Person{奥力振}辩护的作品；你竭尽全力阻止教会谴责一个其信德被殉道者见证的人。事实是，如我之前所述，\Place{凯撒勒雅}主教\Person{欧瑟伯}，他在世时是亚略派的旗手，写了一部六卷的巨著为\Person{奥力振}辩护，用许多证据表明\Person{奥力振}在他（\Person{欧瑟伯}）看来是公教徒，即在我们看来，是亚略派。这六卷中的第一卷，你已翻译，并将其归于殉道者名下。因此，当你将同样的诽谤加之于殉道者时，我毫不奇怪你希望使我这个渺小无足轻重的人显得像是\Person{奥力振}的仰慕者。你更改了少许关于圣子与圣神的论述，你知道这些会冒犯罗马人，而其余从头至尾未加改动；事实上，你在这本你称之为\Person{庞斐禄}的\WorkTitle{辩护}中，正如同你在\Person{奥力振}的\WorkTitle{\OriginalTerm{论原理}{\GreekText{Περὶ} \GreekText{᾿Αρχῶν}}}的翻译中所做的一样。如果那本书是\Person{庞斐禄}的，那么六卷中的哪一卷是\Person{欧瑟伯}的第一卷？在你伪称是\Person{庞斐禄}的那一卷中，提到了后面的书卷。同样，在第二卷及后续卷中，\Person{欧瑟伯}说他在第一卷中已说过某件事，并为自己重复它们而辩解。如果全书是\Person{庞斐禄}的，你为何不翻译其余书卷？如果是另一人的作品，你为何更改作者姓名？你无法回答；但事实自行作答：你以为人们会相信殉道者，尽管他们本会对亚略派的领袖避之唯恐不及。\par

% structure: p0024 source=h2 target=subsection reason=relative-heading-rank; base=h2

\subsection{然而那作品其实是\Person{欧瑟伯}的，他告诉我们\Person{庞斐禄}什么也没写。}

9. 我要直说你的意图是什么吗，我最天真单纯的朋友？你认为我们能相信你无意中把一位异端者的书冠以殉道者的名字，从而使无知的人因信赖基督的见证而成为奥力振的辩护者吗？考虑到你以博学著称，在整个西方被赞誉为杰出的文学家，以致你们同党的人都称你为他们的\OriginalTerm{领袖}{Coryphæus}，我不认为你会不知道欧瑟伯的目录，其中记载着殉道者庞斐禄从未写过一本书。欧瑟伯本人，庞斐禄的挚友与同伴，以及他赞美的宣扬者，以优雅的文字写了三卷书，记述庞斐禄的生平。在这些书中，他以非凡的颂词褒扬他的其他品格，并将其谦卑赞颂到天上；但关于他的文学兴趣，他在第三卷中写道：\Quote{有哪个爱书的人没有在庞斐禄那里找到朋友？如果他得知他们中有人缺乏生活必需品，他会尽力帮助他们。他不仅会把圣经借给他们阅读，而且会最乐意地赠送给他们，不仅对男人，也对女人，如果他看到她们喜欢阅读。因此他存有大量手稿，以便在需要时能够给那些想要的人。然而，他自己却什么也没有写过，除了寄给朋友的私人信件，他对自己如此谦逊。但他对古代作者的作品则极其勤奋地研究，并经常沉浸于其中。}\par

% structure: p0026 source=h2 target=subsection reason=relative-heading-rank; base=h2

\subsection{在德奥斐洛与阿纳斯塔西谴责奥力振之后，鲁菲诺放弃这种伪装的辩护是明智的。}

10. 你瞧，奥力振的斗士，庞斐禄的赞颂者，宣称庞斐禄什么也没有写，没有写过任何他自己的论著。你不能以庞斐禄是在欧瑟伯出版之后才写了这本书为托词，因为欧瑟伯是在庞斐禄获得殉道冠冕之后才写作的。那么你现在能做什么？许多人因你以殉道者之名出版的书而良心受伤；他们不理会谴责奥力振的主教们的权威，因为他们认为一位殉道者赞扬了他。德奥斐洛主教和阿纳斯塔西教宗的信函有何用？他们穷追这个异端者遍及世界各地，但当你的书以庞斐禄之名出现，与他们的信函对立时，殉道者的见证就可以对抗主教们的权威。我认为你最好对这本误题书名的书，像你对\WorkTitle{\OriginalTerm{论原理}{\GreekText{Περὶ} \GreekText{᾿Αρχῶν}}}那本书所做的一样。接受我的友好建议，不要怀疑你技艺的力量；要么说这本书不是你写的，要么说它被长老欧瑟伯篡改过。要证明这本书是你翻译的是不可能的。没有你的笔迹可证明；你的文笔并非无人能够模仿。或者，作为最后的手段，如果事情到了证明的地步，而你的厚颜被众多见证压倒，你就仿效\Person{斯特西克诺斯}唱一首翻案诗。你悔改已做的事，胜于让一位殉道者继续受诬蔑，以及那些受骗者继续在错误中。你不需要为改变你的意见而感到羞耻；你没有那么大的名声或权威，以致因承认错误而蒙羞。以我为榜样吧——你如此爱我，没有我你既不能活也不能死——说我说过的话，当你曾赞扬我而我为自己辩护时。\par

% structure: p0028 source=h2 target=subsection reason=relative-heading-rank; base=h2

\subsection{我曾赞扬欧瑟伯与奥力振仅作为作家；而被迫谴责他们是异端者。为什么这要被误解？}

11. 凯撒勒雅主教\Person{欧瑟伯}，我前面提到过他，在他的\WorkTitle{为奥力振辩护}第六卷中，对主教与殉道者\Person{美多迪}提出了与你在我身上因你的赞美而提出的同样的抱怨。他说：\Quote{美多迪怎敢在说了关于奥力振教义的这样那样的话之后，现在又写文章反对他？}现在不是谈论这位殉道者的时候；一个人不能在所有的场合都同样地讨论每一件事。目前只须提及，一位\OriginalTerm{亚略派}{Arian}人士对一位最杰出且最雄辩的人并一位殉道者说出了同样的抱怨，而您起初作为朋友将其作为赞美的话题，后来当您被冒犯时却将其转变为指控。我在此给了您一个机会，如果您愿意，可以借此构建对我的诽谤。\Quote{这是怎么回事，}您可能会问，\Quote{我现在贬低\Person{欧瑟伯}，而此前在其他场合却赞美他？}的确，\Person{欧瑟伯}的名字不同于\Person{奥力振}，但抱怨的理由在两种情况下是相同的。我因\Person{欧瑟伯}的\WorkTitle{教会史}、\WorkTitle{编年史}和\WorkTitle{圣地志}而赞美他，我通过将这些著作翻译成拉丁文，把它们献给了与我同语言的人。难道应该称我为\OriginalTerm{亚略派}{Arian}，因为那些书的作者\Person{欧瑟伯}是亚略派吗？如果你敢称我为异端者，请回想你的\OriginalTerm{论首要原理}{\GreekText{Περὶ} \GreekText{᾿Αρχῶν}}序言，你在其中为我作证，证明我与你信仰相同：同时我恳求你耐心听一听这位曾是你朋友的人的辩解。你与别人激烈争论，与你同阶层的人互相责骂；此事是对是错，由你来说。但是，即使是对一个兄弟的真实指控，我也是厌恶的。我说这话不是为了责备别人；我只是说我自己不会这样做。我们彼此相隔遥远。我对你犯了什么罪？我的错是什么？是因为我回答说我不是\OriginalTerm{奥力振派}{Origenist}吗？难道因为我为自己辩护，你就应该被视为被指控者吗？如果你说你不是奥力振派，而且从未是过，我相信你的郑重声明：如果你曾经是，我接受你的忏悔。如果我现在是你说你曾经所是的人，你为什么要抱怨？或者我的错在于我竟敢在你之后翻译\OriginalTerm{论首要原理}{\GreekText{Περὶ} \GreekText{᾿Αρχῶν}}，而我的翻译被认为有损于你的作品？但我该怎么办？你对我的赞美或对我的控诉已经传到了我这里。你的\Emph{赞美}是如此强烈且持久，以至于如果我默许它，每个人都会认为我是异端者。看看我从罗马收到的信末尾的话：\Quote{请澄清人们对你的怀疑，并证实你的控告者所言为虚；因为如果你不加理会，你会被认为同意他的指控。}在如此条件逼迫下，我决定翻译这些书，并请你注意我所作的回答。回答如下：\Quote{这就是我的朋友为我设定的立场（请注意，我没有说\InnerQuote{我朋友}，以免显得针对你）；如果我保持沉默，我将被视为有罪；如果我回答，我将被视为敌人。这两个条件都是艰难的；但两者之中我将选择容易的路：因为争吵可以愈合，但亵渎不得赦免。}你注意到，我认为这是加在我身上的重担；我不情愿且抗拒；我只是借着必要性这一理由来平息我对这次翻译工作将导致争吵的预感。如果你翻译\OriginalTerm{论首要原理}{\GreekText{Περὶ} \GreekText{᾿Αρχῶν}}时没有影射我，你就有权抱怨我后来翻译它们对你不利。但现在你无权抱怨，因为我的工作只是对你以赞美为伪装对我进行的攻击的回应；因为凡你所称的赞美，所有人都理解为指控。让我们彼此确认，你指控了我，那么你就不会对我的答复感到愤慨。但现在假设你写作时怀有善意，你不仅清白而且是最忠实的朋友，口中从未出过谎言，你无意中伤害了我。那对我这个受伤者而言又有什么关系？难道因为你无意伤害，我就不该治疗我的伤口吗？我被击倒、被刺穿，胸口的创伤无法平息；我原本白皙的四肢沾满了鲜血；而你对我说：\Quote{请你不要触碰你的伤口，以免我被认为伤害了你。}然而，这部翻译作品是对奥力振的责备，而不是对你的责备。你将你认为由异端者插入的段落改好了。我揭示了整个希腊世界一致归因于他的东西。我们两人中谁的观点更真实，既不是由我，也不是由你来判断；让读者用批评的权杖来触碰它们。我为自己回信的那整封信是针对异端者和我的控告者的。如果我对待异端者稍微严厉一点，并在公众面前揭露他们的诡计，那关你这位自称正教人士和我的崇拜者什么事？你应该为我的抨击之辞高兴；否则，如果你对它们感到恼怒，你可能被认为是异端者本人。当某篇文章针对某种恶习，但没有指名道姓时，如果一个人生气，他就是在控告自己。一个智者，即使感到受伤，也会装作没有意识到过错，并以其面容的宁静驱散心中的阴云。\par

% structure: p0030 source=h2 target=subsection reason=relative-heading-rank; base=h2

\subsection{我写了一封友好的信给\Person{鲁菲努}，却被我的朋友扣留了。}

12. 否则，如果一切反对\Person{奥力振}及其追随者的话都被认为是由我针对你而说，那么我们就必须认为，\Person{德奥斐洛}和\Person{厄比法尼}两位教宗以及其余主教们——我最近应他们的请求所翻译的信函，是意在攻击你并将你撕碎；我们还必须认为，皇帝们下令将奥力振派从\Place{亚历山大里亚}和\Place{埃及}驱逐出去的诏书，是出于我的口授。\Place{罗马城}教宗对这些人的憎恶，无非是我的诡计。在你的译本发表之后，立刻在整个世界燃起的对\Person{奥力振}的仇恨——此前他被人阅读时并未带有偏见——是我的文笔之作。如果我有这么大的力量，我很奇怪你竟不怕我。但我实际上表现得极为节制。在我的公开信中，我采取了所有预防措施，以免你认为其中有任何针对你的内容；但同时我写了一封短信给你，就你的\Quote{praises}向你规劝。这封信我的朋友认为不宜送给你，因为你不在\Place{罗马}，而且据他们告诉我，你和你的同伴们正在散布关于我生活方式的不配基督徒职业的指控。但我已将这封信的副本附于本书之后，以便你能明白你带给我的痛苦，以及我以何等兄弟般的自制承受了它。\par

% structure: p0032 source=h2 target=subsection reason=relative-heading-rank; base=h2

\subsection{我聘请一位犹太人帮助翻译希伯来文，这无可指责。}

13. 我进一步听说，你带着某种批判的尖刻触及我信中的某些要点，并且，你额头上皱起著名的皱纹、双眉紧锁，以堪比\Person{普劳图斯}的机智取笑我，因为我说我有一位名叫\Person{巴拉巴}的犹太人做我的老师。我并不奇怪你把\Person{巴兰尼纳}写成\Person{巴拉巴}，名字的字母有些相似，因为你自己在更改名字本身时竟如此放肆，以至于把\Person{欧瑟伯}变成\Person{庞斐禄}，把异端者变成殉道者。对于像你这样的人必须小心，并敬而远之；否则我可能发现我自己的名字在转眼之间，在我不知情的情况下，从\Person{热罗尼莫}变成\Person{萨尔达纳帕卢斯}。那么请听吧，智慧的柱石，加图式严厉的典型。我从未称他为我的导师；我只是想通过说我曾向这位犹太人请教过，来表明我研读圣经的方法。难道因为我听了\Person{阿波利纳里}和\Person{迪迪慕斯}的课而没有听你的课，我就伤害了你吗？有什么能阻止我在信中提及那位最雄辩的人\Person{额我略}呢？所有拉丁人中，谁能与他相比？我完全可以以他为荣并欣喜。但我只提及那些受责备的人，以表明我阅读\Person{奥力振}就如同我听过他们的课一样，也就是说，并非因为他在信仰上的纯正，而是因为他学识的卓越。\Person{奥力振}本人，以及\Person{克肋孟}、\Person{欧瑟伯}和许多其他人，当他们讨论圣经问题并想为所说的话寻求犹太人的权威时，会写道：\Quote{一位希伯来人告诉了我这个，}或\Quote{我从一位希伯来人那里听说，}或\Quote{这是希伯来人的意见。}\Person{奥力振}确实提到与他同时代的族长\Person{胡伊鲁斯}，并在他的\WorkTitle{依撒意亚诠释}第三十卷的结尾（在该卷末尾他解释了\Quote{祸哉，达味攻取的阿黎耳}这句话）中，引用了\Person{胡伊鲁斯}对这句话的解释，并承认因他的教导而接受了比他先前所持的更真实的意见。他不仅认为第八十九篇圣咏——其标题为\Quote{天主的人梅瑟的祈祷}——是\Person{梅瑟}所写，而且根据\Person{胡伊鲁斯}的意见，认为接下来没有标题的十一篇圣咏也是\Person{梅瑟}所写；并且他毫不迟疑地将希伯来教师的观点插入他对希伯来文圣经的注释中。\par

% structure: p0034 source=h2 target=subsection reason=relative-heading-rank; base=h2

\subsection{在我知道\OriginalTerm{论首要原理}{\GreekText{Περὶ} \GreekText{᾿Αρχῶν}}之前赞美\Person{奥力振}，这并不奇怪。}

14. 据说最近有一次宣读\Person{德奥斐洛}揭露\Person{奥力振}错误的信时，我们的朋友塞住了耳朵，并与在场所有人一起明确谴责了如此多邪恶的始作俑者；并且他说直到那一刻他才知道\Person{奥力振}写过任何如此错误的东西。我对此不置一词：我不作他人可能作出的评论，即他不可能不知道他自己翻译的东西，以及他竟以殉道者的名义出版了一个异端者为其所作的辩护——他也曾在自己书中为此辩护；关于此事，我稍后若有时间，会写一些反对的意见。我只作一个不容反驳的评论。如果他可能误解了他所翻译的，为什么我不可能之前没有读过\OriginalTerm{论首要原理}{\GreekText{Περὶ} \GreekText{᾿Αρχῶν}}这本书，而只读了我翻译的那些讲道辞——在这些讲道辞中，他本人也作证没有任何错误呢？但是，如果与他所表示的意见相反，他现在为那些他以前曾称赞过我的事情而责备我，他将陷入两难：要么他称赞我是因为他相信我是一名异端者，但同时承认他与我意见一致；要么，如果他以前称赞我是正统的，那么他现在的指控就化为乌有，并且完全是出于恶意。但也许，他以前作为我的朋友对我错误保持沉默，而现在因生我的气，就将隐藏的事揭露出来。\par

% structure: p0036 source=h2 target=subsection reason=relative-heading-rank; base=h2

\subsection{十五、指责似乎前后不一，但我只是听说而已。}

15. 这种对友谊的背弃并不足以赢得我的信赖；公开的敌意也带来了作假的嫌疑。但我仍将鼓起勇气去面对他，并问他我表达了何种异端学说，以便我或能像他一样，表示遗憾并宣誓我从未知道奥力振的坏学说，他的不忠如今才由教宗德奥斐洛首次告知我；或至少证明我的观点是正确的，而他，如其一贯所为，未能理解这些观点。在我所写的\WorkTitle{厄弗所书注释}——我听说他以此作为指控的依据——中，我不可能既说得正确又说得错误；同一泉源不可能既涌出甜水又涌出苦水；既然整部著作中我都谴责那些相信灵魂出自天使的人，我不可能突然忘记自己，为我先前谴责的观点辩护。他很难因我愚昧而提出反对，因为他在其作品中称我为最有文化和口才的人；否则，他所归咎于我的那种愚蠢的冗词赘语，更像是讼棍和空谈家所为，而非有口才之人。我不知道他的书面指控的要点何在，因为传到我耳中的只是传言，而非其文字；如宗徒所教导，击打空气是愚妄之事。然而，在真相未明之前，我必须在不确定中作答：我将以晚年之身，先教导我的对手我在年轻时所学的一课：言辞有多种形式，并且根据主题，不仅句子，而且著作中的词语也会变化。\par

% structure: p0038 source=h2 target=subsection reason=relative-heading-rank; base=h2

\subsection{注释者的职责。}

16. 例如，克律西波和安提帕特专注于棘手的问题；德摩斯梯尼和埃斯基涅斯用雷霆般的声音彼此攻击；吕西亚斯和伊索克拉底风格平易而悦人。这些作家之间有奇妙的差异，尽管每位在其领域都完美无缺。再读读\WorkTitle{图利乌斯致赫伦尼乌斯的书}，读读他的\WorkTitle{修辞学}；或者，既然他告诉我们这些书是以最初未完成的状态从他手中流出的，那就翻阅他\WorkTitle{论演说家}的三卷书，其中他引入了克拉苏和安东尼——当时最雄辩的演说家——之间的讨论；还有一本他年老时写给布鲁图的第四卷书\WorkTitle{演说家}：你会认识到，历史著作、演说、对话、书信和注释各有其特殊风格。我们现在处理的是注释。在我所写的\WorkTitle{厄弗所书注释}中，我仅仅追随奥力振、狄第默和阿波利纳里（他们的学说彼此大相径庭），但仅限于与我信仰的真挚一致：因为注释的功用是什么？是诠释他人的话语，将他说得隐晦之处用平白语言表达出来。因此，它列举许多人的意见，并说：有些人这样理解这段经文，有些人那样理解；他们试图用如此这般的证据或理由来支持自己的意见和理解：这样，有智慧的读者在阅读这些不同的解释，并有许多解释摆在他面前供他接受或拒绝之后，可以判断哪种解释最真实，并像一位好钱币兑换商一样，拒绝伪币。难道注释者要为所有这些不同的解释和所有这些相互矛盾的意见负责吗？只因他将许多人在他注释的单部作品中所作的解释写下来？我想，当你年轻时，你读过阿斯佩尔对维吉尔和萨鲁斯特的注释，武尔卡提乌斯对西塞罗演说的注释，维克托里努斯对其对话录和泰伦提乌斯喜剧的注释，以及我的老师多纳图斯对维吉尔的注释，还有其他作家对其他作家如普劳图斯、卢克莱修、弗拉库斯、珀尔西乌斯和卢卡的注释。你会因为那些注释这些作家的人没有坚持单一解释，而是列举他们自己的观点和他人在同一段上的观点而指责他们吗？\par

% structure: p0040 source=h2 target=subsection reason=relative-heading-rank; base=h2

\subsection{我们必须区分写作方法，不应期望有教养的人的各种作品都有俗气的简单。}

17. 关于希腊人，我什么也不说，因为你自夸了解他们，甚至到了这样的程度：你说，在致力于外国文学时，你已经忘记了本国语言。我担心，按古老的谚语，我可能会像猪教密涅瓦，或人扛着柴火进树林。我只感到惊奇，作为我们这个时代的阿里斯塔科斯，你竟对这些每个男孩都知道的事情表现出无知。无疑，这是因为你的心思专注在你所写的内容上，但也有一部分原因是你对制造对我的诽谤如此目光锐利，以致你轻视语法家和修辞学家的训诫，你不试图纠正当句子变得复杂时被颠倒的词语，也不避免某些辅音的粗俗并列，或逃脱充满脱漏的风格。指出一两处伤口是可笑的，因为整个身体都已虚弱破碎。我不会选出部分加以批评；让他选出任何没有错误的部分吧。他一定连苏格拉底的名言\Quote{认识你自己}都不知道。\par

驾驭船只，未受训练的农夫畏惧；\newline{}未经培训的侍从不敢给病人\newline{}服用烈性的苦艾。治疗之药\newline{}只有医生能开。工匠\newline{}唯有工匠能使用工具。但诗歌，\newline{}无论训练与否，我们都随意书写。\par

也许他会发誓自己从未学过读写；没有誓言我也很容易相信。或者，他会躲避到宗徒论及自己时所说的话：\BibleQuote{我虽拙于言辞，却不拙于知识。}但他这样说的原因很明显。他曾受希伯来学训练，在加玛里耳脚前长大，虽然他达到了宗徒的品位，却不以为耻地称加玛里耳为老师；他认为希腊的辞藻无关紧要，或者至少，出于谦逊，他不愿炫耀自己对它的认识。如此，\BibleQuote{他的宣讲不在乎言词的智慧，而在乎所表之事的能力。}他轻看别人的财富，因他自己已很富有。然而，并不是对一个目不识丁、每句话都结结巴巴的人，斐斯托斯在他审判座前喊道：\par

\BibleQuote{保禄，你疯了！学问太多使你疯狂。}你几乎只能用拉丁语喃喃自语，更像乌龟爬行而非步行；你应当要么用希腊文写作，以便在不懂希腊文的人中显得通晓外语；要么，如果你试图写拉丁文，就应先找一位文法老师，畏惧戒尺，作为老学生与孩童一同重新学习说话的艺术。即使一个人拥有克洛苏斯和达理阿的财富，学问也不会随钱袋而来。它们是辛劳与劳动的伴侣，是禁食而非饱食的伙伴，是克己而非放纵的同伴。据说德摩斯梯尼耗油多于饮酒，没有工匠像他那样缩短夜晚。他为了发好字母Rho，甚至愿意以狗为师；而你却将我找人教我希伯来字母视为罪行。这就是使人保持无知的智慧：他们不愿学习他们所不知道的。他们忘记了贺拉斯的诗句：\par

为何因虚假的羞耻我选择无知，\newline{}而不愿去学习？\par

那本读作撒罗满著作的\BibleBook{}{智慧书}也说：\BibleQuote{智慧不进人恶意的心灵，也不住在屈服于罪恶的身体内。因为训导的圣神必逃避欺诈，远离无知的思念。}\par

我想，正是你，在街头用吱吱作响的芦管\newline{}扼杀了那可怜的曲调。\par

如果你想要这些，每个学堂里都有许多卷发的家伙，他们会为你唱米利都的歪诗；或者，你可以去贝西人的表演，看到人们因\WorkTitle{猪的遗嘱}或任何小丑表演中这类蠢事而笑得发抖。没有一天你不看到街头打扮花哨的小丑拍打某个傻瓜的屁股，或用蝎子（他卷起让人咬）半拔出人们的牙齿。无知的书籍拥有众多读者，不足为奇。\par

% structure: p0049 source=h2 target=subsection reason=relative-heading-rank; base=h2

\subsection{我的断言是真的，即奥利振允许使用谎言。}

18. 我们的朋友对我提到奥利振主义者通过虚假誓言聚会感到不快。我提到了写着这话的书，即奥利振\WorkTitle{杂论}第六卷，他在其中试图将我们基督宗教的教义与柏拉图的意见调和。柏拉图在\WorkTitle{理想国}第三卷中的话如下：\Quote{苏格拉底说，真理尤其应当培养。然而，如果正如我刚才所说，谎言对天主是可耻无用的，对人有时却有用，只要它被用作刺激或药物；因为没有人会怀疑，医生可以允许某种这样的自由，但这必须从不熟练的人手中拿走。那很正确，有人回答；如果允许任何人这样做，那必定是城邦统治者的职责，他们有时为了迷惑敌人或造福国家和公民而说谎。另一方面，对那些不知道如何善用谎言的人，则应完全禁止。}现在看奥利振的话：\Quote{当我们考虑诫命\BibleQuote{每人应与邻人说实话}时，不必问：谁是我的邻人？但我们应该仔细考虑哲学家的谨慎意见。他说，对天主来说，谎言是可耻无用的，但对人有时却有用。我们不可认为天主会说谎，即使是在\OriginalTerm{救恩经济}{economy}中；只是，如果听者的益处需要，祂以模棱两可的话说话，并以谜语启示祂所愿意的，同时注意既保持真理的尊严，又使那若赤裸地呈现于群众前会有害的东西，被包裹在面纱中而得以揭示。但一个人因责任而不得不撒谎时，必须非常谨慎，并且像使用刺激或药物那样使用谎言，严格保持其限度，不超出友弟德在对待敖罗斐乃时所遵守的界限，她用智慧伪装话语而战胜了他。他应效法艾斯德尔，她因长久隐瞒自己的种族而改变了阿尔塔薛西斯的心意；还有先祖雅各伯，据记载，他通过诡计和谎言获得了父亲的祝福。由此可见，如果我们以谎言谋求其他目的而非获得大善，我们就会被视为那说\BibleQuote{我是真理}者的敌人。}奥利振写了这些，我们无人能否认。他写在他献给\OriginalTerm{成全者}{perfect}，即他自己的门徒的书中。他的教导是：师傅可以说谎，但门徒不可。由此推论，善于说谎、毫不犹豫地将任何随口编造的东西讲给弟兄的人，表明自己是优秀的老师。\par

% structure: p0051 source=h2 target=subsection reason=relative-heading-rank; base=h2

\subsection{关于误译咏第二篇的指控很容易解释。}

19. 我听说他还挑剔我，因为我翻译了第二篇圣咏中的一个短语。在拉丁文中它是\Quote{学习纪律}，在希伯来文中写的是\OriginalTerm{乃斯库·巴尔}{Nescu Bar}；我在注释中将它翻译为\Quote{朝拜圣子}；后来当我把整部圣咏集翻译成拉丁文时，仿佛我忘记了我先前的解释，我写了\Quote{纯粹朝拜}。当然，没有人能否认这些解释是相互矛盾的；我们必须原谅他，因为他对希伯来文一无所知，甚至常常在拉丁文上犯难。\OriginalTerm{乃斯库}{Nescu}，按字面翻译是\Quote{亲吻}。我不愿给出令人不快的译法，而倾向于遵循意义，给了\Quote{朝拜}这个词；因为那些敬拜的人往往会亲吻自己的手并脱帽，这可以在\Person{约伯}的例子中看到，他宣称从未做过这些事，并说：\Quote{\BibleQuote{如果我看见太阳发光，或月亮在光明中行走，我的心在暗中喜悦，我用口亲吻我的手，这是极大的不义，是对至高天主的谎言。}\BibleRef{约伯传 31:26-28}}希伯来人根据他们语言的特性，用这个词\Quote{亲吻}来表示朝拜；因此我根据我所处理的语言的使用习惯进行了翻译。然而，\OriginalTerm{巴尔}{Bar}这个词在希伯来文中有几个含义。它意为\Quote{儿子}，如在Barjona（鸽子之子）、Bartholomew（Tholomaeus之子）、Bartimaeus、Barjesus、Barabbas等词中。它也意为\Quote{小麦}、\Quote{一捆谷物}、\Quote{蒙拣选的}和\Quote{纯粹的}。那么，当一个词的意义如此不确定时，如果我在不同地方以不同方式翻译它，我犯了什么罪呢？如果在我的注释中，我取了\Quote{朝拜圣子}的意义，在那里有更自由的讨论，而在圣经本身的翻译中，我说\Quote{纯粹地朝拜}或\Quote{选择性地朝拜}，这样我就不会被认为是随意翻译或为犹太人提供挑剔的借口。后一种译法，是\Person{阿奎拉}和\Person{辛马库斯}的；我看不出教会的信仰因读者看到犹太人以多少种不同的方式翻译这句话而受到损害。\par

% structure: p0053 source=h2 target=subsection reason=relative-heading-rank; base=h2

\subsection{在翻译者和注释者的困境中，我们必须尽可能寻求帮助。}

20. 你的\Person{奥利振}允许自己讨论灵魂转生，引入对无数个世界的信仰，使有理性的受造物依次穿上身体，说基督曾多次受苦，并将再次受苦，因为已经有益的事总是有益去承担。你自己也如此放肆，竟然将异端者变成殉道者，并虚构奥利振著作的异端篡改。那么，我为什么不可以讨论词语呢？在做注释者工作时，将我从希伯来人那里学到的教给拉丁人？如果不是一个冗长的过程且带有夸耀的味道，我现在就想向你展示在伟大的导师门前等候、从真正的工匠那里学习技艺是多么有裨益。如果我能这样做，你就会看到希伯来文呈现出怎样一个名字和词语模棱两可的纠结森林。正是这一点为各种译法提供了如此广阔的天地：因为意义不确定，每个人都采取他认为最一致的翻译。我为什么要带你去看任何生僻的作家呢？再通读一遍\Person{亚里士多德}和注释亚里士多德的\Person{亚历山大}；你会从阅读这些中认识到存在多么丰富的未知事物；然后你或许会停止挑剔你的朋友，因为那些事情你连在梦中都从未想过。\par

% structure: p0055 source=h2 target=subsection reason=relative-heading-rank; base=h2

\subsection{在\WorkTitle{厄弗所书注释}中，我坦率地给出了\Person{奥利振}等人的观点。}

21. 我的兄弟\Person{保利尼安}告诉我，我们的朋友指责了我在\WorkTitle{厄弗所书注释}中的某些事情；他记住了其中一些批评，并指出了被指责的具体段落。因此，我不能拒绝回应他的言论，我请求读者，如果我在陈述和反驳他的指控时有些冗长，请允许讨论的必要条件。我不是在指责别人，而是在努力为自己辩护，并反驳强加给我的异端虚假指控。关于《厄弗所书》，\Person{奥利振}写了三卷书。\Person{迪迪慕斯}和\Person{阿波利纳里}也撰写了他们自己的作品。我部分翻译，部分改编；我的方法在我的序言中描述如下：\Quote{这也我想在我的序言中说明。你必须知道，奥利振关于这封书信写了三卷书，我部分地跟随了他。阿波利纳里和迪迪慕斯也发表了关于它的某些注释，我从其中采集了一些东西，虽然很少；并且，按照我认为合适的，我加入了或删除了其他东西；但我这样做的方式是，细心的读者从一开始就可以看出这部作品在多大程度上归功于我，在多大程度上归功于他人。}在这封书信的讲解中发现的任何错误，如果我无法表明它存在于我所说的被翻译成拉丁文的希腊书中，我将承认错误是我自己的，而不是别人的。然而，为了不让人认为我在吹毛求疵，并借这种自我辩解的手段逃避勇敢地面对他，我将列出那些被引为我的错误证据的具体段落。\par

% structure: p0057 source=h2 target=subsection reason=relative-heading-rank; base=h2

\subsection{关于\BibleQuote{祂在创世以前拣选了我们}这段经文。}

22. 首先。在第一卷书中，我引用了\Person{保禄}的话：\BibleQuote{正如祂在创世以前拣选了我们，为使我们成为圣洁无瑕的。}\BibleRef{弗 1:4}我将此解释为并非指，按照\Person{奥利振}的意见，那些存在于先前状态者被拣选，而是指天主的预知；并以这些话结束讨论：\par

\Quote{他的断言，即我们在世界创立之前被拣选，为使我们在他面前，即在天主面前，成为圣洁无瑕疵的，这属于天主的预知，在他眼中，一切将要发生的事已经作成，且在尚未发生之前就已知道。因此，\Person{保禄}在母腹中被预定；\Person{耶肋米亚}在出生之前就被圣化、被拣选、被确认，并作为\Person{基督}的预像，被派遣为外邦人的先知。}\par

这种解释当然没有罪过。\Person{奥力振}以异端意义解释了它，但我跟随了教会的解释。而且，既然注释者的责任是记录许多其他人所表达的意见，而且我在序言中已承诺会这样做，我便记下了\Person{奥力振}的解释，尽管没有提及他那引起恶感的姓名。\par

\Quote{另一位，}我说道，愿意为天主的正义辩护，并显示他并非基于自己的预判和预知拣选人，而是根据被选者的功过，他认为在天空、大地、海洋及其中万物的可见受造物之前，已存在不可见的受造物，其中一部分由灵魂组成，这些灵魂因唯独天主知道的某些原因，被投入这涕泣之谷、这受苦和流徙之地；圣咏作者的祈祷正是适用于此，他正处在这种低下境况，渴望回到他原先的住处：\Quote{我有祸了，因我的客居延长了；我居住在\Place{刻达尔}的帐幕中，我的灵魂经历了长期的流徙。}以及宗徒的话：\Quote{我真不幸啊！谁能救我脱离这死亡的身体？}以及\Quote{归去并与\Person{基督}同在一起更好}；以及\Quote{在我受屈辱之前，我犯了罪。}他还加了更多类似的话。\par

现在请注意，我说了\Quote{另一位愿意为之辩护，}我说道，并没有说\Quote{成功为之辩护}。但如果你因为我将\Person{奥力振}的一篇极长的讨论压缩成简短的陈述，以便让读者一窥其意，就视之为绊脚石；如果你因为我未省略他所说的任何话，就宣称我是他的秘密追随者，那么我要问你：我这样做难道不是必要的吗？以避免你的挑剔。否则，你难道不会宣称，我保持沉默于他大胆谈论的事，并且希腊文本中的断言比我表述的更强烈？因此，我记下了我在希腊文本中发现的一切，尽管以更简短的形式，这样他的门徒就没有什么可以作为新奇之事强加给拉丁人的耳朵；因为对于熟悉之事，我们比对于措手不及之事更容易轻视。但在展示了\Person{奥力振}对这段经文的解释之后，我以这样的话结束这一节，敬请留意：\par

\Quote{\Person{保禄}宗徒并没有说\InnerQuote{他在世界创立之前拣选了我们，是因为当时我们是圣洁无瑕疵的}；而是说\InnerQuote{他拣选了我们，为使我们成为圣洁无瑕疵的}，也就是说，我们这些先前不是圣洁无瑕疵的人，后来可以成为那样。这种说法甚至适用于转向更好的罪人；因此这话仍然真实：\InnerQuote{在你面前，没有一个活着的人能成义}，也就是说，在他的整个生命中，在他存在于世界的整个时期内，没有一个人。如果这段经文这样理解，就反对了在世界创立之前某些灵魂因其圣洁而被拣选、并且他们没有罪人任何败坏的观点。很明显，\Person{保禄}及类似他的人被拣选，不是因为他们是圣洁无瑕疵的，而是他们被拣选和预定，以便在他们后来的生活中，通过他们的行为和德行，成为圣洁无瑕疵的。}\par

那么，在我陈述了我的意见之后，还有人敢指控我赞同\Person{奥力振}的异端吗？我从撰写那些书到现在几乎已十八年，那时\Person{奥力振}的名字在世界上备受推崇，而且他的著作\OriginalTerm{论原理}{\GreekText{Περὶ} \GreekText{᾿Αρχῶν}}尚未传到拉丁人的耳中；然而我明确陈述了我的信仰，并指出了我不赞同之处。因此，即使我的对手本可以在其他地方指出任何异端之处，我也只应承担粗心的过错，而不应承担乖谬的教义，因为无论是此处还是我的其他著作，我都谴责了这些教义。\par

% structure: p0065 source=h2 target=subsection reason=relative-heading-rank; base=h2

\subsection{论及\Quote{远超一切统治、权柄等}的经文}

23. 我将简短处理第二段经文，我的弟兄告诉我这段经文被指摘为应受责备，因为此抱怨极其轻浮，其诽谤性质表露无遗。这段话是\Person{保禄}宣称天主\Quote{使他坐在自己的右边，在天上，远超一切统治、权柄、大能和宰制，以及一切被称呼的名号，不但在现今世界，也在将来世界。}在陈述了已给出的各种解释后，我谈及天主仆人的职务，并论及率领与权柄、大能与宰制，我接着说：\par

\Quote{他们必然有服从于他们的其他人，这些人处于他们的权下，服事他们，并因他们的权威而坚固；这种职务的分配不仅存在于现今世界，也存在于将来世界，以至于每个人会从一级晋升或降落到另一级，有些人上升，有些人下降，并依次处于这些大能、权柄、率领和宰制之下。}\par

接着，我继续描述各种神圣的职务和职份，以地上君王宫殿为喻，我详尽描述了这个比喻，并接着说：\par

\Quote{我们能认为天主——万主之主、万王之王——满足于单一等级的仆人吗？我们谈论天使长，因为还有其他天使，他是他们的首领；因此，除非暗示还有下级，就不会说到率领、权柄和宰制。}\par

但是，如果他以为我之所以成为\Person{奥利振}的追随者，是因为我在释经中提到了这些提升与荣誉、这些上升与下降、增加与减少；那么我必须指出，像\Person{奥利振}那样说天使、革鲁宾和色辣芬变成了魔鬼和人，与说天使们本身被分派了不同的职分——这一教义并不与教会相悖——是截然不同的。正如在人间，因不同的工作而区分出不同的尊严等级，如主教、司铎和其他教会等级各自有其秩序，而他们全都是人；同样，我们可以相信，虽然天使们都保持着尊严，但他们之间有不同程度的卓越，而不必想象天使会变成人，人也重新变成天使。\par

% structure: p0071 source=h2 target=subsection reason=relative-heading-rank; base=h2

\subsection{关于\Quote{在未来的世代等等}}

24. 他指责的第三处经文是我对宗徒的话所作的三重解释：\BibleQuote{为使在未来的世代，他可在基督耶稣内，以他对我们的仁慈，显示他恩宠的极丰富}。第一个是我自己的意见，第二个是\Person{奥利振}所持的对立意见，第三个是\Person{阿波利纳里}给出的简单解释。至于我没有给出他们的名字，我必须请求原谅，理由是出于谦逊。我不愿贬低那些我在一定程度上追随且将其意见译为拉丁文的人。但是，我说，勤勉的读者会立刻查考这些事并形成自己的意见。我在最后重复说：另一个人将\InnerQuote{为使在未来的世代，他可在基督耶稣内，以他对我们的仁慈，显示他恩宠的极丰富}这些话转向不同的意义。\Quote{啊，}你会说，\Quote{我看出你以勤勉的读者的身份阐明了\Person{奥利振}的意见。}我承认我错了。我本应说\Quote{亵渎的读者}而不是\Quote{勤勉的读者}。如果我预料到你会采取这样的措施，我或许会这样做，从而避免你的诽谤之言。我想，称\Person{奥利振}为勤勉的读者是一个大罪，尤其是当我翻译了他的七十卷书并把他捧到天上时——为此我两年前不得不在一篇短文中为自己辩护，以回应你对我的吹捧。在那些你给予我的\Quote{赞美}中，你指责我说\Person{奥利振}是教会的教师，而现在你以敌人的姿态说话，你认为我会害怕你指控我称他为勤勉的读者。为什么，甚至特别节俭的店主、不浪费的奴隶、使我们童年不堪重负的看管者，甚至特别聪明的贼，我们都称之为勤勉；同样，福音中不义管家的行为也被称为聪明。此外，\BibleQuote{世俗之子对于自己的同类，比光明之子更为聪明}，和\BibleQuote{上主天主所造的野兽中，蛇最为狡猾}。\par

% structure: p0073 source=h2 target=subsection reason=relative-heading-rank; base=h2

\subsection{关于\Quote{保禄为耶稣基督的囚犯}}

25. 他指责的第四个理由是在我第二卷书的开头，其中我解释了圣保禄的话：\BibleQuote{因此，我保禄为你们外邦人做了耶稣基督的囚犯}。这段经文本身完全清楚；因此，我只给出其中容易招致恶意评论的那部分注释：\par

\Quote{描述保禄为外邦人做耶稣基督的囚犯的话语，可以理解为他的殉道，因为他正是在罗马被锁链拘禁时写了这封书信，同时写了致\WorkTitle{费肋孟书}、致\WorkTitle{哥罗森书}和致\WorkTitle{斐理伯书}，正如我们先前所表明的。当然我们可以采纳另一种意义，即，既然我们发现这身体在几处被称为灵魂的锁链，灵魂在其中如同被囚禁在狭窄的监狱中，保禄可以说自己被囚禁在身体的锁链中，以至于他不能回去与基督同在；这样他就可以完美地履行向外邦人传教的职务。然而，有些注释者引入了另一种观念，即保禄从他母亲的子宫中，在他出生之前，就被预定并祝圣为外邦人的传道者，后来他承担了血肉的锁链。}\par

在这里，也和先前一样，我对这段经文作了三重阐释：第一个是我自己的观点，第二个是\Person{奥利振}所支持的观点，第三个是\Person{阿波利纳里}与他的教义相反的意见。读读希腊注释吧。如果你发现事实并非如我所述，我将承认我错了。这段经文中我的过错是什么？我认为与之前我回答过的一样，即我没有指出我所引用观点的作者名字。但无需在宗徒的每一处陈述中都举出作者的名字，因为我已在序言中声明我打算翻译他们的作品。此外，说灵魂被束缚在身体中，直到基督再来，并在复活的荣耀中，将我们腐朽可朽的身体改变为不朽坏和不死，这并非一种荒谬的理解方式：因为正是在这个意义上，宗徒使用了这句话：\BibleQuote{我真不幸！谁能救我脱离这死亡的身体呢？}称它为死亡的身体，因为它易受邪恶和疾病的侵害，易受混乱和死亡；直到它同基督一起在荣耀中复活，而它先前不过是脆弱的泥土，却被圣神的热力烘烤成坚实的陶罐，从而改变了它的荣耀等级，而没有改变它的本性。\par

% structure: p0077 source=h2 target=subsection reason=relative-heading-rank; base=h2

\subsection{关于\Quote{身体适当结合等等}}

26. 他所挑出来指责的第五段最为重要，就是我在其中解释宗徒的那句话：\BibleQuote{从祂，全身都靠祂联络得合式，百节各按其职，照着各体的功用彼此相助，便叫身体渐渐增长，在爱中建立自己。}这里我用短短一句话总结了奥力振的解说——他的解说很长，用不同词语重复同一思想，但我没有遗漏他任何一个比喻或断言。当我讲到末尾时，我补充说：\par

\Quote{因此，在万物复兴时，当真正的医生耶稣基督来医治教会这整个身体——如今它分散破碎——每一个人将按他信德的程度和对天主子的认识（\OriginalTerm{认识}{recognize}暗示他从前认识祂，后来又不认识祂了）领受他应有的位置，然后开始成为他从前所是的；然而并非如另一种异端所主张的那样，所有人都被置于同一等级，并通过革新过程都成为天使；而是每个肢体按自己的尺度和职责成为完美的：例如，叛逃的天使将开始成为他受造时的样子，被逐出乐园的人将恢复耕种乐园；}等等。\par

% structure: p0080 source=h2 target=subsection reason=relative-heading-rank; base=h2

\subsection{我引用奥力振的观点，说：\Quote{根据另一种异端。}}

27. 我奇怪以您的卓越智慧竟然没有理解我的解释方法。当我说\Quote{然而并非如另一种异端所主张的那样，所有人都被置于同一等级，即所有人都通过革新过程成为天使}，我清楚表明我提出讨论的是异端观点，且一种异端与另一种不同。您问是哪两种异端？一种是说所有理性受造物将通过革新过程成为天使；另一种是断言在世界的复兴中，每样事物将变成它原本受造时的样子；例如，魔鬼将再次成为天使，人的灵魂将变成它们最初被造时的样子；也就是说，通过革新过程，它们不是变成天使，而是变成天主最初创造的样子，以致义人和罪人平等。最后，为了向您表明我所发展的并非我自己的意见，而是我正在比较的两种异端——我在希腊文献中都找到了它们——我以下面这段话结束我的讨论：\par

\Quote{正如我以前所说，这些在我们的语言中比较晦涩，因为它们在希腊文中是以比喻形式表达的；而在每一个比喻中，当逐字从一种语言翻译成另一种语言时，词语的新生含义就像被荆棘窒息了一样。}\par

如果您在希腊文中找不到我所表达的那个思想，我允许您把我说的一切都当作我自己的。\par

% structure: p0084 source=h2 target=subsection reason=relative-heading-rank; base=h2

\subsection{关于\Quote{丈夫爱妻子如同自己的身子。}}

28. 第六点也是最后一点——我听说他指责我（如果我兄弟没有遗漏什么的话）——就是在解释宗徒的话\BibleQuote{爱妻子的，便是爱自己了，因为从来没有人恨恶自己的肉身，总是保养顾惜，正像基督待教会一样}之后，在我自己的简单解释之后，我提出了奥力振所引发的问题，陈述他的观点，虽未提他的名字，我说：\par

\Quote{可能有人反对说，宗徒的话并不真实，因为他说没有人恨恶自己的肉身，然而那些患有黄疸、痨病、癌症或脓肿的人，宁愿死也不愿活，他们恨恶自己的身体；}接着我自己的意见紧随其后：\Quote{所以，这些话更恰当地应作比喻理解。}\par

当我说比喻时，我的意思是表明所说的并非实际情况，而是真理透过寓言的迷雾隐约显露。然而，我要列出奥力振第三卷中实际的话：我们可以说，灵魂爱那将要看见天主救恩的肉身，它保养顾惜它，用纪律训练它，用天上的食粮满足它，给它喝基督的血；这样它就能因健康的食物而强壮，毫无软弱地跟随她的丈夫。灵魂保养顾惜身体，如同基督保养顾惜教会，这是一个美好的比喻，因为基督曾对耶路撒冷说：\par

\BibleQuote{我多次愿意聚集你的儿女，好像母鸡把小鸡聚集在翅膀底下，只是你们不愿意；}这样这必朽的将穿上不朽的，并且像有了翅膀一样轻盈地更容易升到空中。因此，让我们男人保养顾惜我们的妻子，也让我们的灵魂保养顾惜我们的身体，其方式要使妻子变成男人，身体变成精神，没有性别的差异；因为天使中既没有男性也没有女性，同样，我们——将要成为天使的人——可以在这里开始成为我们被应许在天上将要成为的。\par

29. 在此之前，我以这些话陈述了我自己对这段经文所做的简单解释：\par

\Quote{按照经文的字面意义，我们在夫妇彼此善待的诫命之后，领受了保养顾惜妻子的命令：也就是说，要供给她们所需的食物和衣服。}\par

这是我对这段经文的个人理解。因此，我的话意味着，随后一切可能被拿来反对我的论述，都必须被理解为不是表达我自己的观点，而是表达我对手的观点。但有人可能认为，我对这段经文难题的解决过于简短和武断，并且，按照前面所说的，它用寓言的黑暗包裹了真实含义，从而将其从真实意义降低为不那么真实的意义。因此，我将更贴近这个问题，问一问在另一种解释中有什么是你需要不同意的。我想是这样的：我说过，灵魂应该爱护身体，如同人们爱护妻子，以便这可朽坏的穿上不朽坏的，并且，如同轻松地安放在翅膀上，它能更容易地升入空中。当我说这可朽坏的必须穿上不朽坏的时，我并没有改变身体的本质，而是赋予它在存有等级中更高的地位。同样，对于随后所说的，如同轻松地安放在翅膀上，它能更容易地升入空中：得到翅膀的人，即得到不死，以便更轻快地飞向天堂，并没有停止成为他曾经所是的。但你会说，我被以下的话惊呆了：\par

\Quote{那么，让我们男人爱护我们的妻子，让我们的灵魂爱护我们的身体，这样妻子就变成男人，身体变成精神，不再有性别差异；而是如同在天使中既没有男性也没有女性，我们这些将要像天使一样的人，可以在世上开始成为在天上所应许的那样。}\par

你有理由感到震惊，如果我没有在前面的话之后说过：\Quote{我们可以在天上开始成为所应许的那样。} 当我说：\Quote{我们将要在世上开始成为}时，我并没有消除性别差异；我只是消除了情欲和性交，正如宗徒所说：\BibleQuote{时间是短暂的；从此以后，那有妻子的，要像没有一样；} 也如主所暗示的，当回答那女人会成为七兄弟中哪一个的妻子时，他回答说：\BibleQuote{你们错了，因为你们不明了圣经，也不明了天主的能力；因为复活的时候，他们也不娶也不嫁，而是像天上的天使一样。} 事实上，当男女之间保持贞洁时，就开始真正没有男性也没有女性了；虽然生活在身体中，他们正在变成天使，在天使中既没有男性也没有女性。同一位宗徒在另一处也说：\BibleQuote{你们凡是受洗归于基督的，就是穿上了基督。不再分犹太人或希腊人，不再分奴隶或自由人，不再分男性或女性，因为你们众人在基督耶稣内已是一个。}\par

% structure: p0094 source=h2 target=subsection reason=relative-heading-rank; base=h2

\subsection{对于阅读世俗书籍的指控，我答复说，我仍记得年轻时所学。}

30. 然而现在，既然我的辩护已从这些崎岖破碎之处驶出航程，并且我通过坦然面对控告者而驳斥了他对我提出的异端指控，我将转而讨论他用以来攻击我的其他指控。第一项指控是，我是一个粗鄙的人，一个诽谤所有人的人；我总是在对前辈们咆哮和撕咬。我请他指名道姓说出一个我诋毁其名誉的人，或者一个按照我对手所惯用的伎俩，我通过虚假的赞扬而刺伤的人。但是，如果我针对居心不良的人说话，用我的笔尖刺伤某个\Person{卢修斯·拉努维努斯}或\Person{阿西尼乌斯·波利奥}——科内利家族的人，如果我击退一个夸夸其谈、好奇心强的人的攻击，并把我所有的箭矢都射向一个目标，那他为什么要把专为他一人准备的伤口分散给其他人呢？他又为何如此不明智，通过他对我的攻击所作出的恼怒回答，表明他意识到正是他一人适合这顶帽子？\par

他同时指控我发假誓和亵渎，因为在一本为教导一位基督童贞女而写的书中，我描述了自己有一次在梦中站在审判者的法庭前所作出的承诺，即我再也不阅读世俗文学，而且尽管如此，我有时仍然提到我当时所谴责的学识。我想我在这里遇到了那个人，他以\Person{萨卢斯蒂安·卡尔普尔尼乌斯}的名义，并通过演说家\Person{马格努斯}写给我的信，提出了一个并不太大的问题。我对整个问题的回答包含在我当时写给他的一篇短文中。但现在我必须就我梦中的亵渎和发假誓作出回答。我说过从此不再读世俗书籍：那是对未来的承诺，并非抹去对过去的记忆。你可能会问我，你怎么能记住你这么久没有读过的东西呢？我必须通过回顾这些旧书之一来回答：\par

\Quote{年轻时的习惯最为牢固。}\par

但我这样否认，其实是控告了自己：我搬出\Person{维吉尔}作证，却被我自己的辩护人指控。我想我必须织出一长串言词，来证明每个人自己所意识到的事。我们中间有谁不记得自己的幼年呢？虽然您是极其严肃的人，我却要使您发笑；您最终将不得不像\Person{克拉苏}那样——据\Person{卢齐利乌斯}说，他一生只笑过一次——如果我讲述我童年的记忆：我如何在仆役工作的各个房间之间奔跑；我如何在游戏中度过假期；或者，我如何像俘虏一样被从我祖母的膝上拖去上愤怒的\Person{奥尔比留斯}的课。您可能更加惊讶，如果我说，即使现在我已头发灰白秃顶，我常在梦中看到自己，一个卷发的青年，穿着托加袍，在修辞学老师面前朗诵一篇有争议的论题；当我醒来时，我庆幸自己逃脱了发表演说的危险。请相信我，我们的幼年非常准确地给我们带回许多东西。如果您受过文学教育，您的心智会保留它最初浸染的东西，就像酒桶保留它的气味一样。羊毛上的紫色染料无法用水洗掉。甚至驴子和其他牲畜，无论旅程多长，都认得它们以前停过的客栈。您自己不需要老师就学会了希腊文，却对我没有忘记我的拉丁文书籍感到惊讶吗？我学习了\OriginalTerm{逻辑学要素}{Elements of Logic}中七种三段论形式；我学习了\OriginalTerm{公理}{Axiom}（或可称为\OriginalTerm{确定命题}{Determination}）的意义；我学习了每个句子必须有一个动词和一个名词；我学习了如何堆积\OriginalTerm{堆垛悖论}{Sorites}的步骤，如何识破\OriginalTerm{说谎者悖论}{Pseudomenos}的巧妙转弯以及常见的诡辩的欺诈。我可以发誓，离开学校后我从未读过这些东西中的任何一样。我想，为了避免将我所学的东西当作罪行来指控，我必须按照诗人的寓言去喝\Place{忘川}之水。我召唤您，您因我的学识浅薄而指责我，并自认为是一位\OriginalTerm{文人}{litérateur}和\OriginalTerm{拉比}{Rabbi}，告诉我，您怎么敢写出您所写的一些东西，并以与那位最雄辩之人\Person{格列高利}同等的雄辩光彩翻译他？您从哪里获得那种词语的流畅、陈述的清晰、翻译的多样性——您年轻时几乎只有修辞学的初步尝试？如果我大错特错，您一定在私下研读\Person{西塞罗}。我怀疑，像您这样有教养的人，您禁止我阅读\Person{西塞罗}并加以惩罚，以便您独自在我们教会的作家中夸耀您流畅的雄辩。但我还是要说，您似乎更追随哲学家，因为您的风格类似于\Person{克莱安西斯}那些刺人的句子和\Person{克吕西普}的曲折，这不是出于任何技巧——因为您说您对那技巧一无所知——而是出于天才的共鸣。斯多葛派主张逻辑学为自己的所有物，而您蔑视这门科学，视之为愚蠢；因此，在这一方面，您是伊壁鸠鲁派；您雄辩的原则不是风格，而是内容。因为，确实，如果您只为自己的朋友写作，而不是为所有人，那么别人是否理解您想说的话又有什么关系呢？我必须承认，我自己并不总是理解您所写的东西，而且我认为自己是在读\Person{赫拉克利特}；然而我不抱怨，也不为我自己的迟钝而感到遗憾；因为阅读您所写的东西的麻烦，并不比您写它的麻烦更大。\par

% structure: p0099 source=h2 target=subsection reason=relative-heading-rank; base=h2

\subsection{同样，梦中应许不可强求。老朋友之间为何要翻出这类事情互相攻击？}

31. 即使这是一个完全清醒时做出的承诺，我也完全可以像刚才那样回答。但这是一件新的无耻之事：他用一个梦来指责我。我该如何回答？我没有时间思考自己领域之外的事情。我希望那些聚集在此处的群众和从世界各地汇聚而来的基督徒们没有使我无法阅读圣经。然而，当一个人把梦当作罪行时，我可以引用先知的话对他说：我们不可相信梦；因为即使梦见通奸也不会定我们下地狱，梦见殉道的冠冕也不会提升我们上天堂。我常在梦中看见自己死去并被埋葬；我常在梦中飞越大地，像游泳一样浮在空中，越过山岭和海洋。因此，我的控告者可能会要求我停止生存，或者要我肩上长翅膀，因为我的思维在睡眠中常被这种模糊的幻象所愚弄。有多少人在睡梦中是富人，醒来却发现自己成了乞丐！或者喝凉水解渴，醒来时喉咙干渴灼热！你要求我履行一个梦中给出的承诺。我要以一个更真实、更切近的问题来回应你：你在圣洗时承诺的一切，你都做到了吗？你或我是否实现了隐修士身份所要求的一切？我请求你，想想你是否在看着自己眼中的大梁时，却在挑剔我眼中的木屑？我是不情愿说这些话的；是悲伤迫使我勉为其难地说出。至于你，你捏造我清醒时的行为还不够，还要因我的梦来控告我。你对我的行为如此感兴趣，以至于要讨论我在睡眠中说了什么或做了什么。我不想详述你如何因热衷攻击我而玷污了你的身份，以及你如何用言语和行为尽一切可能侮辱了整个基督徒团体。但我郑重警告你，并且会一再重复。你正在攻击一个长着角的生物；如果不是因为我牢记宗徒的话：\BibleQuote{诽谤者不能继承天主的国}，以及\BibleQuote{你们彼此仇恨，彼此消耗}，我会让你感受到你在轻微而虚假的和解之后激起了多么巨大的不和。你在朋友和陌生人面前堆积对我的诽谤，对你有什么好处？难道因为我不是奥利振主义者，不相信我在天上犯了罪，就因此在地上被指控为罪人吗？我们重新和好的结果，难道是要我不再反对异端，以免我对他们的注意被当作对你的攻击吗？只要我不拒绝被你赞美，你就以导师的身份追随我，称我为朋友和兄弟，并承认我在各方面都是大公的。但当我请求你免去赞美，并自认不配有这样一位伟人为我吹嘘时，你立刻用笔划掉你所写的，开始辱骂你以前赞美过的一切，并同一张嘴吐出甜苦两种言语。但愿你能理解我在克制自己，没有让言语配合胸中的怒火；并且我像圣咏作者一样祈祷：\BibleQuote{上主，求你在我的口前设警卫，把守我嘴唇的门户。不要使我的心偏向恶言；}以及他在别处说的：\BibleQuote{当恶人在我面前时，我缄默不语，自卑自谦，甚至连好话也不出口；}又说：\BibleQuote{我成了好像一个听不见的人，口中没有辩驳。}但为我申冤的上主必要回答，正如他藉先知所说的：\BibleQuote{报复在我，我必报应，上主说}；又在另一处：\BibleQuote{你坐着毁谤你的兄弟，诬蔑你母亲的儿子。你做了这些，我默不作声；你便以为我像你一样；但我要谴责你，并将这些摆在你眼前；}好叫你看见自己犯下了你诬告他人所犯的罪。\par

% structure: p0101 source=h2 target=subsection reason=relative-heading-rank; base=h2

\subsection{我认为所有罪都在圣洗中得到赦免，我的主张是正确的。}

32. 据说，再说一点，他的一个追随者\Person{克里索戈努斯}挑剔我说，在圣洗中所有的罪都被除去，而对于那个两次结婚的人，说他已死去并在基督内复活成为一个新人；并且还有几个这样的人在教会中担任主教。我要给他一个简短的答复。他和他朋友们手中有我的信，正是为此他们指责我。让他来答复这信，让他用自己的推理推翻它的推理，用他的著作证明我的著作是错误的。他为什么皱起眉头，收缩并皱起鼻孔，掂量他空洞的话语，在普通人中间假装一种与他的行为相矛盾的圣德？让我再一次在他的耳边宣告我的原则：旧亚当在圣洗的水中完全死去，然后一个新人同基督一起复活；属土的人灭亡，属天的人被举起来。我说这话，并不是因为我本人通过基督的仁慈对这个问题有特别的兴趣，而是因为我回答了我的弟兄们询问我的意见，既不是打算为别人规定他们可能认为应该相信什么，也不是要用我的意见推翻他们的决定。因为我们这些隐藏在斗室中的人并不觊觎主教的职位。我们不象一些人，他们藐视一切谦卑，急于用黄金购买主教职；我们也不愿意以叛逆的心理去压制天主所选的主教；我们也不通过偏袒异端来表明我们自己就是异端。至于金钱，我们没有，也不想要。\BibleQuote{只要有吃有穿，就应当知足。}同时，我们不断吟唱描述那将上主的山的人的话：\BibleQuote{不放债取利，也不受贿赂加害无辜；这样行事的人，永不动摇。}我们还可以说，做相反这些事的人将永远跌倒。\par

这最后一章中几乎每一句话都是对\Person{鲁菲努斯}的恶意影射。他的\Quote{皱起的眉头}和\Quote{翘起的鼻子}，他掂量字句，他所谓的财富，都在其他地方被提到，尤其是在\Person{热罗尼莫}致\Person{鲁斯蒂古斯}的信中（第125信，第18章）对他死后所作的讽刺性描述中。\par

\SourceRecord{Translated by W.H. Fremantle.；Nicene and Post-Nicene Fathers, Second Series；Vol. 3.；Edited by Philip Schaff and Henry Wace.；Buffalo, NY: Christian Literature Publishing Co.,；1892.；<http://www.newadvent.org/fathers/27101.htm>.}

% structure: p0001 source=h1 target=section reason=page-title

\section{驳斥鲁菲努斯辩护辞（卷二）}

% structure: p0002 source=h2 target=subsection reason=relative-heading-rank; base=h2

\subsection{对鲁菲努斯致阿纳斯塔修斯辩护辞的批评。他为自己不来罗马所找的借口荒谬可笑。他的父母已去世，旅途亦甚为便捷。此前从未有人听说过他曾因信德而被囚禁或流放。}

1. 至此，我已回答了关于我的罪行的指控，确切地说，是为我的罪行辩护——这些指控是那位狡诈的颂扬者先前对我提出的，而他的门徒至今仍在不断强调。我的回应虽未达到应有的水准，但已是尽力而为；我克制了怨言，因为我的目的与其说是控告他人，不如说是为自己辩护。现在，我将转而讨论他的辩护辞，他在其中试图向罗马城主教阿纳斯塔修斯表明自己的清白，并且为了自我辩护，对我罗织了大量诽谤之词。他对我的爱，恰似一个被风暴卷起、在深水中险些淹死的人，对那位他所紧抓其脚踝的强壮游泳者的爱：他决意要我与自己一同沉浮。\par

2. 首先，他宣称自己是在回应罗马方面对其正统性的质疑，而他本人是一个在神圣信德与爱德方面都得到充分认可的人。他说他本愿亲自前来，只是因为他刚在阔别三十年后回到了父母身边，如果这么快就离开他们，似乎显得冷酷无情；此外，他也因长途跋涉而身体欠佳，难以再次劳顿。既然他本人无法前来，他便将他的辩护辞作为一根文字棍棒送来，希望主教手持它驱赶那些向他狂吠的恶犬。如果他在信德与爱德上为众人所认可，尤其是为他所写信的主教所认可，那么他在罗马为何会遭到攻击和辱骂，而关于他名誉受损的传闻却日益增多？再者，一个人自称在神圣信德与爱德上为众人所认可，这是何等谦卑？宗徒们曾祈祷：\Quote{主，请增加我们的信德，}得到的回答是：\Quote{如果你们有信德如同芥菜籽一样大；}甚至对伯多禄也说：\Quote{小信德的人哪，你为什么怀疑？}至于爱德，它比信德和望德都大，保禄说他期盼而非自以为拥有；没有爱德，即使殉道时流血、身体被焚，也得不到冠冕。然而，我们的朋友却将这两者都据为己有：即便如此，仍有生物向他吠叫，并且会继续吠叫，除非那位杰出的教宗用他的棍子将它们驱走。但他提出的那个借口——阔别三十年后回到父母身边——是何等荒谬！他既无父亲也无母亲！他年轻时离开时他们尚在人世，如今他年迈了，反而思念已故的双亲。不过，也许他所说的\Quote{父母，}是指士兵和百姓口中所说的亲属和亲戚；他说他不愿被认为如此冷酷无情，竟至于抛弃他们；因此他离开家园，住在阿奎莱亚。他那备受认可的信德在罗马处境危险，然而他却仰面躺着，因为三十年来有些疲倦，竟不能乘坐马车沿那弗拉米尼亚大道完成极为便捷的旅程。他以长途跋涉后的疲倦为借口，仿佛三十年来他一直在奔波，或者仿佛在阿奎莱亚休息两年后，他仍因过去的旅途劳顿而疲惫不堪。\par

3. 我将触及其余各点，并抄录他信中的原话：\par

\Quote{尽管我的信德已在异端迫害时期得到证明，当时我生活在亚历山大的神圣教会中，因信德而遭受囚禁和流放。}\par

我仅诧异他为何没有加上\Quote{耶稣基督的囚犯}，或\Quote{我被从狮子的口中救出}，或\Quote{我在亚历山大与野兽搏斗}，或\Quote{我已跑完赛程，保守了信德。从此以后，正义的冠冕已为我预备好了。}他所描述的这些流放和囚禁是什么？我为这赤裸裸的谎言感到脸红。难道囚禁和流放无需经过司法判决就会执行吗？我希望看到他列出这些囚禁和他说他被迫流放的各种省份的清单。然而，似乎有多次囚禁和无数次流放，以至于他至少可以举出其中一次。请让他拿出他作为宣信者的行为记录吧，因为此前我们对这些一无所知；这样我们就能满意地将他的事迹与亚历山大的其他殉道者一同诵读，他也能以这些话来面对那些向他狂吠的人们：\Quote{从今以后，但愿没有人再烦扰我，因为我身上带着主耶稣基督的印记。}\par

% structure: p0008 source=h2 target=subsection reason=relative-heading-rank; base=h2

\subsection{他的信仰宣认不能令人满意。无人询问他关于天主圣三的事，而是关于奥利振关于复活、灵魂起源及撒殚可得救的学说。关于复活和撒殚，他语焉不详。关于灵魂，他声称自己无知。}

4. 他接着说：\par

\Quote{尽管如此，或许有些人愿意证明我的信德，或者愿意聆听并学习我的信德，我将宣告我对天主圣三的看法如此：}\par

等等。起初你说，你把你的信德托付给主教，如同交给他一根棍子，好使他藉此为你抵御那些吠叫的狗。现在你说话就不那么自信了：\Quote{可能有些人希望考验我的信德。}当如此众多的吠叫声传到你的耳中时，你便开始犹豫。我不打算讨论你所使用的措辞形式，因为这些你轻视并谴责：我只按意义回答。有人问你这事，你却为另一件事为自己辩护。关于亚略的教义，你早在亚历山大里亚以监禁和流放，不是用言语，而是用鲜血与它们争战。但现在的问题涉及奥力振的异端，以及因此事激起的反对你的情绪。我很遗憾你竟然费心去医治已经愈合的创伤。你宣认同一神性中的天主圣三。现在普世都宣认这一点，我认为连魔鬼也宣认天主子由童贞玛利亚诞生，并取了属于人性的肉体和灵魂。但我恳求你不要认为我是一个好争论的人，如果我更严格地考问你。你说天主子取了属于人性的肉体和灵魂。那么，我问你，请不要对我生气，而要回答这个问题：耶稣所取的那个灵魂，在它由玛利亚诞生之前就已存在吗？它是与身体一起在原始的童贞本性中受造，而那本性是由圣神所生的吗？或者，当身体已在母腹中形成时，它才被立即创造，并从天上降下？我希望知道你认为哪一种。如果它在由玛利亚诞生之前就已存在，那么它就还不是耶稣的灵魂；它是以某种方式被使用，并因其德行的报酬而成了他的灵魂。如果它是由于传衍产生的，那么人类的灵魂（我们相信是永恒的）就与兽类的灵魂处于同一状况，兽类的灵魂与身体一同消亡。但如果它是在身体形成之后被创造并送入身体的，就直截了当地告诉我们，使我们不再焦虑。\par

5. 这些答案你一个也不会给我们。你转向其他事情，用你的诡计和言辞展示阻止我们仔细关注问题。什么！你会说，问题不是关于肉身的复活和魔鬼的惩罚吗？确实如此；因此我要求一个简短而真诚的答案。我不对你所说的正是我们生活的这个肉身复活，不失任何肢体，不切除身体任何部分（因为这是你亲口的话）提出质疑。但我想知道你是否相信——奥力振所否认的——身体以与死亡时相同的性别复活；玛利亚仍是玛利亚，若翰仍是若翰；还是性别会如此混合和混淆，以致没有了男人或女人，而是某种既是又不是的东西；以及你是否相信身体保持不朽坏和不死，并且，正如你精明地追随宗徒所暗示的，是永远属灵的身体；不仅是身体，而且是实际的肉身，其中有血液注入，通过管道穿过血管和骨骼——正如多默所触摸的那种肉身；还是逐渐分解为无，还原成组成它们的四种元素。这些你应当要么承认，要么否认，而不是说出奥力振也说过的话，但却不真诚，仿佛他在愚弄傻瓜和孩子的软弱一般：\Quote{不失任何肢体，不切除身体任何部分。} 你以为我们所担心的是我们会没有鼻子和耳朵复活，会发现我们的生殖器官被切除或残废，在新耶路撒冷建起一座阉人城吗？\par

6. 关于魔鬼，他如此陈述他的意见：\par

\Quote{我们也宣认将来的审判，在那审判中，每个人都要按他所行的善或恶，从他肉身的生活中领受应得的报应。而且，如果在人的情况下，赏报是按他们的行为而定，那么对于作为罪恶普遍原因的魔鬼，岂不更是如此呢？至于魔鬼本身，我们的信仰正如福音所载，即他和他所有的天使都要承受永火为他们的份，并与他们一起那些行他作为的人，也就是那些成为弟兄控告者的人。因此，若有人否认魔鬼应受永火，就愿他与魔鬼一同在永火中有份，好使他亲身经历他现在所否认的事实。}\par

我将逐字重复。\Quote{我们也宣认将来的审判，在那审判中等等。} 我本来决定对言语上的缺陷不置一词。但是，既然他的门徒仰慕他们老师的口才，我就对之略加批评。他已经说过\Quote{将来的审判}；但是，作为一个谨慎的人，他害怕只说\Quote{在其中}，因此写了\Quote{在那审判中}；恐怕若他没有第二次说\Quote{审判}，我们会忘记前面的话，而补上\Quote{驴}这个词。他后来所引入的\Quote{那些成为弟兄控告者的人将与他一同在永火中有份}，其风格也同样优美。谁曾听过\Quote{拥有火焰}？这就像\Quote{享受酷刑}一样。我想，他现在已是希腊人，曾试图自己翻译，因此对于\OriginalTerm{将承受}{\GreekText{κληρονομήσουσιν}}这个词——它可以用一个拉丁词\OriginalTerm{将继承}{Hæreditabunt}来表达——他用了\OriginalTerm{将获得继承}{Hæreditate potientur}，以为这样更精致华丽。他的整篇言论充满了这种琐碎和不当的措辞。但让我们回到他话的意思上。\par

7. 继续：\par

\Quote{这是一柄长矛，刺穿了魔鬼——\InnerQuote{他是罪恶的普遍原因}——如果他像人一样要为自己的行为交账，并且\InnerQuote{与他的天使一同承受永火的产业}。毫无疑问，这正是他所缺乏的：在使人类陷入痛苦之后，他自己应当\InnerQuote{承受永火}，那是他一直渴望的。}\par

在我看来，你对魔鬼的言辞似乎有些过于严厉，以不实的指控攻击那控告众人的。你说：\Quote{他是罪的普遍原因。}你既将他定为一切恶行的元凶，又免除了人的过失，取消了自由意志。我们的主说：\BibleQuote{从心里发出来的，有恶念、凶杀、奸淫、苟合、偷盗、妄证、毁谤。}关于犹大，我们在\BibleBook{若望}{福音}中读到：\BibleQuote{耶稣蘸了饼以后，撒殚就进入了他。}这是因为在蘸饼以前，他已自愿犯罪，并未因自卑或救主的容忍而悔改。宗徒也说：\BibleQuote{这样的人我已交给撒殚，使他们受教训不再亵渎。}他将那些人交给撒殚如交给行刑者，以惩罚他们；在他们被交给他以前，他们已凭自己的意志学会了亵渎。\Person{达味}也用简短的言语区分了出于自身意志的过犯和罪恶的诱惑，他说：\BibleQuote{求你也赦免我隐而未显的罪愆，求你不要让骄傲的权势辖制我。}在\BibleBook{训道}{篇}中也读到：\BibleQuote{如果有官长起来攻击你，不要擅离岗位。}由此我们清楚看到，如果我们给那起来攻击的权势以机会，如果我们不将那攀爬我们城墙的敌人一举摔倒，我们就会犯罪。至于你威胁你的弟兄们，即那些控告你的人，要他们与魔鬼同受永火，在我看来，你与其说是拖你的弟兄们下去，不如说是提升了魔鬼，因为按你的说法，他仅与基督徒同受同样的火刑。但我想你很清楚，按照\Person{奥力振}的观念，永火是什么意思，即罪人的良心和他们内心折磨的悔恨。他认为这些意义包含在\BibleBook{依撒意亚}{书}的话中：\BibleQuote{他们的虫子不死，他们的火也不灭。}以及向巴比伦说的话：\BibleQuote{你有炭火，你要坐在其上，这些要作你的帮助。}同样，在\BibleBook{圣咏}{集}中对忏悔者说：\BibleQuote{有什么能给你，有什么能加给你，以对抗诡诈的舌头？勇士的利箭，加上荒野的炭火。}按他的解释，这意味着天主旨意的箭（关于这箭，先知在别处说：\BibleQuote{我生活在痛苦中，有刺扎我。）}要创伤并击穿诡诈的舌头，在其中除去罪恶。他还解释主证言说：\BibleQuote{我来是为把火投在地上，我是多么希望它已经燃烧起来！}其意思是：我愿意众人都悔改，通过圣神烧尽他们的恶习和罪恶；因为我是那经上所写的：\BibleQuote{我们的天主是吞噬之火。}由此可见，这样谈论魔鬼不是什么了不起的事，因为这火也是为人预备的。如果你要避免被怀疑相信魔鬼得救，你更应该这样说：\BibleQuote{你已成了毁灭，将永远不再存留；}正如主对\Person{约伯}论及魔鬼时所说：\BibleQuote{看，他的希望要落空，在众人面前他要被摔倒。我不激发他，如同激发一个残忍者，因为谁能抗拒我的面目？谁先给了我，使我偿还？因为天下万物都是我的。我不顾惜他和他那强大而善于平息忿怒的言语。}因此，这些事可能被视为一个质朴之人的作品。它们的含义对于明白这些事的人是显而易见的；但对于无知的人，它们可能显得天真无邪。\par

8. 但接下来关于灵魂状况的论述，无论如何也不能原谅。他说：\par

\Quote{我接着被告知，关于灵魂的本性问题已引起一些骚动。是否应当受理这类申诉而非搁置一边，须由您自己决定。然而，如果您愿意知道我对这个问题的看法，我将坦率陈述。我读过许多作者关于这个问题的著作，发现他们的意见各异。我所读过的作者中，有些人认为灵魂是通过人类种子的途径与物质身体一同注入的，并尽其所能力证这一点。我认为这是拉丁作者中\Person{戴尔都良}或\Person{拉克唐提乌斯}的意见，或许还有少数其他作者。另一些人断言，天主每天都在创造新的灵魂，并注入那些在子宫中形成的身体；还有一些人相信，灵魂都是很久以前被造的，那时天主从无中创造了万物，他现在所做的只是按照他的意愿派遣每个灵魂降生在身体中。这是\Person{奥力振}和某些希腊作者的看法。至于我自己，我在天主面前声明，在读过这些意见之后，我无法将其中任何一个视为确定无疑的；这个问题的真理的断定我留给天主，以及他愿意启示的任何人。因此，我在这点上的表白是：首先，这些不同的意见是我在书籍中找到的；其次，我至今仍对此无知，除了这一点：教会信理宣示天主是灵魂和身体的创造者。}\par

% structure: p0021 source=h2 target=subsection reason=relative-heading-rank; base=h2

\subsection{什么样的拉丁语！可怜的灵魂，必被他的野蛮折磨。}

9. 在我进入这段文字的主题之前，我必须钦佩\Person{泰奥弗拉斯托斯}式的措辞：\par

\Quote{我被告知，关于灵魂的本性问题已引起一些骚动。是否应当受理这类申诉而非搁置一边，须由您自己决定。}\par

如果这些关于灵魂来源的问题在罗马引起了争论，那么对于是否应该接纳这些问题（这完全属于主教的裁量权）的抱怨和怨言又有什么意义呢？但他或许认为问题和抱怨是同一回事，因为他发现这种说法出现在\WorkTitle{卡珀注释}中。接着他写道：\Quote{我读过的某些人认为，灵魂是通过人类种子的渠道与物质身体一同注入的；对于这些，他们尽其所能地给出了证明。}这里的说法是多么随意啊！语气和时态是多么混杂啊！\Quote{我读过一些说法——他们用尽他们所能的论断来证实它们。}接下来：\Quote{另一些人断言，天主每天都在创造新的灵魂，并将它们注入在母腹中已成形身体里；还有一些人相信，灵魂早已在万物之初由天主从无中创造而成，祂现在所做的只是将每一个灵魂按其旨意派遣出去，在它自己的肉身中出生。}这里也有一个非常优美的安排。他说，有些人断言这样那样；有些人宣称灵魂早已被造，即当年天主从无中创造万物时，祂现在按祂的喜好将它们派遣出去，在它们自己的肉身中出生。他讲得如此令人不快、如此混乱，以至于我纠正他的错误比他写这些错误更加费劲。最后他说：\Quote{然而，我虽然读过这些事；}而且，句子还未完结，他仿佛提出了什么新东西一样补充道：\Quote{然而，我不否认我读过每一件这样的事，并且至今承认我是无知的。}\par

% structure: p0025 source=h2 target=subsection reason=relative-heading-rank; base=h2

\subsection{你不可对这样一件众教会都知道的事无知。}

10. 不幸的灵魂！被所有这些粗野的言辞刺穿，如同被许多长矛刺中！我怀疑，即使按照奥力振的错误理论，他们从天上坠落到地上并穿上这些粗重的身体时，所受的麻烦，也不及现在被这些奇怪的词句四面击打来得更多：更不必说那个不祥的词，它说灵魂是通过人的种子被注入的。我知道在基督徒著作中，通常不批评纯风格的错误；但我认为通过几个例子表明，讲授你所不知道的事、写你所不懂的东西是何等鲁莽，是好的；这样，当我们进入主题时，我们才能准备好发现同等程度的智慧。他寄了一封信，他称之为一根非常结实的棍子，作为武器送给罗马主教；而正是关于狗们对他吠叫的这个主题，他声称自己对这个问题完全无知。如果他对引起他恶评的主题无知，他何必寄送一份辩白书，其中没有为他辩护，只是承认他的无知？这种做法只会滋生怀疑，而不是平息怀疑。他提出了关于灵魂来源的三种意见；最后得出结论：\Quote{我不否认我读过每一种意见，我承认我仍然无知。}你会以为他是阿尔克西劳或卡涅阿德斯，他们宣称没有确定性；尽管他在谨慎方面甚至超过他们；因为他们因将一切真理从人类生活中剔除而在哲学家中激起了无法忍受的恶意，被迫发明了概率论，以便通过他们的可能断言来调和他们的不可知论；但他仅仅说他是不确定的，不知道这些意见中哪一个是真的。如果这就是他所能回答的全部，是什么促使他召唤如此伟大的教宗来见证他缺乏神学修养？我猜想这就是他告诉我们的那种倦怠，他因三十年的旅程而筋疲力尽，无法来到罗马。有许多事情我们大家都不知道；但我们不为自己无知寻求见证。至于圣父、圣子和圣神，关于我们的主和救主的诞生，关于依撒意亚呼喊说：\BibleQuote{谁能述说他的世代？}他大胆发言，并且一个过去所有时代都不知道的奥秘，他声称完全了解：唯独这一件事他不知道，而不知道这件事就会使人跌倒。至于童贞女如何成为天主之母，他完全知道；至于他自己如何出生，他一无所知。他承认天主是灵魂和身体的创造者，无论灵魂是先在身体之前存在，还是与身体的胚芽同时产生，或者是在子宫中已成形后才被送入。在任何情况下，我们都承认天主是它们的作者。问题不在于灵魂是由天主还是由另一个创造的，而在于他陈述的三个意见中哪一个是真的。对此他声称无知。当心！你可能会发现人们说，你承认对三者无知的原因，是你不想被迫谴责其中之一。你饶恕\Person{戴尔都良}和\Person{拉克坦提乌斯}，以便不与他们一起谴责\Person{奥力振}。据我所记（虽然我可能记错），我不记得读过拉克坦提乌斯说过灵魂是与身体同时植入的。但是，既然你说你读过，请告诉我可以在哪本书中找到，这样你就不会被以为在他死后诽谤了他，就像你在我睡觉时诽谤了我一样。然而即使在这里，你也以谨慎和犹豫的步伐行走。你说：\Quote{我认为，在拉丁人中，戴尔都良或拉克坦提乌斯持这种意见，也许还有其他人。}你不仅对灵魂的起源有疑问，而且对于每位作者持什么意见，你也只是\InnerQuote{认为}；然而这件事有些重要性。不过在灵魂的问题上，你公开宣称无知，并承认你未受教导；关于作者，你的知识仅限于\InnerQuote{认为}，几乎不能\InnerQuote{假定}。但唯独对奥力振，你很明确。你说：\Quote{这是奥力振的意见。}但是，我问你：这个意见正确与否？你回答说：\Quote{我不知道。}那么，你为什么给我派来不断来往的信使和送信人，仅仅是为了告诉我你很无知？为了防止我怀疑你的无知是否如你所说的那样严重，并认为你有可能狡猾地隐藏你所知道的一切，你在天主面前发誓，直到此刻你对这个问题还没有任何确定和明确的看法，你把知道真理的事留给天主，也留给任何他愿意启示的人。什么！难道在你看来，通过所有这些时代，没有一个人值得接受这个启示吗？无论是圣祖、先知、宗徒、还是殉道者？你自己住在王子和流放者中间时，这些奥秘不是也向你启示了吗？主在福音中说：\BibleQuote{父啊，我已将你的名启示给人们。}那位启示了父的，难道对灵魂的起源保持沉默吗？如果你的兄弟们因你发誓你不知道基督教会声称知道的事而惊骇，你岂不惊奇？\par

% structure: p0027 source=h2 target=subsection reason=relative-heading-rank; base=h2

\subsection{关于翻译\OriginalTerm{论原理}{\GreekText{Περὶ} \GreekText{᾿Αρχῶν}}，这不是一个问题，而是一个指控，即你无理地篡改了那本书。}

11. 在阐述了他的信仰，或者更确切地说，他的无知之后，他转入另一个问题；并试图为他将\OriginalTerm{论原理}{\GreekText{Περὶ} \GreekText{᾿Αρχῶν}}翻译成拉丁文找借口。我将逐字记下他的话：\par

\Quote{我听说，有人对我提出反对，因为我应几位弟兄的请求，将奥力振的某些作品从希腊文译成了拉丁文。我想每个人都会明白，这不过是出于恶意才被当作指责的借口。因为，如果原作者有什么令人生厌的言论，这为何要转嫁成译者的过错呢？有人请我将希腊文本中已有的内容用拉丁文呈现出来；我所做的不过是把拉丁词汇与希腊思想对应起来。因此，如果这些思想中有什么值得称赞，那称赞并不属于我；同样，如果有什么该受指责，也与我无关。}\par

\Quote{我听说，}他说，\Quote{从此\Emph{纷争}产生了。}这多么巧妙，将其实是对他的指控说成是纷争。\Quote{我应弟兄们的请求，将奥力振的某些言论译成了拉丁文。}是的，但这些\Quote{\Emph{某些言论}}是什么？难道没有名字吗？你沉默吗？那么指控者提出的诉状就会替你说话。\Quote{我想，}他说，\Quote{每个人都会明白，这不过是因为嫉妒才被当作指责的借口。}什么嫉妒？人们嫉妒你的口才吗？还是你做了别人从未能做的事？我在这里，翻译了许多奥力振的作品；然而，除你之外，没有人对此向我表示嫉妒或诽谤。\Quote{如果原作者有什么令人生厌的言论，为何要转嫁成译者的过错呢？有人请我将希腊文本中已有的内容用拉丁文呈现出来；我所做的不过是把拉丁词汇与希腊思想对应起来。因此，如果这些思想中有什么值得称赞，那称赞并不属于我；同样，如果有什么该受指责，也与我无关。}你还能对人们对你怀有恶感感到惊讶吗？当你对公开的亵渎言论仅仅说\Quote{如果原作者有什么令人生厌的言论}？那些书中所说的令所有人厌恶；而你独自一人怀疑并抱怨这被\Quote{转嫁成译者的过错}，而你在序言中却赞扬了它们。\Quote{有人请我将希腊文本中已有的内容译成拉丁文。}我希望你真的做了你声称被请求的事。那样你就不会成为任何恶意的对象。如果你保持了译者应有的忠实，我就不必用我的真实翻译来对抗你的虚假翻译。你自己的良心知道你添加了什么，删除了什么，以及你凭自己的判断在一方面或另一方面改变了什么；在此之后，你竟敢告诉我们，好坏不应归于你，而应归于原作者。你通过再次缓和你的言辞，显示出你对激起恶意的感觉；如同你行走在空气或谷穗尖上，你说：\Quote{\Emph{无论}这些意见中有称赞或指责。}你不敢为他辩护，但也不愿定他的罪。请选择两者中任何一个；选择权在你；如果你所翻译的是好的，就称赞它；如果是坏的，就定它的罪。但他却找借口，并编织另一个诡计。他说：\par

\Quote{我承认，我在作品中加入了自己的东西：正如我在序言中所声明的，我凭自己的判断删去了不少段落；但只是那些我开始怀疑并非奥力振本人如此陈述的段落，而且在我看来，这些陈述是由他人插入的，因为在其他地方我发现作者以公教的意义阐述了同样的事情。}\par

何等出色的口才！还点缀着阿提卡风格的华丽辞藻。\Quote{此外还！}和\Quote{那些使我产生怀疑的事情！}我惊讶他竟然敢把这样的文字怪物送到罗马。人们会以为这人的舌头被锁住了，被无法解开的绳索捆绑着，以至于几乎无法迸发出人类的言语。不过，我将回到正题。\par

11 (\Emph{a}). 我想知道，是谁允许你从你正在翻译的作品中删除若干段落的？你被要求将一本希腊文著作译成拉丁文，而非修改它；转述别人的话，而非写你自己的书。你承认，由于删减了这么多，你没有做你被要求的事。而且，我希望你所删去的都是坏的部分，而你没有加入许多你自己的东西来支持坏的部分。我将举一个例子，人们可以从中判断其余。在奥力振所写的亵渎不敬的话，即子不能见父的\OriginalTerm{论首要原理}{\GreekText{Περὶ} \GreekText{᾿Αρχῶν}}第一卷中，你作为作者的代言人补充了理由，并翻译了狄迪默的注释，他徒劳地试图为别人的错误辩护，证明奥力振说得对；但我们这些可怜、单纯的人，就像恩尼乌斯所说的温顺动物一样，既不能理解他的智慧，也不能理解他译者的智慧。你在解释中所引用的序言，其中你对我如此吹捧和赞美，表明你犯了最严重的翻译错误。你说你从希腊文中删去了许多东西，但对你加入了什么却只字不提。删去的部分是好的还是坏的？我猜想是坏的。你保留的部分是好的还是坏的？我假设是好的；因为你不可能翻译坏的部分。那么，我猜想你删去了坏的部分，留下了好的部分？当然。但你的译文几乎全部可以被证明是坏的。因此，在你的译文中我能证明是坏的东西，都必须算在你的账上，因为你把它当作好的来翻译。如果你要像一位不公正的审查官那样行事，而他自己也犯了罪，可以随意将一些人逐出元老院而保留另一些人，那真是怪事。但你说：\Quote{不可能改变一切。我只想删去异端者所添加的东西。} 很好。那么，如果你删去了所有你认为异端者所添加的东西，那么你留下的所有东西都属于你所翻译的作品。那么回答我，这些是好的还是坏的？你不可能翻译坏的东西，因为你一劳永逸地删去了异端者所添加的东西，也就是说，除非你认为你的职责是删去异端者所添加的坏部分，而将奥力振本人的错误原封不动地译成拉丁文。那么告诉我，你为什么将奥力振的异端翻译成拉丁文？是为了揭露邪恶的作者，还是为了赞扬他？如果你的目的是揭露他，为什么在序言中赞扬他？如果你赞扬他，你就被证明是异端者。剩下的唯一假设是，你把这些当作好的东西出版。但如果它们被证明是坏的，那么作者和译者都卷入同样的罪行，圣咏作者的话就应验了：\BibleQuote{你见了盗贼，就与他同伙，你又与奸夫分赃。}\BibleRef{咏 50:18} 不必通过争论使一个清楚的问题变得可疑。关于以下内容，让他回答，他心中对异端者添加这些内容的怀疑从何而来。他说：\Quote{因为我发现这位作者在其他地方以公教的思想讨论了同样的内容。}\par

% structure: p0034 source=h2 target=subsection reason=relative-heading-rank; base=h2

\subsection{奥力振断言基督是受造物，并主张万有复兴。他在何处反驳了这一点？}

12. 我们必须首先考虑事实，然后按顺序得出后来的推论。现在，我在奥力振所写的许多坏东西中，发现以下最明显属于异端：天主子是受造物，圣神是服役者；有无穷的世界，在水恒的世代中相继；天使变成了人的灵魂；救主的灵魂在由玛利亚诞生之前就存在，正是这个灵魂\BibleQuote{他虽具有天主的形体，却没有以自己与天主同等为应当把持不舍的，反而空虚了自己，取了奴仆的形体}\BibleRef{斐 2:6-7}；我们身体的复活将是这样，我们不会有相同的肢体，因为当肢体的功能停止后，它们将变得多余；我们的身体本身将变得像空气和灵体，逐渐消散于稀薄的空气和虚无之中；在万有复兴之时，当赦免的圆满达到时，革鲁宾和色辣芬、宝座、宰制者、异能者、掌权者、率领者、总领天使和天使、魔鬼、恶魔以及不论基督徒、犹太人或外邦人的灵魂，都将有相同的状况和等级；当他们达到真正的形式和重量，以及整个种族从世界流亡归回的新军队呈现出理性受造物的一团，留下所有渣滓时，那么将有一个新世界以新的起源开始，以及其他的身体，使从天上坠落的灵魂穿上；因此我们可能害怕：我们现在是男人，以后可能生为女人；现在是一位贞女，那时可能成为一个妓女。这些是我在奥力振的书中指出的异端。由你来指出，你在他的哪一本书中找到了对这些观点的反驳。\par

13. 不要告诉我\Quote{你发现同一作者在其他地方以公教的思想讨论了同样的内容}，然后打发我去搜索奥力振的六千卷书——你指控最可敬的厄丕法尼主教曾读过这些书；而要准确地引用段落；这还不够，你必须逐字逐句地拿出句子来。奥力振不傻，据我所知；他不会自相矛盾。那么所有这些计算的最终结果是：你所删去的不是异端者造成的，而是奥力振本人造成的，而你翻译了他所写的坏东西，因为你认为它们是好的；书中好的和坏的东西都要算在你的账上，因为你在序言中认可了他的著作。\par

% structure: p0037 source=h2 target=subsection reason=relative-heading-rank; base=h2

\subsection{问题是，如亚纳大削对耶路撒冷的若望所说，你怀着何种动机翻译了\OriginalTerm{论首要原理}{\GreekText{Περὶ} \GreekText{᾿Αρχῶν}}？}

14. 这本辩护书的下一段如下：\par

\Quote{我既非奥力振的提倡者或辩护者，也不是第一个翻译他著作的人。我前面已有别人做了同样的事；而我是在许多人之后，应弟兄们的请求才作的。如果规定不准作这样的翻译，这种规定只对未来有效，不追溯既往；但若要在发布规定前作翻译的人受责难，那责难必须从第一个迈出这一步的人开始。}\par

他终于呕吐出他想要说的话，整个激动的情绪发泄为对我的恶意指控。当他翻译\OriginalTerm{《论原理》}{\GreekText{Περὶ} \GreekText{᾿Αρχῶν}}时，他声称是在效法我。当他因做了这事而受责难时，他以我为榜样：无论他是在危险中还是安全中，他都离不开我。因此，让我告诉他他所自称不知道的事。没有人指责你翻译了奥力振，否则希拉略和盎博罗削就会被定罪；但你翻译了一部异端著作，并试图在序言中赞美我来支持它，这才是你被谴责的理由。至于我自己，你所控告的我，翻译了奥力振的七十篇讲道以及他的部分注释，以便通过翻译他最好的作品使最坏的作品不被注意；我也公开翻译了\OriginalTerm{《论原理》}{\GreekText{Περὶ} \GreekText{᾿Αρχῶν}}，以证明你译本的谬误，好让读者知道应避免什么。如果你想把奥力振译成拉丁文，你手边有很多他的讲道和注释，其中涉及道德主题或解释圣经的隐晦段落。翻译这些吧，把这些给你所求的人。为什么你的第一份工作要从可耻的事情开始？为什么你准备翻译一部异端著作时，要用一位殉道者所谓的书来作序和支持它，并把一本译本曾引起世界恐慌的书强加给罗马人的耳朵？无论如何，如果你翻译这样的作品是为了表明作者是异端，就不要改动希腊原文的任何内容，并在序言中说明这一点。这正是教宗安纳斯塔西乌斯在写给若望主教反对你的信中最明智地包含的；他使如此行事的我免受一切指责，而谴责不肯这样做的你。你或许会否认这封信的存在；所以我附上一份副本，以便你若不听从兄弟的劝告，至少会听主教定罪的话。\par

% structure: p0041 source=h2 target=subsection reason=relative-heading-rank; base=h2

\subsection{你假装不是奥力振的辩护者，却为他出版并扩充辩护书，并声称异端者篡改了他的著作。}

15. 你说你不是奥力振的辩护者或提倡者；但我立刻用你自己的书来面对你，你那本著名的作品在其臭名昭著的序言中这样提到：\par

\Quote{这种差异的原因，我已在庞斐禄所写的辩护书中的论著里为你更充分地阐述了，并加上了一篇我自己很短的文献，在其中我用在我看来明显的证据表明，他的著作在许多地方被异端者和恶人篡改，尤其是现在正在翻译的\OriginalTerm{《论原理》}{\GreekText{Περὶ} \GreekText{᾿Αρχῶν}}。}\par

他刚离开港口就把船撞上了岩石。他引用了殉道者庞斐禄的辩护书（然而，我已证明那是亚略派首领欧瑟比的作品）中的话，他曾说过：\Quote{我尽我所能，按照主题的要求，把它译为拉丁文；}接着他又说：\Quote{正是这一点，玛加略，渴望的人啊，我愿意托付给你，好使你确信，我从他的书中上面陈述的这个信仰准则，是应当拥抱和持守的：所有书中显然都有大公含义。}尽管他从欧瑟比的书中删去了许多内容，并试图将关于圣子和圣神的表述在善义上加以修改，但其中仍有许多冒犯之处，甚至公开的亵渎，而我们的朋友既然宣称它们是大公的，就不能拒绝接受。欧瑟比（或你愿意的庞斐禄）在那本书中说：圣子是圣父的仆人，圣神与圣父圣子不是同一实体；人的灵魂从天上坠落；既然我们从天使的状态改变，在万有复兴时天使、魔鬼和人将全都平等；还有许多其他如此不敬和可恶的事，连重复它们都是罪过。奥力振的拥护者和庞斐禄的翻译者处于一种奇怪的境地：如果在他已修改的这些部分中有如此多的亵渎，那么在他假称已被异端者篡改的部分中，又会有什么亵渎神明的东西呢！他之所以持此观点，他说，是因为一个既非愚笨也非疯狂的人不可能说出相互矛盾的话；而且，为避免我们以为他在不同时期写了不同的东西，或根据写作时间发表了相反的观点，他补充道：\par

\Quote{当我们看到，有时在同一处，几乎可以说是紧接着的下一段中，出现了一句意思相反的句子，我们该说什么呢？我们能否相信，在同一部著作、同一卷书中，有时正如我所说的，就在紧接着的句子里，他会忘记自己的话？例如，那个之前说过\Quote{我们在整个圣经中找不到任何一处说圣神是被创造的或被造的}的人，立刻又加上圣神是在其他受造物中被造的？或者，那个将圣父和圣子定义为同一性体（希腊文称为\OriginalTerm{同质}{Homoousion}）的人，紧接着又说他属于另一个性体，并且他是被造的，而就在不久前他还宣称他是从天主圣父的本性中所生的？}\par

% structure: p0046 source=h2 target=subsection reason=relative-heading-rank; base=h2

\subsection{你的辩护无法从\Person{欧瑟伯}或\Person{狄迪莫}那里得到支持，他们各自出于自己的理由，为现存的\OriginalTerm{论首要原理}{\GreekText{Περὶ} \GreekText{᾿Αρχῶν}}进行辩护。}

16. 这是他自己的话，他无法否认。现在，我不希望被诸如\Quote{他上文说过}这样的说法搪塞，而是想知道那本书的名字，在那本书中他先说话正确，后来说话错误：在那本书中，他先说圣神和圣子出于天主的性体，紧接着又宣称他们是受造物。难道你不知道我拥有\Person{奥利金}的全部著作，并且读了其中很大一部分吗？\par

你向民众展示的装饰！我深知你；\newline{}我看见里面和皮上的一切。\par

\Person{欧瑟伯}是一个非常博学的人，（注意，我说他有学问，但不是公教信徒；你不可以像往常一样以此作为诽谤我的借口。）他用了整整六卷书，只试图表明\Person{奥利金}和他信仰一致，也就是亚流派之背信。他举出了许多经文证据，并确实证明了他的论点。你是在亚历山大城的监狱中做了什么梦，得到启示，使你认定这些被他接受为真的段落是伪造的？但也许，作为亚流派，他接受这些异端添加的内容以支持自己的错误，这样他就不会被认为是唯一持有违背教会的错误观点的人。那么，关于\Person{狄迪莫}，你又如何回答？他无疑在圣三论上是公教的。你知道我将他的\WorkTitle{论圣神}一书译成了拉丁文。他肯定不可能赞同那些被异端者添加到奥利金作品中的段落；然而，他为你翻译的\OriginalTerm{论首要原理}{\GreekText{Περὶ} \GreekText{᾿Αρχῶν}}写了简短的注释；在这些注释中，他从未否认那里所写的是出于奥利金，而只是试图说服我们这些愚钝的人，我们没有理解他的意思，以及这些段落应当从好的意义上理解。以上是关于圣子和圣神单独说的。但关于奥利金的其他教义，欧瑟伯和狄迪莫都遵循他的观点，并以公教和基督徒的意义来辩护那些所有教会都谴责的教义。\par

% structure: p0050 source=h2 target=subsection reason=relative-heading-rank; base=h2

\subsection{如果我们动辄断定伪造，就会使所有过往文献陷入混乱。}

17. 但让我们考察一下他试图证明\Person{奥利金}的著作被异端者篡改时所用的论据。\par

他说：\Quote{克肋孟是宗徒的门徒，在主教职和殉道上继承宗徒，他写了名为\WorkTitle{认主}（即\WorkTitle{识别}）的书。在这些书中，尽管总体上以宗徒伯多禄名义提出的教义是真实的，但在某些段落中，欧诺弥的教义被引入，以至于你会以为欧诺弥本人在进行辩论并主张圣子是从无中受造的观点。}\par

接着，在跳过一段太长的段落之后，他补充说：\par

\Quote{那么我们对这些事实该作何感想？我们必须相信一位宗徒传人写了异端吗？还是更可能是一些心术不正的人，希望为自己的教义获得支持并更易取得信任，将圣人不可能持有或写下的观点托于圣人之名？}\par

他告诉我们，亚历山大城的司铎\Person{克肋孟}也是一位公教信徒，有时在他的著作中写道天主子是被造的；还有亚历山大城主教\Person{狄约尼削}，一位极其博学的人，在他反驳撒伯里乌教义的四卷书中，陷入了亚流的教条。他引用这些例子的目的，不是要表明教会人士和公教信徒犯了错误，而是他们的著作被异端者篡改了，他以下面这些话结束讨论：\par

\Quote{当我们发现奥利金在教义上有某种不一致，正如我们在上述诸位身上发现的那样，难道我们不应该以同样的态度相信，在他身上的情形与我们回顾过的公教人士身上的情形相同吗？当情形相同时，同样的辩护不是也成立吗？}\par

我回答说，如果我们承认，书中一切令人反感的内容都是别人篡改插入的，那么就不会有任何东西属于该书署名作者，而一切都可以归于那些被认为篡改的人。但那样的话，这些内容也不会属于他们，因为我们不知道他们是谁：结果就是每本书都属于所有人，而没有任何东西属于任何一个特定的人。在这种辩护方法引入的混乱中，我们将无法定\Person{玛尔西翁}、\Person{摩尼}、\Person{亚流}或\Person{欧诺弥}的错误；因为，一旦我们指出他们不信的言论，他们的门徒就会回答说，那不是老师写的，而是被对手篡改插入的。按照这个原则，你这本书本身将既不是你的，也不是我的。至于我反驳你指控的这本书，你在其中发现的任何错误，都将被认为不是我写的，而是现在挑剔它的你写的。此外，当你把一切归于异端者时，就没有任何东西可以归于教会人士作为他们自己的了。\par

但你会问：既然如此，他们的著作中为何会出现一些错误的观点？好吧，如果我回答说我并不知道这些错误观点来自何方，那绝不能认为我是说他们是异端者。他们有可能是无意中陷入错误，或者那些话有别的含义，或者它们逐渐被不熟练的抄写者所篡改。必须承认，在亚略作为南方的恶魔出现在亚历山大里亚之前，有些话是不经意间说出的，无法抵御恶意的批评。但是，当奥力振的明显错误被揭露时，你并不为他辩护，反而指控别人；你不否认错误，却召集一群罪犯。如果有人要你指出谁是奥力振异端中的同伴，那么把这些其他人叫来倒是合适的。但现在要求你告诉我们的是：奥力振著作中的那些陈述是好是坏？你什么也不说，反而提出不相干的事情，诸如：这是克莱孟说的；这是狄约尼削被认定犯下的错误；这是亚太纳削主教为狄约尼削的错误所作的辩护；宗徒的著作也曾以类似方式被篡改。然后，当异端的指控加在你身上时，你却不为自己辩护，反而承认关于我的事。我不指控任何人，只满足于为自己辩护。我不是你试图证明的那样：你是不是像你所被指控的那样，这由你自己考虑。我被免于责备这一事实并不证明我无罪，也不证明你被指控就表明你有罪。\par

% structure: p0059 source=h2 target=subsection reason=relative-heading-rank; base=h2

\subsection{奥力振那封他只翻译一部分的信，其目的不是要显示他的著作被篡改，而是要谴责那些谴责他的主教们。}

18. 在关于异端者篡改宗徒、两位克莱孟和狄约尼削著作的这一开场白之后，他终于谈到了奥力振，他的话是这样的：\par

\Quote{我已从他自己的话和著作中表明他本人如何抱怨和痛惜此事：他在写给亚历山大里亚一些密友的信中清楚地解释了他住在肉身中、感官完全清醒时所遭受的，即他的书和论文被篡改，或者出现了伪作。}\par

他附加了这封信的副本；而这位向异端者恳求奥力振著作被篡改的人，自己却以篡改开始，因为他没有按照他在希腊文中找到的原文翻译这封信，也没有将奥力振在信中所陈述的内容传达给拉丁人。整封信的目的是攻击亚历山大里亚的教宗德默特琉，并谴责世界各地的主教们，告诉他们，他们对他的绝罚是无效的；他还进一步说，他无意报复他们的恶言；事实上，他如此害怕恶言，以至于连对魔鬼也不敢说恶言；以至于他给瓦伦廷谬误的追随者康迪杜斯提供了机会，谎称他说魔鬼的本性是可以得救的。但我们的朋友没有注意信的真实主旨，而为奥力振编造了一个他没有使用的论点。因此，我从信的一部分翻译了一小段，从已经提到的地方稍往下开始，并把它附在他用删节和不诚实方式翻译的部分之后，以便读者可以看出他隐瞒前一部分的目的。于是，他普遍地反对教会的主教们，因为他们判定他不配与其共融；他接着说：\par

\Quote{何必提先知们不断威胁和责备牧者、长老、司祭和首领的话呢？这些事你们自己无需我的帮助就可以从圣经中引出，并且可以清楚地看到，现在可能正是那所说的\BibleQuote{不要信赖朋友，也不要指望首领}的时候，而且那预言如今正在应验：\BibleQuote{我人民的领袖不认识我；我的儿子是愚昧的，不聪明；他们善于作恶，却不知行善。}我们应当怜悯他们，而不是恨他们；应当为他们祈祷，而不是诅咒他们。因为我们被造是为了祝福，而不是为了诅咒。因此，当弥额尔就梅瑟的身体与魔鬼争辩时，连这么严重的恶，他也不敢用毁谤的罪状，只说：\BibleQuote{愿主责斥你。}我们在匝加利亚也读到类似的话：\BibleQuote{愿主责斥你，撒殚！愿拣选耶路撒冷的主责斥你。}同样，我们也希望那些不愿谦卑接受邻人责备的人，能受到主的责备。但是，既然弥额尔说：\BibleQuote{愿主责斥你，撒殚！}匝加利亚也这样说，魔鬼很清楚主是否责斥他，也必定承认责斥的方式。}\par

然后，在太长得无法插入此处的一段之后，他补充说：\par

\Quote{我们相信，不仅犯了大罪的人将被逐出天国，例如行淫、通奸、与男人同寝的、偷窃的，而且那些犯下不太明显恶行的人也是如此，因为经上记载：\BibleQuote{醉酒的和毁谤人的，不能继承天主的国。}并且，人受审判的标准既是天主的仁慈，也是天主的严厉。因此，我们努力在一切事上谨慎行事，在饮酒和节制言语上，使我们不敢毁谤任何人。现在，由于对天主的敬畏，我们小心不对任何人发出咒骂，记着\BibleQuote{他不敢用毁谤的罪状控告他}这句话是针对弥额尔与魔鬼打交道时说的，正如别处所说：\BibleQuote{他们藐视上主，亵渎众尊荣者。}某些寻找争论理由的人，便把这种亵渎归给我们和我们的教训（他们必须把适用于他们自己的话放在心上：\BibleQuote{醉酒的和毁谤人的，不能继承天主的国？}），即认为那些将被逐出天国的人的邪恶和灭亡之父可以得救；这种事情连疯子也不至于说。}\par

同一封信中其餘部分，他取代了我所翻譯的奧力振後面的話：\Quote{現在，因為對天主的敬畏，我們謹慎不咒罵任何人}等等；他欺詐地截斷了前段，而後段依賴於前段，並開始翻譯那封信，好像前段以此陳述開始，並說：\par

\Quote{有些喜歡控告鄰人的人，將我們從未說過的褻瀆罪歸於我們和我們的教導（關於此，他們必須考慮是否願意遵守那說\InnerQuote{惡言者不得繼承天主的国}的誡命），因為他們說我斷言那將被逐出天主的国的邪惡與滅亡之父，即魔鬼，將得救；這事即使一個人失去理智並明顯瘋狂也不可能說出。}\par

% structure: p0068 source=h2 target=subsection reason=relative-heading-rank; base=h2

\subsection{惟有在與坎迪都斯辯論中的某一點，奧力振才指控這種偽造。希拉里因著作被偽造而受譴責的故事並無根據。}

19. 現在請比較我逐字翻譯的奧力振原文，與被他轉譯成拉丁文、或者說顛覆的這些文字；你將清楚看到它們之間有多大的差異，不僅是詞語，更是意義。我請求你別因我的譯文較長而厭煩；因為我翻譯整段的目的，是要揭示他壓制前段的意圖。希臘文中存有奧力振與坎迪都斯（瓦倫提尼安異端的辯護者）之間的對話錄，我承認當我閱讀時，彷彿在看兩位安達巴提亞角鬥士的對決。坎迪都斯主張聖子是出於聖父的實體，陷入斷言\OriginalTerm{流溢}{Probolé or Production}或產出的錯誤。另一方面，奧力振像亞略和歐諾米一樣，拒絕承認祂是被產生或出生的，以免天主聖父因而被分割；但他說聖父按自己的意願，像其他受造物一樣產生了一個崇高且至卓越的受造物。然後他們進入第二個問題。坎迪都斯斷言魔鬼的本質全然邪惡，永不能得救。對此，奧力振正確地主張魔鬼不是可朽壞的實體，而是因自己的意志墮落，並且可以得救。坎迪都斯卻虛假地將其轉為對奧力振的指責，好像他說過魔鬼的本質可以得救。因此，奧力振駁斥坎迪都斯虛假指控他的內容。但我們看到，惟獨在此對話錄中，奧力振指控異端者偽造了他的著作，而不是在其他從未被質疑的書卷中。否則，如果我們相信所有異端內容並非出於奧力振，而是出於異端者，而他的幾乎所有著作都充滿了這些錯誤，那麼奧力振的作品將一無所剩，一切都必是那些我們不知其名的人的作品。\par

他誹謗希臘人和古人還不夠，對於他們，時間或空間的距離讓他能夠隨意編造謊言。他來到拉丁人中間，首先舉起宣信者希拉里的案例，聲稱他的書在阿里米努姆大公會議後被異端者偽造。為此，在一次主教會議中，關於他的問題被提出，他隨後命人從自己家中取來那書。那本具有異端形式的書就在他的書桌裡，他卻不知情；當被出示時，該書的作者被定為異端並被絕罰，離開了會議廳。這個故事，只是他對親密者講述的自己的夢幻；他以為自己的權威如此之大，以至於當他說這些話時，沒有人敢反駁他。我要問他幾個問題。希拉里被絕罰的會議是在哪個城市舉行的？在場的主教們叫什麼名字？誰簽署了判決？誰同意，誰不同意？那一年的執政官是誰？哪位皇帝下令召集會議？在場的主教只是高盧的，還是也包括意大利和西班牙的？會議是為了什麼目的召集的？你沒有告訴我們這些；然而，為了辯護奧力振，你卻將一位最具雄辯力、拉丁語反對亞略派的號角般的人物視為罪犯和被絕罰者。但我們面對的是一位宣信者，甚至他的誹謗也必須耐心忍受。他接著轉向著名的殉道者西彼廉，告訴我們，在君士坦丁堡，馬其頓異端的支持者將特土良的一本名為\WorkTitle{論天主聖三}的書當作西彼廉的著作來讀。在這項指控中，他說了兩個謊言。那本書不是特土良的，也不是以西彼廉之名流傳的；它是諾瓦提安的，並以其名字稱呼；文風的獨特性證明了該著作的作者。\par

% structure: p0071 source=h2 target=subsection reason=relative-heading-rank; base=h2

\subsection{你在達瑪甦的議會中關於我所講的，不過是飯後的閒談。}

20. 这是何等胡言乱语，他们竟从中捏造对我的指控！似乎不值一提。这是我关于罗马主教达玛苏斯所主持的会议的讲述，而我自己被以他某位朋友的名义受到攻击。他曾给我一些关于教会事务的文件去抄写；记述中提到亚玻里纳留派耍的一个花招：他们借走了其中一份文件，即亚大纳削的一部著作，其中出现了\Quote{\OriginalTerm{主的人}{Dominicus homo}}一词，他们先是将其刮去，然后在擦痕处重新写上，这样看起来就不是他们篡改了那书，而是我添加了那词。我恳求你，我至亲爱的朋友，在这些与教会密切相关、涉及教义真理、且我们为灵魂的益处寻求先贤权威的事务中，丢弃这类无聊的东西，不要把饭后的谈资当作论据。因为很可能，即便你从我这里听到了真实的故事，另一个不知情的人仍可能宣称它是编造的，并且是你像菲利斯提翁的寓言或兰图卢斯、马尔凯卢斯的歌谣一样用优雅的语言写成的。\par

% structure: p0073 source=h2 target=subsection reason=relative-heading-rank; base=h2

\subsection{攻击厄丕法尼作为奥力振的剽窃者是对主教们的普遍侮辱。奥力振从未写过六千卷书。}

21. 一旦放纵的缰绳松开，鲁莽会达到何种地步？在讲述了希拉里的绝罚、那部假借西彼廉之名的异端著作、以及我在睡着时对亚大纳削著作的连续刮擦和插入之后，他最后爆发了对教宗厄丕法尼的攻击：他因厄丕法尼在写给若望主教的信中称他为异端者而在心中产生的懊恼，在他为奥力振的辩护中倾泻而出，并用以下话语安慰自己：\par

\Quote{隐藏的真相必须在此揭露。任何人都不可能行使如此不义的判断，以致在案件相同时作出不平等的裁决。但事实是，那些诽谤奥力振的人的唆使者，要么是习惯在教会中长篇大论的人，要么是写作书籍的人，而这些书籍和讲论全都取自奥力振。因此，为了防止人们发现他们的抄袭——只要他们对恩师忘恩负义，这一罪行就能隐藏——他们阻止所有单纯的人阅读他。其中一人，自认为负有在万民万语中谈论奥力振的恶行的必要，仿佛那就是宣讲福音，曾在一大群弟兄面前宣称他读过奥力振的六千卷书。如果他像他惯常宣称的那样，是为了知道奥力振有什么害处而读，十卷或二十卷，至多三十卷就足以知道。读六千卷书不像一个想知道他里面的害处和错误的人，而像一个将几乎全部生命都奉献给他所指导的研究的人。那么，他怎敢要求别人听从他指责那些为求教导而读了他一小部分著作、同时注意保持自己整个信仰体系和虔诚的人呢？}\par

22. 这些惯常在教会中进行如此长篇的辩论、写作书籍、并且其讲论和著作全部取自奥力振的人是谁？这些担心自己文字盗窃败露、对恩师忘恩负义、从而阻止单纯之人阅读他的人？你应当指名道姓，指出这些人。难道可敬的主教阿纳斯塔西、德敖斐罗、味内留、克罗马提，以及东西方整个公教会的主教会议——他们都公开指认他为异端者——要被看作他著作的剽窃者吗？难道我们要相信，当他们在教会中宣讲时，他们不是宣讲圣经的奥秘，而只是重复他们从奥力振那里偷来的东西？你笼统地贬低他们所有人还不够，还非要用你的笔之矛特别瞄准教会中一位可敬而杰出的主教？自认为有责任在各民族和万语中辱骂奥力振、如同必须宣讲的福音一样的人是谁？这个人曾在一大群弟兄面前宣称他读过奥力振的六千卷书。在那群弟兄中，你自己就在正中央，当时，正如他在信中所抱怨的，你正在大谈奥力振的怪诞学说。难道懂得希腊语、叙利亚语、希伯来语、埃及语、以及部分拉丁语要算作他的罪过吗？那么，我想，那些说各种方言的宗徒和宗徒时代的人也要被定罪；而你懂得两种语言，可以嘲笑我懂得三种。至于你假装他读过的六千卷书，谁会相信你说的是实话，或者他能够撒这么大的谎？如果奥力振真的写了六千卷书，那么一个学识渊博、自幼受神圣文学训练的人，有可能因为好奇心和对学问的热爱而阅读与自己信念相异的书。但他怎么可能读奥力振从未写过的书？数一数欧瑟比乌斯第三卷中收录的索引，即他的\WorkTitle{邦斐禄生平}，你会发现，我不说六千卷，连那个数字的三分之一都不到。我手头有上述教宗的信函，他在信中回答了你仍在东方时发出的这一诽谤；它以真理的坦诚面容驳斥了这个极其明显的谎言。\par

% structure: p0077 source=h2 target=subsection reason=relative-heading-rank; base=h2

\subsection{我在凯撒勒雅的图书馆确信，你当作邦斐禄作品引用的\WorkTitle{辩护词}是欧瑟比乌斯的著作。}

23. 在此之后，您竟敢在您的\WorkTitle{护教篇}中说，您并非\Person{奥利振}的辩护人或捍卫者，尽管您认为\Person{欧瑟伯}和\Person{庞斐禄}为他的辩护说得太少。我将在另一篇论文中尝试回复那些作品，若天主赐我足够的寿命。目前，只需说明我已回应了您的断言，并通过声明我从未见过被称为\Person{庞斐禄}著作的成书，直至我在您的手稿中读到它，以此提醒细心的读者提防。知晓为异端者所写的内容并非我所关心，因此我一直以为\Person{庞斐禄}的著作与\Person{欧瑟伯}的不同；但在问题提出后，我愿回复他们的作品，为此我阅读了两人为\Person{奥利振}辩护的言论；于是我清楚辨识出\Person{欧瑟伯}六卷书中的第一卷，正是您以希腊文和拉丁文出版的单卷\Person{庞斐禄}著作，仅修改了关于圣子和圣神的意见，而这些意见明显带有公开亵渎的标记。因此，当我的朋友\Person{德克斯特}，担任禁卫军长官之职，十年前请我为他编列我们信仰的作者名单时，我将这本归于\Person{庞斐禄}名下的书列入各作者的不同论文中，以为事实正如您和您的门徒所公开的那样。但是，既然\Person{欧瑟伯}本人说\Person{庞斐禄}除了一些写给朋友的短信外什么也没写，而他的六卷书的第一卷包含了您假借\Person{庞斐禄}之名虚构的话，那么显然，您传播这本书的目的是以殉道者的权威引入异端。我不能允许您用我的错误来掩盖您的欺骗，您先是假装这本书是\Person{庞斐禄}所著，然后又歪曲其中的许多段落，使拉丁文与希腊文不同。我相信这本书是由其署名作者所著，正如我相信\OriginalTerm{论首要原理}{\GreekText{Περὶ} \GreekText{᾿Αρχῶν}}及其他许多\Person{奥利振}和其他希腊作家的作品一样，这些我以前从未读过，现在被迫阅读，因为异端问题已被提出，我希望知道什么应当避免、什么应当反对。因此，我年轻时只翻译了他公开讲授的讲道，其中较少引起反感；并且是在无知和他人请求下进行的：我没有试图利用人们赞同的部分来诱使人们接受那些明显是异端的部分。无论如何，为长话短说，我可以指出我从何处得到\OriginalTerm{论首要原理}{\GreekText{Περὶ} \GreekText{᾿Αρχῶν}}，即从那些从您的手稿抄写的人那里。我们同样想知道您的抄本从何而来；因为如果您不能指出任何其他来源，您自己将被证明伪造了它。\Quote{善人从他心中的善库中取出善来。}好树由它果实的甘甜可知。\par

% structure: p0079 source=h2 target=subsection reason=relative-heading-rank; base=h2

\subsection{在非洲被假托为我所写、并对我从希伯来文翻译\WorkTitle{旧约}表示遗憾的那封信，带有您的手笔。我一直尊崇\OriginalTerm{七十贤士}{Seventy Translators}。}

24. 我的弟兄\Person{欧瑟伯}写信给我说，当他参加一次为某些教会事务召集的非洲主教会议时，发现那里有一封据称是我写的信，信中我表示忏悔，承认是因年轻时的出版压力而将圣经从希伯来文译成拉丁文；这些全无一句真话。我听到这事，惊愕不已。但一个证人不够；连\Person{卡托}也不能仅凭他的孤证而取信：\BibleQuote{凭两三个证人，一切话都可成立。} 不久许多罗马的弟兄给我来信询问此事，问事实是否如所述；他们指出了一个让我哭泣的人，就是那在百姓中散布这封信的人。他敢于做这事，还有什么不敢的？幸而恶意没有与意图相等的力量。如果邪恶总是与权力结合，毁谤总能得逞其所求，无辜早就灭亡了。他尽管多才多艺，也无法模仿任何一种写作风格和方式，无论其价值如何；在他所有的诡计和欺诈地冒充另一个人格中，他是谁显而易见。那么，就是这个人，以我的名义写了这封虚构的撤回信，声称我将希伯来文书籍译成拉丁文是坏的，我们现在听说，他指控我翻译圣经是为了贬低\OriginalTerm{七十贤士译本}{Septuagint}。无论如何，我的翻译是对是错，我都要受谴责：我必须要么承认我的新作是错的，要么承认我的新译本旨在打击旧译本。我奇怪在这封信中他没有给我加上杀人的罪名，或奸淫、或亵渎、或弑亲，或任何心灵沉默运作所能酝酿的卑鄙之事。的确，我应该感谢他，从整个罪行林中只给我加上一次错误或欺诈的作为。我是否可能说过任何贬低\OriginalTerm{七十贤士译本}{Septuagint}的话？多年前我曾仔细清除了该译本的错讹，将其交给拉丁读者，并在我们的修院聚会中每日讲解；该译本的圣咏集曾是我长期默想和歌唱的对象。我竟如此愚蠢，想在老年忘记青年时所学的吗？我所有的论著都是根据他们译本所证实的言论织成的。我对十二先知书的注释是对他们译本和我自己译本的解说。人的劳作是多么无常！研究的结果与其初衷有时多么相反。我以为通过这译本惠及拉丁同胞，给了他们学习的激励；因为它现在被重译成希腊文后，也不被希腊人轻视；然而现在它成了控诉我的理由；我发现给他们的食物反胃了。如果无辜成为控告的对象，人生还有什么安全可言？我是家主，发现仇敌趁我睡觉时在麦子中撒了莠子。\BibleQuote{森林中的野猪把我葡萄园掘了，田野中的野兽把它吞吃了。} 我保持沉默，但一封不是我写的信却在控告我。我不知道加给我的罪行，却被迫在全地承认罪行。\BibleRef{耶 15:10}\BibleQuote{祸哉，我的母亲！你生了我这个人，使我在全地受审判和定罪。}\par

% structure: p0081 source=h2 target=subsection reason=relative-heading-rank; base=h2

\subsection{作为证明，我提出我所译\BibleBook{}{创世纪}至\BibleBook{}{依撒意亚}各卷的序言。}

25. 我对旧约各卷的所有序言——以下我附上一些例证——在这点上为我作证；无需以与其中所述不同的方式陈述此事。因此，我将从\BibleBook{}{创世纪}开始。序言如下：\par

我从亲爱的\Person{德西德里乌}得了长久以来热切盼望的信，他仿佛预知了未来，与\Person{达尼尔}同名，恳求我将\OriginalTerm{梅瑟五书}{Pentateuch}从希伯来文译成拉丁文，以交给我们的朋友。这工作确属冒险，暴露于诽谤者的攻击之下，他们断言我是出于对七十贤士译本的鄙视，才着手造一个新译本以取代旧译本。他们这样测试能力，如同测试酒；而我已再三声明，我忠心地在天主的会幕中献上我所能的，并指出一人带来的大恩赐并不因另一人较劣的恩赐而减损。但我是受\Person{奥力振}的热忱所激励而承担此任务，他将\Person{特敖多提翁}的译本与旧版本混合，并在全书中使用星标*和剑标作为区分标记。\par

\SourceRecord{Translated by W.H. Fremantle.；Nicene and Post-Nicene Fathers, Second Series；Vol. 3.；Edited by Philip Schaff and Henry Wace.；Buffalo, NY: Christian Literature Publishing Co.,；1892.；<http://www.newadvent.org/fathers/27102.htm>.}

% structure: p0001 source=h1 target=section reason=page-title

\section{驳\Person{鲁菲努斯}辩护书（卷三）}

\EditorialNote{前两卷构成了一个完整整体，但有暗示，当\Person{热罗尼莫}收到\Person{鲁菲努斯}的全部作品后，可能还有后续。前两卷由一艘与\Place{阿奎莱亚}通商的商船船长带给\Person{鲁菲努斯}，并附有一份\Person{热罗尼莫}曾被\Person{帕玛丘}扣押的友好信件的副本。信使（按\Person{鲁菲努斯}所说，\Person{热罗尼莫}嘲笑这不可能）只等了不到两天。\Place{阿奎莱亚}的\Person{克罗玛蒂乌斯}主教敦促斗争应就此结束，并成功使\Person{鲁菲努斯}未作公开答复。然而，他私下写了一封信给\Person{热罗尼莫}，这封信未能传世，而且从第4章、第6章等处所引的片段来看，似乎并非和平的基调。其详情可从\Person{热罗尼莫}的答复中窥知。\Person{热罗尼莫}暗示，此信试图将他卷入异端，重申并加剧了先前的指控，以只适用于最底层舞台人物的语言谈论他；并且宣称，如果作者有此意，他可以提出一些事实，这些事实将使他的对手毁灭。\Person{热罗尼莫}虽然收到了\Person{克罗玛蒂乌斯}希望他不作答复的表示（可能还有\Person{奥古斯丁}的遗憾劝诫——\Person{热罗尼莫}书信第110章第6节，\Person{奥古斯丁}书信第73章），但他声明自己不可能让步。他无法不为死罪指控辩护，也不能宽恕异端者。唯有在信仰中的合一才能带来和平。}\par

% structure: p0003 source=h2 target=subsection reason=relative-heading-rank; base=h2

\subsection{你的信充满谎言与暴力。我将努力不采用同样的口气。}

1. 我读了您凭智慧写给我的信。您抨击我，尽管您曾称赞我，称我为真正的伙伴和兄弟，而现在您却写书来召唤我回应那些恐吓我的指控。我看出在您身上应验了\Person{撒罗满}的话：\BibleQuote{愚昧人口中有傲慢的杖}，以及\BibleQuote{愚昧人不接受明智的话语，除非你说出他心中的事}；也如\Person{依撒意亚}所说：\BibleQuote{愚昧人讲论愚昧，他的心怀思念虚假之事，以行邪恶并说谎言反对主。}因为您何必将满篇指控和诅咒的整卷书送给我，并公之于众，而在信的末尾您却以死威胁我，如果我胆敢回应您的诽谤——请原谅——回应您的赞美？因为您的赞美和您的指控是一回事；同一个泉源涌出甜水和苦水。我请求您以您推荐给我的谦逊和羞耻之心为榜样；您指责他人说谎：请停止您自己的谎言。我不愿给人绊脚石，也不会成为您的指控者；因为我不仅要考虑您应得的，也要考虑什么是相称于我自己的。我战栗于我们救主的话：\BibleQuote{谁若使这些信我的小子中的一个跌倒，倒不如拿一块驴拉的磨石，系在他的颈上，沉在海的深处更好}；以及\BibleQuote{世界因了恶表是有祸的，恶表固然免不了要来，但立恶表的那人是有祸的。}我本也可以堆积谎言反对您，说我曾听到或看到无人察觉的事，以便在无知者中间我的厚颜被称为诚实，我的暴力被称为决心。但愿我决不效法您，不做我自己所谴责的事。能做肮脏事的人才会用肮脏话。\BibleQuote{恶人从邪恶的宝库中发出邪恶来，因为心里充满什么，口里就说什么。}您可以视为幸运，因为您曾称为朋友而现在指控的那个人无意对您做卑鄙的指控。我说这话并非因为害怕您指控的刀剑，而是因为我宁愿被指控也不愿指控别人，宁愿受伤害也不愿伤害别人。我知道宗徒的训诫：\BibleQuote{亲爱的，你们不可为自己复仇，宁可给天主的义怒留有余地，因为经上记载：复仇是我的事，我必要报应，主说。所以，你的仇人若饿了，就给他吃；若渴了，就给他喝；因为你这样作，就是将炭火堆在他的头上。}因为为自己复仇的人不能声请主的辩护。\par

% structure: p0005 source=h2 target=subsection reason=relative-heading-rank; base=h2

\subsection{我们为何不能存异而仍为朋友？为何您以死亡的威胁强迫我答复？}

2. 然而，在我答复您的信之前，我必须与您理论；您作为隐修士中长者，善司铎，基督的追随者；您竟可能愿意杀死您的兄弟，而即使恨他已是杀人？您从救主学到的教导难道不是：若有人打你一边的脸，你应把另一边也转给他？祂不是对打祂的人回答说：\BibleQuote{我若说得不对，你指证那里不对；若说得对，你为什么打我？}您以死亡威胁我，而死亡是连蛇也能加给我们的。死亡是众人的命运，只有软弱者才杀人。那么怎样？如果您不杀我，我就永不死吗？也许我该感激您将这必然性变成了德行。我们读到宗徒们争吵，即是\Person{保禄}和\Person{巴尔纳伯}，他们因\Person{若望}（其别号\Person{马尔谷}）而彼此生气；那些被基督福音的纽带联合起来的人，竟因航行而分离；但他们仍是朋友。同一位\Person{保禄}岂不是当面反对\Person{伯多禄}，因为他没有按福音的真理而行？然而\Person{保禄}称他为他在福音中的前辈，教会的柱石；并向他陈述他的宣讲方式，\BibleQuote{免得我白白奔跑，或曾奔跑。}难道在宗教事务上，儿女与父母、妻子与丈夫没有分歧吗？而家庭情感仍然无损。如果您和我一样，为何恨我？即使您相信不同，为何要杀我？难道与您不同者都当杀？我呼求\Person{耶稣}——祂将审判我现在所写的和您的信——为我的良心作证，当可敬的\Person{克罗玛蒂乌斯}主教恳请我保持沉默时，我本想沉默，从而结束我们的争端，以善胜恶。但现在您以消灭我相威胁，我不得不答复；否则，我的沉默将被视为承认罪行，您会将我的节制解释为良心有愧的迹象。\par

% structure: p0007 source=h2 target=subsection reason=relative-heading-rank; base=h2

\subsection{您可耻的嘲弄，说我想通过贿赂您的秘书获得您\WorkTitle{辩护书}的副本，这是将您自己的行为归咎于我。}

3. 我所处的两难之境是您造成的：它并非从您所不知的辩证法资源中产生，而是出自凶手的工具及其意图。如果我保持沉默，我便有罪；如果我开口，我便成了恶言者。您同时禁止我答复又强迫我。那么我必须避免两方面的极端。我不会说任何伤人的话；但我必须消解对我的指控，因为对准备杀您的人不可能不害怕。我将按您现在摆在我面前的内容顺序来做，其余内容保留在您那些最有学识的书中，我在阅读它们之前已经驳斥了它们。\par

你说：\Quote{你把对我的指控不是寄给许多人，而是只寄给那些被我所说的话冒犯的人；因为向基督徒说话，不应是为炫耀，而是为建树。}那么，我请你考虑，你是如何得知我已写了这些书的消息？是谁把它们在\Place{罗马}、\Place{意大利}以及\Place{达尔马提亚}沿岸的岛屿上广泛散播？如果这些指控只藏在你的书桌和你朋友们的书桌里，它们又怎会传到我的耳中？你怎么敢说你是在按基督徒的方式说话，不为炫耀而为建树，却在自己成熟的年纪，针对与你同辈的人说出那些话，就连凶手对窃贼、娼妓对同行、小丑对闹剧演员都难以出口的话？你长久以来一直辛劳地制造出这些山一样的指控来攻击我，磨利这些剑来刺穿我的喉咙。你的喊声如同\Person{刻瑞斯}的抱怨或车夫对马的吆喝一样响亮。这是为了让这些喊声回荡过的所有省份都读到你写给我的赞美吗？并且在每条街道、每个角落，甚至妇女的织布作坊里背诵你对我的颂词？这就是你所说的宗教节制和基督徒的建树。你的缄默、你的沉默如此这般，使得人们从西方成群结队地来到我这里，告诉我你对我的辱骂；尽管只是凭记忆，却如此一致，以至于我不得不作出回应，不是回应我那时尚未读过的你的著作，而是回应据称其中所含的内容，并用真理之盾阻挡那正在全世界飞行的谎言之箭。\par

4. 你的来信接着说：\par

\Quote{我恳请你不要费心拿出大笔金币贿赂我的秘书，就像你的朋友们在我那些包含\OriginalTerm{论原理}{\GreekText{Περὶ} \GreekText{᾿Αρχῶν}}的文稿尚未修改和完成时所做的那样，以便他们更容易篡改那些无人拥有、或至少极少人拥有的文件。请接受我免费送给你的这份文稿，尽管你本会愿意付一大笔钱买下它。}\par

我本以为你会为你作品的这种开端感到羞耻。什么！我贿赂你的秘书！有谁会试图与\Person{克里萨斯}和\Person{大流士}的财富比肩？有谁突然面对\Person{德马拉托斯}或\Person{克拉苏}时不战栗？你变得如此厚颜无耻，以致你信任谎言，以为谎言会保护你，并且我们会相信你编造的任何虚构？那么，是谁从我们的兄弟\Person{欧瑟伯}的小室里偷走了那封高度赞扬你的信？是谁的诡计和同伙，使一份文件在基督徒妇女\Person{法比奥拉}和智者\Person{欧加诺}的住所被发现，而他们自己从未见过？你是否因为能把所有本属于你的指控转嫁给别人，就以为自己无罪？每一个冒犯你的人，无论他多么清白无害，是否立即被视为罪犯？我想你之所以这样认为，是因为你拥有那种使\Person{达那厄}的贞洁被摧毁的东西，那东西比他的主人的神圣品格对\Person{革哈齐}更有力量，\Person{犹达斯}也因它出卖了他的主。\par

% structure: p0013 source=h2 target=subsection reason=relative-heading-rank; base=h2

\subsection{\Person{欧瑟伯}不该指控你；但你对他指控不能成立。}

5. 让我们了解一下我朋友所做的不当之事，你说他\Quote{篡改了你那些尚未修改和完善的文稿的部分内容，而且由于几乎无人拥有这些文稿，所以更容易篡改}。我已经说过，现在在天主面前再次重申，我不赞成他指控你，也不赞成任何基督徒指控另一位基督徒；因为那些可以在私下纠正或改正的事情，何必要公开宣扬，使多人绊倒和跌倒呢？然而既然各人为自己的口腹而活，一个人不会因为成为你的朋友就成为你意志的主人；尽管我责备指控弟兄的行为，即使指控属实，但我同样不能接受对一位圣德之人关于篡改你文稿的指控。一个只懂拉丁文的人，如何能在希腊文的翻译中改变任何东西？或如何在像\OriginalTerm{论原理}{\GreekText{Περὶ} \GreekText{᾿Αρχῶν}}这样的书中增删任何内容？其中一切都如此紧密相连，一部分依赖另一部分，任何按你意愿加入或删除的内容都会立刻像衣服上的补丁一样显眼。你要求我做的事，你自己应当去做。至少在你的断言中，即使没有基督徒的谦逊，也要有一点点天然的谦逊；不要轻视并践踏自己的良心，不要以为凭言辞的表象就能证明自己有理，而事实却对你不利。如果\Person{欧瑟伯}花钱买了你未经修改的文稿来篡改，那就拿出未被篡改的原稿来；如果你能证明其中没有任何异端内容，那么他就将承担伪造指控的责任。但是，无论你如何修改或订正，你都无法使它们变成大公的。如果错误只在于措辞或某些陈述，就可以砍掉坏的，代之以好的。但若整个讨论基于一个单一原则，即所有理性受造物都因自己的意志堕落，并且将来会回到合一状态，然后从那起点再次开始另一个堕落：那你还有什么可以修正的，除非你改变整本书？但如果你考虑这样做，你就不再是翻译别人的作品，而是创作自己的作品了。\par

然而，我几乎看不出你的论证意在何处。我猜你是想说，这些文稿未经修改且未经过最终修订，因此更容易被\Person{欧瑟伯}篡改。也许我愚钝，但这论证在我看来有些愚蠢且不着边际。若文稿未经修改且未经过最终修订，其中的错误应归咎于你的懈怠与拖延，而非\Person{欧瑟伯}；你所能归咎于他的，无非是他将你打算日后修改的文稿在基督徒团体中传播。但若如你所断言，\Person{欧瑟伯}篡改了它们，那你为何又提出它们未经修改、在最终修订前就已流传出去的指控？因为无论文稿是否经修改，同样可能被篡改。但你却说，\Quote{这些书卷无人拥有，或只有极少数人拥有}。这一句话本身就充满矛盾！若无人拥有这些书卷，它们又怎会落在少数人手中？若少数人拥有它们，你为何又虚假声称无人拥有？再者，当你说少数人拥有它们，且你自己承认\Quote{无人拥有}的说法不成立时，你抱怨你的秘书被金钱收买又有什么意义？告诉我们秘书的名字、贿赂的金额、地点、中间人、受贿者。当然，这个叛徒已被你抛弃，犯下如此大罪的人已与你断绝一切往来。\Person{欧瑟伯}所得到的作品副本，难道不更可能是你提及的那少数几位朋友给他的吗？尤其是这些副本彼此完全吻合，连一笔之差都没有。我们还可以问一问，将你尚未修改的副本交给他人是否明智？这些文献尚未得到最终修订，而其他人却拥有了这些需要修正的错误。难道你看不出你的谎言无法自圆其说吗？此外，在那特定时刻，你从中能得到什么好处——这又如何能帮助你逃脱主教们的谴责——使讨论中的书卷公开给所有人，让你被自己的话驳斥？由此看来，正如那位著名演说家的警句所说，你固然有说谎的意愿，却无编造的艺术。\par

% structure: p0016 source=h2 target=subsection reason=relative-heading-rank; base=h2

\subsection{你讥讽我夸耀自己的口才，难道你要夸耀自己的无知吗？}

6. 我将按你信件的顺序，并附上你亲口所说的话。\Quote{我承认，如你所说，在我的序言中我赞扬了你的口才；现在我也乐意再次赞扬，只是你违背了你的\Person{图利乌斯}的劝言，因过分自夸而使其变得可憎。} 我哪里夸耀过自己的口才？我甚至没有欣然接受你所给予的赞美。也许你说这话的原因是，你自己不愿被心怀诡计的公开赞美所奉承。请放心，你将被公开指控；你拒绝赞美你的人；你将领教一个公开指控你的人。我还不至于愚蠢到批评你那粗陋的文风；没有人能像你每次写作时那样强烈地将其暴露于谴责之下。我只是想让那些与你一样缺乏文学修养的同门看看，你在\Place{东方}三十年里取得了何等进步——一个不学无术的作者，将厚颜无耻当作口才，将遍及四方的恶语视为良心的标志。我并不打算举起戒尺；我并不像你所说的那样，自以为能挥动皮鞭教一个年迈的学生识字。但事实上，你的口才和教导如此闪耀，以至于我们这些区区小册子作家无法承受，你用才华的敏锐眩惑了我们的眼睛，以至于我们都必须显得嫉妒你；而且我们必须真的联手压制你，因为一旦你作为作家在我们中间取得首位，站在修辞拱门的顶端，我们所有声称略知一二的人将不被允许嘟囔一个字。按照你的说法，我是一个哲学家、演说家、文法家、辩证学家，一个通晓希伯来语、希腊语和拉丁语的人，一个\OriginalTerm{三语者}{trilingual}。依此估算，你也是一个\OriginalTerm{双语者}{bilingual}，你的拉丁语和希腊语知识足以让希腊人以为你是拉丁学者，让拉丁人以为你是希腊人；而\Person{埃皮法尼乌斯}主教将是一个\OriginalTerm{五语者}{pentaglossic}，因为他用五种语言说话反对你和你的宠儿。但我惊讶于你的鲁莽，你竟敢对我这样一个你自称认为如此博学的人说：\Quote{你，你的成就赋予你如此多的警惕之眼，如果你犯错，怎能被宽恕？你怎能不被掩埋于无尽耻辱的沉默？} 当我读到此处，并想到我必定在某处言语上有了过失（因为\BibleQuote{若有人在言语上不犯罪，他便是完人}），正期待他要揭露我的一些过失时，我突然看到了这样的话：\Quote{在这封信的送信人出发前两天，你那篇攻击我的演说辞交到了我手中。} 那么你那些威胁以及你的话：\Quote{如果你犯错，怎能被宽恕？你怎能不被永远耻辱的沉默所覆盖？} 又到哪里去了？然而，也许尽管时间短暂，你仍能将其整理好；或者你打算雇用一个博学之类的人，他期望在我的作品中发现你那类口才的装饰和宝石。你之前写道：\Quote{请接受我寄去的文件，你本想以高价购买的；}但现在你却带着谦卑的伪装说话。\Quote{我本想效法你的榜样；但是，由于返回你处的信使又匆匆折回，我认为简短地写信给你，比长篇大论地写给其他人更好。} 同时，你厚颜无耻地以你的不学为乐。事实上，你曾承认这一点，宣称\Quote{既然承认每部分都有错误，再注意文风上的几个小错便是多余的。} 因此，我不会因你把\Quote{获得}写成\Quote{购买}而指责你（尽管\Quote{获得}用于同类事物，而\Quote{购买}意味着数钱）；也不会指责像\Quote{那位返回你处的人又匆匆折回}这样的句子——这是一种与最拙劣文体相称的冗余。我只回应论点，并将定你的罪，不是出于\OriginalTerm{语法错误和粗鄙语}{solæcisms and barbarisms}，而是出于谎言、狡诈和厚颜无耻。\par

% structure: p0018 source=h2 target=subsection reason=relative-heading-rank; base=h2

\subsection{你希望先赞扬我，然后改正我，但两者都用拳头；并使我无法保持沉默。}

7. 若你真写信给我以劝诫我，因为你希望我改过，且不愿人们受绊脚石阻路，也不愿有人疯狂、有人缄默；那么你为何写信给他人反对我，又差遣手下散布于全世界让人阅读？你所设的困境——你试图以此困住我——又如何收场？\Quote{最高明的大师啊，你以为要纠正谁呢？若你写信给那些人，他们并无过错；若你指控的是我，你却没有写信给我。}我也要用你自己的话回答你：无学问的大师啊，你希望纠正谁呢？是那些没有犯错的人？还是你未曾写信给我的我呢？你以为你的首领们愚蠢迟钝，完全无法理解你的狡诈，或者说你的恶意——(因为蛇正是以此比乐园中的一切野兽更狡猾)——你要我私下劝诫你，你却公开指控我。你不以称此指控为\WorkTitle{辩词}为耻；你抱怨我以盾牌对抗你的匕首，并满怀虔诚与假正经，戴上谦卑的面具说：\Quote{若我犯了错，你为何写信给他人，而不尝试驳斥我呢？}我要反问你这一点。你抱怨我未做的事，你自己为何不做？就如一人用拳脚攻击他人，见对方要还手，便对他说：\Quote{你不知道这命令：\BibleQuote{若有人打你的右脸，连左脸也转过来由他打}吗？}于是，我的好先生，你决心要打我、挖我眼睛；我稍一动弹，你便援引福音的教训。你愿意我揭露你的一切狡诈伎俩吗？——就是那些居住在废墟中的狐狸的伎俩，\Person{厄则克耳}论及它们写道：\BibleQuote{你的先知们像旷野中的狐狸，哦，\Place{以色列}。}让我使你明白你所做的事。你在序言中赞美我，其赞美方式反而成了指控我的根据；若我未曾声明自己与那仰慕者毫无关系，我就会被判为异端。在我反驳了你的指控（即你的赞美），且没有对你个人怀恶意，而是回应了指控而非指控者，并痛斥异端人士，以表明尽管你诽谤我，我仍是公教徒之后，你发怒、狂躁，并写出了最宏伟的作品反对我；你将其交给所有人阅读传诵，此时，来自\Place{意大利}、\Place{罗马}和\Place{达尔马提亚}的信件也到了我手中，一封比一封更清楚地表明，你先前在颂词中加给我的所有赞美究竟价值几何。\par

8. 我承认，我立即着手回应针对我的指责，竭尽所能证明我不是异端，并将我的这些\WorkTitle{辩词}书卷寄给那些因你的书而痛苦的人，使你的毒药能被我的解毒剂所追随。对此，你寄给了我你早前的书卷，如今又寄来这最后一封信，满含侮辱与指控。我的好友，你期望我做什么？保持沉默？那便等于认罪。开口说话？但你将剑悬在我头顶，威胁要指控我，不再是在教会前，而是在法庭上。我做了什么应受惩罚的事？我在何处伤害了你？是因为我证明自己不是异端吗？还是因为我不能认为自己配得上你的赞美？或是因为我以直白的话语揭露了异端者的诡计和伪证？这一切与你何干？你自称是一位真人和公教徒，却在攻击我时比为自己辩护时更为热心。难道因为我为自己辩护，就必须被认为是攻击你吗？还是除非你证明我是异端，否则你就不可能是正统？与我关联能给你什么帮助？你行动的意义何在？你被一群人指控，却只通过攻击另一个人来回应。你发现一人攻击你，你却转身背对他，去攻击另一个本来要任由你独处的人。\par

% structure: p0021 source=h2 target=subsection reason=relative-heading-rank; base=h2

\subsection{为何你不能与我一同谴责\Person{奥力振}，以结束我们的争吵？}

9. 我呼求中保耶稣作证，我进入这场舌战的竞技场是违背自己意愿的，与必要性抗争；若不是你向我挑战，我绝不会打破沉默。即便如此，现在你若停止对我的指控，我的辩护也将停止。因为呈现给读者的一幕并不令人振奋：两位老人因一位异端者而角斗；尤其两人都希望被认为是公教徒。让我们停止对异端的一切偏袒，我们之间便再无纷争。我们曾热切赞美\Person{奥力振}；如今他被全世界谴责，让我们同样热切地谴责他。让我们携手同心，紧随东西方两位胜利者，稳步前进。我们年轻时走错了路，让我们年老时改正。若你是我的弟兄，就请为我看见自己的错误而高兴；若我是你的朋友，我必为你的悔改而喜悦。只要我们持续争斗，人们便会认为我们持有正确信仰并非出于自愿，而是出于必要。我们的敌意使我们无法展现真正的悔改。如果我们的信仰是同一个，如果我们都接受并拒绝相同的事物——（正如\Person{喀提林}所证实的，正是由此产生坚固的友谊）——如果我们同样憎恨异端，同样谴责我们过去的错误，那么我们为何要互相开战？我们攻击和辩护的对象是相同的。请原谅我在年轻时代，在完全了解他的异端之前，曾称赞过\Person{奥力振}研读圣经的热诚；我也将原谅你在鬓发斑白时，竟然为他的作品写了\WorkTitle{辩词}。\par

% structure: p0023 source=h2 target=subsection reason=relative-heading-rank; base=h2

\subsection{你声称只有两天时间回答是虚构的。}

10. 你声称，我的书在你写信给我之前两天才到你手上，因此你没有足够闲暇来答复。否则，若是你经过充分思考和准备之后才发言反对我，我们或许会认为你是在喷吐闪电而非指控。但是，像你这样诚实的人，也很难让人相信你告诉我们：一个经营东方货物的商人，他的生意是卖掉他从这些地方带来的货物，再购买意大利商品运回这里出售，只在\Place{Aquileia}停留了两天，以至于你不得不匆忙即兴地给我写信。因为那些你花了三年时间才彻底成形的著作，也未必比这写得更用心。不过，也许那时你身边没有人为你修改那些拙劣的作品，这就是为什么你的文学之旅缺少帕拉斯的助力，并且一路被文风瑕疵所中断，如同处处布满崎岖和裂缝。很明显，关于两天的说法是假的；你甚至不可能在那段时间里读完我所写的东西，更不用说回复了；因此，显然要么你写那封信花了好几天（其精心雕琢的文风也表明了这一点）；要么，如果这就是你匆忙写作的风格，而且你能如此即兴地写得这么好，那么你在有时间思考时反而写得差得多，就实在是太疏忽了。\par

% structure: p0025 source=h2 target=subsection reason=relative-heading-rank; base=h2

\subsection{你的翻译，与我的注释相反，为\Person{奥力振}的正统性作证。}

11. 你有些敷衍地说，你从希腊文翻译了我先前译成拉丁文的东西；但我并不清楚你指的是什么，除非你还是一直在对我提起关于\WorkTitle{关于 \BibleBook{厄弗所}{书} 的注释}，并且顽固地厚着脸皮，好像你没有得到任何答复似的。你塞住耳朵，听不见那用魔笛引诱者的声音。我在这部注释和其他注释中所做的，是既阐述我自己的意见，也阐述别人的意见，清楚指明哪些是公教的，哪些是异端的。这是通常的规则和习惯，凡是承担在注释中解释书卷的人都这样做：他们在讲解中详细列出各种意见，并说明自己和他人所想的。这不仅是在解释圣经的人中如此，在解释世俗著作的人中也是如此，无论希腊文还是拉丁文。然而，你不能因此就以\OriginalTerm{论原理}{\GreekText{Περὶ} \GreekText{᾿Αρχῶν}}来为自己辩护；因为你的自序已定你的罪，你在其中保证，邪恶部分和异端者增添的部分已被删去，但所有最好的部分都保留下来；因此，你所写的，无论是好是坏，都必须被视为不是你翻译的作者的作品，而是你自己——译者——的作品。也许，你本应纠正异端者的错误，并公开展示奥力振的错谬。但在这一点上（因为你把我引向文件本身），我在读你的信之前就已经答复了你。\par

% structure: p0027 source=h2 target=subsection reason=relative-heading-rank; base=h2

\subsection{你试图通过错误地将为\Person{奥力振}的辩护归给\Person{彭斐略}来庇护他。}

12. 关于帕姆菲鲁斯的书，我的遭遇，并非如你所说的滑稽，而或许是可笑：即在我断言此书系欧瑟比乌斯而非帕姆菲鲁斯所作之后，我在讨论结束时声明，我多年来一直相信它是帕姆菲鲁斯所作，并且我曾从你处借阅过此书的一个抄本。你可从以下事实判断我多么不惧怕你的嘲笑：即使现在，我仍作同样的声明。我取用了你那份抄本，以为是帕姆菲鲁斯著作的副本。我信任你，视你为基督徒和隐修士：我未曾想象你会犯下如此邪恶的欺诈。但是，由于你翻译了奥力振的著作，世人开始议论奥力振的异端问题，我在查考该书抄本时更加谨慎，并在凯撒勒雅的图书馆中找到了欧瑟比乌斯所著\WorkTitle{为奥力振辩护}六卷。我一翻阅它们，便立刻发现你独自以殉道者之名出版的那本论圣子与圣神的书，并将其中大部分亵渎之言改为更佳意义的措辞。我看此事必定是由狄迪慕斯或你或别的什么人所作（很明显，你所作的是针对\OriginalTerm{《论首要原理》}{\GreekText{Περὶ} \GreekText{᾿Αρχῶν}}），此有决定性的证据：欧瑟比乌斯告诉我们，帕姆菲鲁斯本人并未发表过任何作品。因此，你必须说明你是从哪里得到你的抄本的；且不可为了避免我的指控而说是从某个已故之人那里得到的，或，因为你无人可指认，便点名一个不能为自己辩护的人。如果这条小溪源于你的书桌，那么其推论已足够明显，无需我明言。但假设此书的标题和作者姓名被其他某个奥力振的爱好者更改了，你将其转译成拉丁文有何动机？无疑是：藉由殉道者为奥力振所作的见证，众人应信赖奥力振的著作，因为它们已事先获得如此权威的证人的保证。但这最为博学之人的辩护尚不能满足你；你还必须亲自撰写一篇论著为他辩护，当这两份文件广泛流传后，你便放心地着手把\OriginalTerm{《论首要原理》}{\GreekText{Περὶ} \GreekText{᾿Αρχῶν}}本身从希腊文翻译出来，并在前言中称赞它，你说其中有些内容已被异端者败坏，但你通过研究奥力振的其他著作将其纠正了。随后，你又对我加以赞扬，目的是防止我的任何朋友指责你。你把我推为奥力振的吹鼓手，将我的口才夸上天，好把信仰拖入泥沼；你称我为同僚和兄弟，自称是我作品的效仿者。于是，一方面你吹嘘我翻译了奥力振的七十篇讲道以及他关于宗徒的一些短论，你说我如此润色，以致拉丁文读者在其中找不出任何与公教信仰相违之处；如今另方面你却指认这些书为异端；并且抹去你从前的赞美，指责这个当初你以为会充当你的盟友而加以赞誉的人，因为你发现他是你不忠之敌。我们两人之中，谁是殉道者的诽谤者？是我，说他不曾为异端者，且那本被众人谴责的书不是他所写；还是你，出版了一本由一个亚略异端者所写、将名字改为殉道者之名的书？希腊世界已起反感，这尚且不够；你还必须把这本书强加给拉丁人听，并尽你所能以翻译侮辱一位卓越的殉道者。你的意图无疑并非如此；不是要指控我，而是要叫我服务于为奥力振著作的辩护。但让我告诉你，那曾由一位宗徒之声赞扬过的罗马信仰，并不承认此类伎俩。一种曾由宗徒权威保证的信仰，即使天使宣报别的福音，也不能更改。因此，我的兄弟，无论这本书的伪造是出自你（如许多人相信的），还是出自他人（你或许试图说服我们如此，这样你就只犯了轻信之罪，相信了异端者的作品是殉道者的著作），请改变标题，使罗马人的无辜脱离这巨大的危险。使你成为将一位最卓越的殉道者谴责为异端者的工具，于你并无益处；使一位为基督流血的人骄傲地成为基督教信仰的敌人，也无益处。采取另一条途径吧：说，我发现了一本书，相信是殉道者的作品。不要害怕成为悔改者。我不会进一步逼迫你。我不会问你是从谁那里得到的；你可以随意点名某个已故之人，或说是在街上从一个不认识的人那里买来的：因为我不愿看到你被定罪，而是愿看到你悔改。你出现错误，比殉道者是异端者更可接受。无论如何，设法把你的脚从当前的纠缠中拔出：想一想，在将来的审判中，当殉道者们向你提出控诉时，你将如何回答。\par

% structure: p0029 source=h2 target=subsection reason=relative-heading-rank; base=h2

\subsection{在我的注释中，引用相反的意见表明两者都不是我的。}

13. 再者，你提出一项无人控告你的指控，并在无人责备你之处辩解。你说你在我的信里读到这些话：\Quote{我想要知道，是谁准许你在翻译一本书时，删去一些内容，更改另一些，再添加其他内容？} 然后你自行回答，并反驳我说：\Quote{我对你说：我请问，是谁准许你在你的\WorkTitle{注释}中，记下一些出自奥力振的话，一些出自阿波利纳里的话，一些是你自己的话，而不是全部出自奥力振，或全部出自你自己，或全部出自其他人？} 自始至终，当你针对别的事时，你对自己提出了非常强烈的指控；你忘了那句古老的谚语：说谎的人应有好的记忆力。你说我在我的\WorkTitle{注释}中，写下了一些出自奥力振的话，一些出自阿波利纳里的话，一些是我自己的话。那么，如果我以他人名义写下的话是阿波利纳里和奥力振的，你加给我的指控又有什么意义呢？——当我说\Quote{另一个人这样说}，\Quote{以下是某人的推测}，那个\Quote{另一个人}或\Quote{某人}就是指我自己？奥力振和阿波利纳里在解释、风格和教义上有着巨大的差异。当我在同一段经文下列出不同意见时，难道我应该被认为同时接受了这些矛盾的观点？但此事后文再提。\par

% structure: p0031 source=h2 target=subsection reason=relative-heading-rank; base=h2

\subsection{若你翻译诚实，奥力振的异端就不会归咎于你。}

14. 现在我问你这一点：谁曾指责你在奥力振的书中添加、更改或删去某些内容，并像拷问架上的人一样质问你：你翻译中所写的内容是坏还是好？你假装无辜，并用一个愚蠢的问题来回避真正问题的核心，这是没有用的。我从未因你为满足自己而翻译奥力振而指责你；我也做过同样的事，维克多利诺、希拉里和盎博罗削也都做过。但我指责你，是因为你写了一篇赞同他的序言，以巩固你对一部异端作品的翻译。你迫使我重走同样的路，在我自己划定的路线上行走。因为你在那篇序言中说，你删除了异端者添加的内容，并代之以好的内容。如果你除去了异端者的错误陈述，那么你留下或添加的内容必定是奥力振的，或是你自己的，而且你大概认为它们是好的。但你无法否认其中许多是坏的。你会说：\Quote{那与我有什么关系？} 你必须把它归咎于奥力振；因为我只是更改了异端者添加的内容。那么请告诉我们，你为什么删除异端者写的坏东西，却留下奥力振写的原封不动？难道不明显是你以异端学说的名义谴责了奥力振错误教义的部分，而接受了其他部分，因为你判断它们不是错误的，而是真实的，且与你的信仰一致？正是后者使我询问：你在序言中称赞的那些东西是好是坏？正是这些东西，你承认在删除了所有最坏的部分后，将它们作为完全好的东西留下；这样，我就如我所说，把你放在拷问架上，以致如果你说它们是好的，你就会被证明是异端者；但如果你说它们是坏的，你立刻会被问：\Quote{那你为什么在你的序言中称赞这些坏东西？} 并且我没有加上你狡猾地假装我问的那个问题：\Quote{你为什么通过你的翻译把邪恶的教义带到拉丁人的耳中？} 因为展示坏的东西有时不是为了教导人，而是为了警告人防备它们：这样读者可以警惕不要追随错误，而是轻视他所知道的邪恶，否则，如果它们不为人知，可能会成为他惊叹的对象。然而在此之后，你竟敢说我是这类作品的作者，而你作为一个纯粹的译者，如果你选择纠正任何东西，就超出了译者的职责；但如果你没有纠正任何东西，你就只是一个译者。如果你翻译\OriginalTerm{论原理}{\GreekText{Περὶ} \GreekText{᾿Αρχῶν}}时没有序言，你这样说就完全正确；正如希拉里在翻译奥力振的讲道时所做的那样，他小心地做到使其中的好坏都归于作者本人，而非译者。如果你没有夸耀你删除了最坏的部分、留下了最好的部分，你也许会以某种方式逃脱泥潭。但正是这一点使你的发明诡计归于乌有，并使你四面受缚，无法逃脱。我必须请你不必过低估计你的读者的智慧，也不要以为所有阅读你作品的人都如此迟钝，以至于看到你把真正的伤口弄成腐烂，却把膏药贴在健康的身体上时，不会嘲笑你。\par

% structure: p0033 source=h2 target=subsection reason=relative-heading-rank; base=h2

\subsection{你说意大利的主教们接受你关于复活的看法。我对此表示怀疑。}

15. 关于肉身的复活，你的看法我们已从你的\WorkTitle{辩护书}中得知。\Quote{没有肢体会被切除，也没有身体部分会被毁坏。} 这是你天真无邪地做出的明确公开声明，你说意大利所有主教都接受它。我本应相信你的话，但那本不是庞斐禄所写的书的内容使我怀疑你。而且我感到奇怪，意大利竟会赞同罗马所拒绝的；主教们竟会接受宗座所谴责的。\par

% structure: p0035 source=h2 target=subsection reason=relative-heading-rank; base=h2

\subsection{你轻率地说，无论德奥斐洛制定什么，你都同意。现在你得考虑你与他的敌人依西多尔的友谊。}

16. 你另写信说，你是通过我的信函得知德奥斐洛教宗最近发表了一份信仰阐释，但尚未送达你处，你承诺接受他可能写的一切。我不记得曾说过这话，或曾寄过此类信函。但你却同意你仍不确定的事，以及你不知其为何物、将如何发展的事，为的是避免谈论你十分了解的事，并且不受你已给予的同意的约束。德奥斐洛有两封书信，一封是主教会议信，一封是巴斯卦信，反对奥力振及其门徒，还有其他反对阿波利纳里和也反对奥力振的信，在过去两年左右，我翻译并交给了说我们语言的人，以建立教会。我不记得翻译过他的其他任何东西。但是，当你说你完全同意德奥斐洛教宗的意见时，你必须小心，不要让你的师傅们和门徒们听见，也不要冒犯那些称我为强盗、称你为殉道者的众多人士，也不要激怒那位给你写信反对厄丕法尼主教、并劝你坚守信仰真理、不因任何恐惧而改变意见的人。这封书信的完整版本由收到它的人保存着。此后，你照你的方式说：\Quote{即使你对我发怒，我也要满足你，如同你先前所说的那件事一样。}但你又说道：\Quote{你想要什么？你还有什么可以用你的演说之弓来射击吗？}然而，当我批评你令人不快的说话方式时，你却很气愤，而你却采用喜剧演员最低级的表达，在写作教会事务时采用只适合舞台上妓女及其情人的语言。\par

% structure: p0037 source=h2 target=subsection reason=relative-heading-rank; base=h2

\subsection{你谈到埃及主教保禄。我们接待了他，虽然他是奥力振派，却作客旅看待；他已与正统信仰合而为一。不仅德奥斐洛，连皇帝们也谴责奥力振。}

17. 现在，关于你提出的问题，即我何时开始承认德奥斐洛教宗的权威，并与他在信仰上保持共融。你自己回答说：\Quote{那么，我想，是在你支持他谴责的保禄，并尽最大努力帮助他，怂恿他通过皇帝的谕令恢复被主教法庭免去的主教职位的时候。}我不先为自己辩护，而是首先谈谈你在此处对另一个人造成的伤害。有什么人性或爱德可言，为别人的不幸而幸灾乐祸，并将他们的创伤暴露给世人？这就是你从那个撒玛黎雅人学来的功课吗？他把半死的人带回客栈，你理解这是把油倒进他的伤口，并支付店主的花费吗？你就是这样解释带回圈中的羊、找回的钱币、回头的荡子吗？假设你有权诽谤我，因为我对你造成了伤害，并且用你的话来说，激怒了你，刺激你口出恶言：一个默默无闻的人对你做了什么伤害，你要暴露他的伤疤，并在它们结痂时，通过施加这种不必要的痛苦来撕裂它们？即使他配得上你的谴责，你这样做就正当吗？如果我没有弄错，你想通过他打击的人（我讲的是许多人的公开意见）是奥力振派的敌人；你利用他们中一人的麻烦来显示你对两者的暴力。如果德奥斐洛教宗的决定如此令你喜悦，你认为废除主教法令是亵渎，那么你对其他被他谴责的人怎么说？你对阿纳斯塔西教宗又怎么说？你极其真实地断言，没有人认为，作为如此大城的主教，他会伤害一个无辜或不在场的人。我这样说，并非因为我自居为主教决定的审判官，或希望他们所作的决定被废除；而是说，让每个人按自己的风险做认为正确的事；他必须独自考虑他的审判将如何被审判。我们在隐修院中的职责是款待；我们以人性的微笑欢迎所有来我们这里的人。我们必须小心，免得再次发生玛利亚和若瑟在客栈找不到地方，耶稣被关在门外，并对我们说：\BibleQuote{我作客，你们没有收留我。}我们唯一不欢迎的人是异端者，而你们却唯独特的异端者是迎；因为我们的会规要求我们为来者洗脚，而不是讨论他们的功德。我的兄弟，请回想：我们所说的人曾宣认基督；想想那被鞭子抽打的胸膛；回忆他所受的监禁、黑暗、流放、矿场劳役，那么你就不会惊讶我们把他当作过往的客人来接待。难道我们应被你们视为叛逆，因为我们因基督的名\BibleQuote{给口渴的人一杯凉水喝}吗？\par

18. 我能告诉你一件事，这可能使他对我们更显可爱，却令你更为可憎。不久前，被逐出埃及和\Place{亚历山大}的异端派别来到\Place{耶路撒冷}，愿与他结盟，这样他们一同受苦，便可被归于同一异端。但他拒绝了他们的提议，蔑视并驱逐了他们：他告诉他们，他不是信仰的敌人，不会拿起武器反对教会；他先前的行动是出于恼怒，而非信仰不纯；他只想证明自己的清白，并非攻击他人。你声称认为一道推翻主教谕令的皇帝谕旨是一种渎神行为。那是获取该谕旨的人的责任。但你对那些自己已被定罪，却出入权贵宫邸，形成密集队列攻击一位代表基督信仰的人，作何感想？然而，关于我与\Person{德奥斐洛}的共融，我只需以你假称我伤害过的那人为证。他的信始终是寄给我的，你很清楚，甚至在你阻止它们转交给我时，你每天派信使到他那里，激烈地重复说他的对手是我最亲密的朋友，并讲述你现在无耻地写下的同样谎言，以便激起他对我的仇恨，并使他因假想的受伤而痛苦，进而在信仰的事上压迫我。但他是个审慎而有智慧的人，经过时间和经验，逐渐明白了我们对他的忠诚以及你对我们设下的阴谋。如果如你所宣称，我的追随者在罗马暗中策划反对你，并在你熟睡时偷走了你未校对的文稿；那么是谁煽动\Person{德奥斐洛}反对埃及的公敌？是谁获得了公侯们反对他们的谕旨，以及整个这一地区的同意？然而你夸口说你自幼就是德奥斐洛的听众和弟子，尽管他在成为主教之前，因天生的谦逊，从未公开教导；而你在成为主教之后，也从未到过亚历山大。然而，为了打击我，你竟敢说：\Quote{我不指控或更改我的师傅们。}如果那是真的，依我之见，那将对你基督徒的立场产生严重怀疑。至于我自己，你无权指责我谴责我以前的老师；但我敬畏\Person{依撒意亚}的这些话：\BibleQuote{祸哉，那些称恶为善、称善为恶，以暗为光、以光为暗，以苦为甜、以甜为苦的人！}然而你却饮你师傅们的蜜酒和毒药，你已背弃了你真正的师傅——宗徒，他教导说，若他自己或天使在信仰的事上犯错，也不可跟随。\par

% structure: p0040 source=h2 target=subsection reason=relative-heading-rank; base=h2

\subsection{对\Person{维吉兰修斯}的回应，我仅写了正确之言。我知道是谁煽动他反对我。}

19. 你提到了\Person{维吉兰修斯}。我不知道你对他做了什么样的梦。我在哪里说过他在\Place{亚历山大}因与异端者共融而被玷污？告诉我那本书，拿出那封信：但你绝对找不到任何这样的陈述。然而，你以一贯的轻率陈述，更确切地说是无耻的谎言，使你以为每个人都会相信你说的话，你补充道：\Quote{当你引用一段圣经经文反对他时，方式如此侮辱，以至我不敢亲口重复。}你不敢重复，是因为你可以通过沉默使指控显得更严重；并且，因为你的指控没有事实依据，你假装谦逊，以使读者以为你是出于对我的考虑，尽管你的谎言表明你不顾自己的灵魂。这段圣经经文是什么，竟然羞于从你那张最无耻的嘴里说出来？确实，你在神圣的经卷中能提到什么可耻的东西？如果你羞于开口，至少你可以写下来，那么我就会因我自己的话而被认定是轻率。我可能对所有其他点保持沉默，但我仍可以通过这一段证明你的厚颜无耻有多大胆。你知道我多么不害怕你的控告。如果你拿出你威胁我的证据，那么现在落在你身上的所有责备都会落在我身上。当我处理维吉兰修斯时，我已经向你作出了答复；因为他对我提出了同样的指控，而你先以友好的颂词伪装，后以敌意的控诉提出。我知道是谁煽动他对我狂言；我知道你的阴谋和恶习；我并非不知道他的头脑简单，这一点人人皆知。通过他的愚蠢，你对我的仇恨找到了发泄其怒火的出口；而我写信抑制它，以免你被看作是唯一拥有文学棍棒的人，这并不证明你有权捏造羞耻的表达，而这些表达在我的任何著作中都找不到。你必须接受并承认这个事实：那回应他狂躁的文件同时也激起了你的诽谤。\par

% structure: p0042 source=h2 target=subsection reason=relative-heading-rank; base=h2

\subsection{关于教宗\Person{安纳斯塔修斯}谴责你的信函，你会发现它是真实的。}

20. 关于教宗阿纳斯塔修斯的信函问题，你似乎站在了一块滑溜的地方；你步履蹒跚，不知该把脚放在哪里。你一会儿说那信必是我写的，一会儿又说它应由收信人转交给你。接着你又指责写信人不公；或者抗议说信是否他写的与你无关，因为你持有他前任的证明书，而且当罗马请求你给她亲临的荣誉时，你因对自己小镇的爱而藐视了她。你若怀疑那信是我伪造的，为何不向罗马宗座的档案库索取，然后发现主教并未写那信，就证明我显然有罪？那时你便不是试图用蛛网绑我，而是用粗绳的网紧紧捆住我了。但若那信确为罗马主教所写，你向未收信的人索取副本，而不向发信人索取，又往东方寻求证据而其来源在你本国，这实是愚蠢之举。你最好去罗马，向他当面理论，因他当着不在场又无辜的你的面施加了指责。你首先可以指出，他拒绝接受你的信仰说明——如你所说，全意大利都赞同的——并且没有使用你的文学棍棒去对付你所说的那些狗。其次，你可以抱怨他往东方寄了一封针对你的信，给你烙上异端的印记，并说因你翻译了奥力振的著作 \OriginalTerm{论原理}{\GreekText{Περὶ} \GreekText{᾿Αρχῶν}}，罗马教会——它原本单纯地接受了这部作品——面临失去从宗徒所学信仰纯正的危险；而且他竟敢谴责这书本身，尽管它因你的序言之证明而得到加强，这更激起了对你的敌意。如此伟大城市的教宗将此项指控加于你，或冒失地听信他人而接纳它，这绝非小事。你应该在街头上大声呼喊，一遍又一遍地叫道：\Quote{那不是我的书，或者，若是我的书，未校正的稿页被欧瑟伯偷走了。我出版的是不同的版本，我甚至根本没有出版；我没有给任何人，或者最多给了少数人；我的敌人如此无顾忌，我的朋友如此疏忽，以至所有副本都被他篡改了。} 我最亲爱的弟兄，这才是你该做的，而不是背对他，把谩骂的箭跨海射向我；因为我的受伤如何能治愈你的伤口？一个濒死的人看见朋友同死，就能得安慰吗？\par

% structure: p0044 source=h2 target=subsection reason=relative-heading-rank; base=h2

\subsection{已逝的西里修可能写了信支持你；活着的阿纳斯塔修斯则往东方写信反对你。}

21. 你出示了如今在基督内安眠的西里修的信，却轻视活着的阿纳斯塔修斯的信。你问：他写了（或者也许根本没写）信，而你不知道，这能对你有什么损害？若他真写了，你仍可满足，因你有全世界的证词支持你，且无人认为如此伟大城市的主教会伤害一个无辜的人，甚或只是一个不在场的人。你说自己无辜，但你的翻译使整个罗马战栗；你说你不在场，但这只是因为你受指控时不敢答辩。你如此畏惧罗马城的审判，以致宁愿将自己置于蛮族的入侵之下，也不愿接受一个和平城市的意见。假设去年的信是我伪造的，那么最近在东方收到的那些信又是谁写的呢？然而在后来的信中，教宗阿纳斯塔修斯对你如此恭维，你读到后，会更倾向于努力自卫而非控告我。\par

我希望你思考一下，你所回避的智慧以及雅提卡式的机锋和宗教写作中的辞藻是何等不可避免。你受他人攻击，被他们的谴责洞穿，然而你却狂暴地向我冲来，说：\Quote{我可以讲述你离开罗马的方式，以及当时关于你的意见和后来关于你的书写，你的誓言、你登船的地方、你虔诚地避免发虚誓的方式；这一切我都可以详述，但我决定隐瞒的比讲述的更多。} 这便是你悦人言论的标本。若此后我说了什么尖锐的话作为回答，你便以立即的放逐和刀剑威胁我。你是一个极雄辩的人，通晓修辞的一切伎俩；你假装要略过一些事，实则透露出来，这样你无法通过公开指控证明的事，便可以借看似搁置而使之成为罪行。这一切就是你的单纯；这就是你所谓饶恕朋友并把陈述留给审判庭的意思；你通过堆砌大量指控来饶恕我。\par

% structure: p0047 source=h2 target=subsection reason=relative-heading-rank; base=h2

\subsection{我离开罗马前往东方，并无如你所暗示的任何可指责之处。}

22. 若有人想听我从罗马出发的行程安排，情况如下：在八月，当季风刮起时，我与可敬的司铎文森提乌斯、我的年轻弟弟以及现今住在耶路撒冷的其他隐修士一起，在罗马港口上了船。我选择了自己的时间，在众多圣徒的陪伴下抵达了雷焦。我在斯库拉海岸稍作停留，听到了关于足智多谋的尤利西斯快速航行、塞壬的歌声以及卡律布狄斯无底漩涡的古老传说。那里的居民给我讲了许多故事，并建议我不要航向普罗透斯之柱，而要驶向约纳上岸的港口，因为前者适于仓皇逃命之人，后者则最适合被囚禁之人；但我宁愿取道马莱阿和基克拉泽斯前往塞浦路斯。在那里，我受到可敬的厄丕法尼主教的接待，你以他对你的证词为荣。我来到安提约基雅，在那里与宣信者主教保利尼乌斯共融，并由他护送前往耶路撒冷，我在严冬酷寒中进入该城。我在那里看到了许多奇妙之事，并亲眼证实了先前传闻之事。从那里我前往埃及，参观了尼特利亚的隐修院，并察觉了潜伏在隐修士唱诗班中的毒蛇。然后我急忙返回白冷，那里现在是我的家，并在救主的马槽和摇篮上倾注了我的香膏。我还看到了那不祥之湖。我没有贪图安逸和懒散，而是学到了许多先前不知道的事。至于罗马对我作了什么判断，或后来写了什么，你尽可畅所欲言，尤其是你有文件可依；因为我不应由你的言语来审判，你随意将其隐藏在谜语中或公然说谎，而应由教会的文件来审判。你可以看出我多么不怕你。如果你能拿出罗马主教或任何其他教会哪怕一份针对我的记录，我就承认自己犯有我在你身上所发现的所有不义。我可以轻易讲述你离开时的情形、你的年龄、航行的日期、你居住的地方、你的同伴。但我绝不会做我所责备你做的事，在教会人士之间的讨论中，编造出只配吵架泼妇的胡言乱语。这一点足以让你智慧的心灵记住：不要用那种会立刻反过来针对你自己的方法来对付别人。\par

% structure: p0049 source=h2 target=subsection reason=relative-heading-rank; base=h2

\subsection{厄丕法尼的确给了你平安之吻，但后来表明他已开始不信任你。}

23. 关于我们可敬的朋友厄丕法尼，你的这种奇怪的搪塞令人费解，你说他在给了你平安之吻并与你一同祈祷之后，不可能写信反对你。这就好像你争辩说，如果他在不久前还活着，他就不会死；或者好像同样肯定的是，他先谴责了你，然后在平安之吻之后，又绝罚了你。经上说：\BibleQuote{他们从我们中间出去，却不属于我们；否则他们必会留下来与我们同在。}宗徒吩咐我们在第一次和第二次劝戒之后要避开异端者；当然，这意味着他在被避开或被谴责之前，曾是教会羊群的一员。我不得不承认，当你在某个聪明人的提示下，开始唱赞美厄丕法尼的颂歌时，我忍俊不禁。怎么，这就是那个\Quote{老糊涂}、\Quote{拟人派}，就是那个在你面前夸耀自己读过奥力振六千卷书的人，那个\Quote{自以为受托在所有民族中用自己的语言宣讲反对奥力振的福音}，那个\Quote{不让别人读奥力振，唯恐他们发现他从奥力振那里偷了什么东西}的人。读读他所写的，以及那封信，或者更确切地说，那些信，其中一封我将作为你正统信仰的证明提出，这样就能看出他为你今日的赞美是多么相称。\Quote{我的兄弟，以及委托给你的基督的神圣子民，以及所有与你同在的弟兄，尤其是司铎鲁菲努斯，愿天主使你从奥力振的异端和所有其他异端中，以及它们所带来的毁灭中，获得自由。因为如果教会因一句或两句违背信仰的话而谴责了许多异端，那么一个发明了这么多乖谬之事、这么多虚假教义的人，更应被视为异端者！他显明是天主和教会的敌人。}这就是这位圣人给你的证词。这是他给你炫耀的赞美花环。这就是你的金币从我们的兄弟欧瑟比乌斯的房间中提取出来的信的内容，以便你能诽谤信的译者，并将一个最明显的罪责加在我身上——把一个希腊词译为\Quote{最亲爱的}，而它本应译为\Quote{可敬的}！但这与你何干？你以你的精明方法能左右一切，能在不同可能性之间调整你的道路：首先，如果你能找到人相信你，你会说阿纳斯塔修斯和厄丕法尼从未写过一句反对你的话；其次，当他们的信函本身向你呐喊，击垮你厚颜无耻的狂妄时，你又轻视他们两人的判断，说他们写不写对你无关紧要，因为他们不可能反对一个无辜和不在场的人写东西。\par

那么，你也没有权利说那位圣人\Person{圣埃皮法尼乌斯}的坏话，像你那样说\Quote{可见他口里和亲吻中给了我平安，心里却怀着恶意和欺骗}——因为这是你的推理，你这样为自己辩护。这是\Person{圣埃皮法尼乌斯}的信，且对你不利，天下皆知；它以其本来面目落入你手中，我们可以证明；因此你否认你所不能怀疑的真实，是惊人的无耻，或者说完全是厚颜无耻。什么！难道要让\Person{圣埃皮法尼乌斯}蒙受他给你和平之吻却心怀欺骗的污名吗？难道不是更真实地相信他起初劝诫你，是因为他愿意救你脱离错误，引你回到正道；因此他没有拒绝你的\Person{犹达斯}的亲吻，希望以他的宽容瓦解信仰的背叛者——但后来发现他的一切劳苦白费，豹子不能改变斑点，埃塞俄比亚人不能改变皮肤，他就在信中宣布了他先前心中只是怀疑的事？\par

% structure: p0052 source=h2 target=subsection reason=relative-heading-rank; base=h2

\subsection{当我们作为朋友分别时，我相信你是一个真正的信徒；没有人被派往\Place{罗马}去伤害你。}

24. 你用来攻击教宗\Person{亚纳斯大削}的论证也大致相同，即既然你持有主教\Person{西利奇乌斯}的信函，他就不可能写信反对你。我担心你以为你受到了一些伤害。我不明白以你这样的敏锐和能力，怎么会降尊纡贵到这种荒谬的地步；你以为你的读者是愚昧的，但你自己却显得愚昧。然后在这非凡的论证之后，你加上了这样一小句：\Quote{这样的行为绝不属于这些可敬的人。正是出自你们这一学派。在我们分别时，你们给了我们一切和平的记号，然后从我们背后投掷充满毒液的武器。}在这一句，或者更确切地说，在夸张的演讲中，你无疑想展示你的修辞技巧。确实，我们给了你和平的记号，但不是为了拥抱异端；我们握了手，送你上路，是基于你是大公的，而不是我们是异端的。但我想知道这些中毒的武器是什么，你抱怨我从你背后投掷。我派了司铎\Person{文森爵}、\Person{保利尼亚诺}、\Person{欧瑟伯}、\Person{鲁菲诺}。其中，\Person{文森爵}远在你之前就去了\Place{罗马}；\Person{保利尼亚诺}和\Person{欧瑟伯}在你航行一年后出发；\Person{鲁菲诺}两年后，为了\Person{克劳狄}的事；他们所有人要么出于私人原因，要么因为另一个人生命危险。我怎么可能知道，当你进入\Place{罗马}时，一位贵族梦见一艘满载货物的船张满帆驶入？或者所有关于命运的问题都被一种本身并非愚昧的解答解决了？或者你正在翻译\Person{欧瑟伯}的书，仿佛它是\Person{庞斐禄}的？或者你正在用自己华丽的文采把这个臭名昭著的著作\OriginalTerm{论原理}{\GreekText{Περὶ} \GreekText{᾿Αρχῶν}}的翻译当作盖子盖在\Person{奥利振}的有毒菜肴上？这是一种诽谤人的新方法。我们在你犯罪之前就派出了控告者。我再说一遍，不是出于我们的计划，而是出于\OriginalTerm{天主眷顾}{providence}，这些因其他原因被派出去的人，前来对抗正在兴起的异端。他们被派去，就像\Person{若瑟}一样，要藉着他们信德的热情来缓解即将来临的饥荒。\par

% structure: p0054 source=h2 target=subsection reason=relative-heading-rank; base=h2

\subsection{你发誓说你没有写我的所谓撤回书。你的文风暴露了你，关于我的翻译我已经充分回答。}

25. 一旦挣脱了约束，厚颜无耻会爆发到何种地步？他把针对别人的指控归于自己，以使人们以为是我们捏造的。我对某个未指名的人提出了指控，而他当作是针对自己的；他洗净自己别人的罪，只确信自己的无辜。因为他发誓说他并没有写那封以我的名义发给\Place{阿非利加}主教们的信，在那信中我承认自己因犹太人的影响而错误地翻译了圣经；他给我送来写有所有这些他所声称不知情的东西的文件。值得注意的是，他的精明如何与另一个人的恶意相吻合，以至于这个在\Place{阿非利加}所说的谎言，他一致地宣称是真的；而且他那优雅的文风如何能被某个偶然的、没有经验的人模仿。你独享翻译异端者的毒素的特权，使万民喝\Place{巴比伦}杯中的酒。你可以从希腊文修正拉丁文圣经，并将与宗徒传给教会的不同读物交给教会阅读；但我却不允许超越我多年前为同语言的人严格校正后翻译的\OriginalTerm{七十贤士译本}{Septuagint}，并为驳斥犹太人而翻译他们所承认为最准确的圣经真本，这样当犹太人与基督徒之间发生争论时，他们就无处可退和遁词，只能被自己的矛最有效地刺中。如果我记得不错，我在许多其他地方，特别是在我第二卷的末尾，已经相当充分地写到了这一点；我用清楚的事实陈述制止了你那追逐声誉之举，你企图以此在无辜和无经验的人中激起我的恶意。我认为将读者引向那里就足够了。\par

% structure: p0056 source=h2 target=subsection reason=relative-heading-rank; base=h2

\subsection{你提醒我提防篡改和背叛。你是在警告我防你自己。}

26. 我认为这一点不应被忽略：你没有权利抱怨，你文件的伪造者在我眼中享有宣信者的荣耀地位，因为你自己犯有同样的罪行，却被奥力振的所有党羽称为殉道者和宗徒，由于你在亚历山大的流放和监禁。关于你所谓不谙拉丁文写作一事，我在上面已经答复了你。但是，既然你重复同样的事情，并且仿佛忘记了你先前的辩护，再次提醒我应当知道，三十年来你一直忙于吞噬希腊文书籍，因此不谙拉丁文，我希望你注意，我并不是挑剔你几个词的问题，尽管你所有的著作都值得被销毁。我想要做的是向你的追随者表明——你花费如此多心血教导他们什么都不知道——让他们明白，一个人教导他所不知道的、写作他所无知的东西，他会有多少谦逊；这样他们可以期望在他的意见中找到同样的智慧。至于你所添加的\Quote{它并非言辞的过错令人反感，而是罪过，如说谎、诽谤、贬低、假见证，以及一切恶言，说谎的口杀害灵魂}，以及你的恳求\Quote{不要让那恶臭到达我的鼻孔}；我会相信你的话，要不是我发现事实与此不符。这正如一个漂洗工或鞣皮工对颜料商说话，警告他最好在经过他们的店铺时捏住鼻子。我将按照你的建议行事；我将捏住我的鼻子，以免它因你那说真话和祝福的愉快气味而受折磨。\par

% structure: p0058 source=h2 target=subsection reason=relative-heading-rank; base=h2

\subsection{在某事上称赞而在另一事上责备一个人并无矛盾。你在我身上就是这样做的。}

27. 关于你对我交替的称赞和贬低，你非常尖锐地争辩说，你对我有同样的权利说好说坏，正如我有权利挑剔奥力振和狄迪莫——我曾赞扬过他们。那么，我必须教导你，最有智慧的人，罗马辩证学家的首领，在某一事上赞扬一个人而在另一事上责备他并无逻辑上的过失，只有在对同一件事既赞同又不赞同时才成问题。我将举一个例子，以便虽然你可能不明白，但明智的读者能与我一同理解这一点。在戴尔都良身上，我们赞扬他的伟大才华，但谴责他的异端。在奥力振身上，我们钦佩他对圣经的知识，但我们不接受他的虚假教义。至于狄迪莫，我们却赞扬他记忆的才能和他对天主圣三信仰的纯正，而在另一方面，他在信任奥力振时所犯的错误，我们则与他分离。我们不应效法老师的恶习，而应效法他们的美德。在罗马有一个人，他的文法教师是一个非洲人，非常博学；他以为模仿教师尖锐的声音和错误的发音就能与教师相等。你在你为\OriginalTerm{论原理}{\GreekText{Περὶ} \GreekText{᾿Αρχῶν}}所作的序言中，称我为你的兄弟，称我为你最雄辩的同僚，并宣告我信仰纯正。从这三点你无法后退；其他一切方面你尽管可以挑剔我，只要你没有公开反驳你为我所作的这份见证；因为当你称我为朋友和同僚时，你承认我值得你的友谊；当你宣告我是雄辩之人时，你就不能再指责我无知；当你宣告我在各方面都是公教徒时，你就不能把异端的罪归给我。除了这三点，你可以随意指控我任何事而不公开自相矛盾。所有这些计算最终的结果是，你在我身上责难你从前赞扬的事，这是错误的；但我并没错，因为对同一些人，我赞扬可赞扬的，责难可责难的。\par

% structure: p0060 source=h2 target=subsection reason=relative-heading-rank; base=h2

\subsection{我对许多自然现象的无知不能成为你在灵魂起源问题上无知的借口。根据你自夸的梦，你应当知道一切。最重要的东西在你那船东方货物中被遗忘了。}

28. 你接着谈到灵魂的起源，并长篇大论地反对我所说的你扬起的烟雾。你希望被允许在这一点上表达无知，而你却故意隐藏你的知识；因此开始问我关于天使和总领天使的事；关于他们存在的方式、他们居所的位置和性质、他们之间是否存在差异；然后关于太阳的运行、月亮的盈亏、星辰的性状和运行。奇怪的是你没有列出所有的诗行：\par

地震从何而来，澎湃的潮水从何而来？\newline{}冲破堤岸，然后又平息；\newline{}日蚀，月亮的辛劳；\newline{}人类和兽类的种族，风暴，火焰，\newline{}大角的雨星毕宿，以及大熊星：\newline{}为何冬日的太阳匆匆沐浴于波浪之下？\newline{}又是什么挡住了它的长夜？\par

然后，离开天上的事，屈尊就卑论及地上的事，你开始在琐碎问题上卖弄哲学。你说：\Quote{告诉我们泉源和风的原因是什么；什么造成冰雹和阵雨；为什么海是咸的，河是甜的；云和风暴、雷电、雷声和闪电如何解释。}你的意思是，如果我不知道这一切，你就有权说你对于\OriginalTerm{灵魂的起源}{origin of souls}一无所知。你希望用我在许多事上的无知来平衡你在单一问题上的无知。但是，你在页页升起的所谓我的烟雾中，难道不明白我能看见你的迷雾和旋风吗？你希望被认为是一个知识渊博的人，并在\Person{卡尔普尔尼乌斯}的门徒中享有智慧的盛名，因此你在我面前抬起了整个物质世界，仿佛\Person{苏格拉底}在转向伦理学研究时所说的：\Quote{在我们之上的事与我们毫无关系。}是白说了一样。那么，如果我无法告诉你，蚂蚁这种小生物，身体不过是一个点，为何有六只脚，而大象体积庞大却只有四只脚行走；为何蛇和蟒蛇用胸部和腹部滑行；为何通常称为千足虫的虫子有那么多的脚；我就被禁止知道关于灵魂的起源的任何事情！你问我关于灵魂我知道什么，以便我一有所陈述，你就立刻攻击。如果我说教会的道理是天主每日形成灵魂，并将它们投入出生者的身体，你就立刻拿出你的导师所发明的网罗，问道：如果天主将灵魂赐给由奸淫或乱伦所生的人，天主的公义何在？祂岂不是成了人罪的共犯？如果祂为那些制造身体的奸淫者创造灵魂。就好像，当你听说种子被偷了，你就认为过错在于种子的性质，而不在于偷麦子的人；因此大地就不该在怀抱中滋养种子，因为撒种者的手不洁一样。由此也引出你神秘的问题：为什么婴儿会死？既然按你的观点，他们是因为自己的罪才获得了身体。有一篇\Person{狄迪莫斯}写给你的论文，他在其中回应了你的这一询问，回答说他们并未犯许多罪，因此仅仅接触他们肉身的牢狱就是足够的惩罚了。他——曾是你我的导师——当你问这个问题时，应我的请求写了三卷关于\Person{欧瑟亚}先知的注释，并题献给我。这表明我们各自接受了他教导的哪些部分。\par

29. 你敦促我表达对事物本性的看法。若有余裕，我可以向你复述\Person{卢克莱修}追随\Person{伊壁鸠鲁}的观点，或\Person{亚里士多德}为\OriginalTerm{逍遥学派}{Peripatetics}所教导的观点，或\Person{柏拉图}与\Person{芝诺}为\OriginalTerm{学园派}{Academics}与\OriginalTerm{斯多亚派}{Stoics}所教导的观点。转到教会，那里我们有真理的标准，\BibleBook{}{创世纪}、\BibleBook{}{先知书}和\BibleBook{}{训道篇}在这些问题上给了我们许多信息。但如果我们承认对所有这一切——同样也包括对\OriginalTerm{灵魂的起源}{origin of souls}——一无所知，那么你理应在你的\WorkTitle{辩护书}中承认你对这一切都同样无知，并质问你的诽谤者，他们自己对这些重大之事一无所知，为何竟敢迫使你就这单一问题作答。但是，哦！那条满载埃及和东方一切货物、前来充实\Place{罗马}城之贫穷的三层桨战船，所蕴藏的财富是何等巨大啊！\par

你就是那位英雄，名副其实的\Person{马克西穆斯}，\newline{}你独自以写作拯救了国家。\par

除非你从东方来，那位博学之士仍会困在\OriginalTerm{占星术士}{mathematici}之中，所有基督徒仍会对反对\OriginalTerm{宿命论}{fatalism}的论点一无所知。当你将满载这类货物的一艘船带上岸时，你就有权用关于占星术以及天空和星辰的原因的问题来烦扰我。我承认我的贫穷；我不像你那样在东方变得如此富有。你在\Place{法罗斯岛}的荫庇下长期居住所学到的东西，\Place{罗马}从未知道：\Place{埃及}将\Place{意大利}直到现在才拥有的知识教导了你。\par

30. 你的\WorkTitle{辩护书}说关于\OriginalTerm{灵魂的起源}{origin of souls}有三种意见：一种由\Person{奥利振}所持，第二种由\Person{德尔图良}和\Person{拉克坦提乌斯}所持（至于\Person{拉克坦提乌斯}，你所说的显然是假的），第三种由我们这些愚钝之人所持，他们看不出，如果我们的意见是真的，天主就被显为不义。之后你说你不知道什么是真理。那么，我说，告诉我，你是否认为在这三种意见之外可以找到任何真理，以致这三种可能都是错误的？或者你是否认为这三种之中有一种是真的？如果还有其他可能，你为什么将讨论的自由限制在窄小范围之内？你为什么只提出错误的观点，而对正确的保持沉默？但如果三者中有一真而二假，你为什么将真假混同在一句无知的断言中？也许你假装不知道哪一个是真，以便你可以在任何喜欢的时候安全地维护错误。这就是烟，这就是雾，你企图用它们来遮蔽人们眼中的光明。你是我们时代的\Person{亚里斯提卜}：你将你的船驶入\Place{罗马}港，满载各种货物；你将教授座椅高高架起，向我们展现\Person{赫尔玛戈拉斯}和莱昂提尼的\Person{高尔吉亚}：只是，你如此急于起航，竟将一件小货物、一个小问题遗忘在东方。你反复大声说，在\Place{阿奎莱亚}和\Place{亚历山大港}你都学到天主是我们身体和灵魂的创造者。那么，说真的，这紧迫的问题是：我们的灵魂是由天主创造的还是由魔鬼创造的？而不是\Person{奥利振}的意见是否正确——我们的灵魂在身体之前存在，并犯了某些罪，因此被束缚在这些粗重的身体里？或者，再次，它们是否像睡鼠一样处于麻木沉睡的状态？每个人都在问这个问题，但你却闭口不谈；没有人问另一个问题，但你的回答却针对那个。\par

31. 你所常讥笑的我之\Quote{烟}的另一部分，依你所说，是我假装知道我所不知道的事，并且炫耀伟大的导师来欺骗无知的平民。当然，你是一个不是烟而是火焰——甚或闪电——的人；你说话时雷霆万钧；你无法抑制口中生出的火焰，就像犹太叛乱的领袖\Person{巴尔科赫巴}，他常将一根点燃的稻草含在口中，吹出来以显出喷火的样子；同样，你也像第二个\Person{萨尔莫纽斯}，照亮你所踏上的整条路，并斥责我们仅为烟人，或许那句话可以应用在我们身上：\Quote{你触摸山岳，它们就冒烟。}你不明白先知提到蝗虫之烟时的寓意；无疑是你眼睛之美使你无法忍受我们烟的辛辣。\par

% structure: p0069 source=h2 target=subsection reason=relative-heading-rank; base=h2

\subsection{你的梦是自夸：而你所指责我的梦却使我谦卑。}

32. 至于你对发假誓的指控，既然你引我参考你的书；既然我在前几卷中已向你及\Person{卡尔普尼俄斯}作了答复，这里只须指出：你要求我在睡梦中做到的事，你醒着时却从未实现过。似乎我犯了大罪，因为我告诉基督的少女和贞女们最好不读世俗作品，而且我曾在一场梦中受警告承诺不再读它们。但你的船——藉启示向罗马城宣告——承诺一件事，却行另一件事。它来是为了解除数学家的难题：它所做的却是解除基督徒的信德。它满帆航行穿过爱奥尼亚海、爱琴海、亚得里亚海和第勒尼安海，却只在罗马港口搁浅。你难道不羞愧于搜罗这类废话，并使我费力去提出类似的事反对你吗？假设有人梦见关于你的事，可能会使你虚荣；你明智且谦逊的做法是不显得知道它，而不是吹嘘别人的梦作为对你自己的严肃见证。你的梦和我的梦多么不同！我的梦告诉我如何被谦卑和抑制；你的梦一再夸耀你如何被赞美。你不能说：别人梦见什么与我无关，因为在你那些最有启发性的书中，你告诉我们这正是促使你翻译的动机；你不能容忍一位杰出的人白白做梦。这便是你的一切努力。如果你能证明我犯了假誓之罪，你以为你就会被认为不是异端者了。\par

% structure: p0071 source=h2 target=subsection reason=relative-heading-rank; base=h2

\subsection{首先揭露你异端的不是我，而是很久以前的\Person{埃皮法尼乌斯}和在他之前的\Person{阿特比乌斯}。}

33. 现在我来到最严重的指控，即你指责我在我们友谊恢复之后不忠。我承认，在你对我提出或威胁我的所有谴责中，我最不愿接受的就是欺诈、欺骗和背信。犯罪是人性，设陷阱是魔鬼性的。什么！难道我为了这个缘故，在复活教堂内与你在被宰杀的羔羊上握手，好让我在罗马偷你的手稿吗？或者好让我放出我的狗去咬碎你尚未校订的文稿？谁能相信我们在你犯罪之前已经准备好了控告者？我们岂能知道你在心中酝酿什么计划？或者另一个人梦见什么？或者希腊谚语在你身上应验：\Quote{猪教\Person{密涅瓦}}？如果我派\Person{欧瑟比乌斯}去攻击你，那么又是谁激起了\Person{阿特比乌斯}和其他人对你的愤怒？难道不是事实：他认为我因与你的友谊也是异端者？而当我向他解释了\Person{奥力振}的异端之后，你把自己关在家里，从不敢见他，唯恐你不得不谴责你不想谴责的，或者因公开抵抗他而遭受异端的谴责。你以为他不能作为见证人反对你，因为他是你的控告者？在可敬的\Person{埃皮法尼乌斯}主教来到耶路撒冷，用言语和亲吻给你和平的记号之前，他心中却怀着恶念和诡诈；在我为他翻译那封谴责你的信，他在信中写你为异端者，尽管他先前曾认可你为正统；\Person{阿特比乌斯}已在耶路撒冷攻击你，如果不是他迅速离去，他本会感受到的不是你文学上的棍棒，而是你右手挥舞以驱赶狗的手杖。\par

% structure: p0073 source=h2 target=subsection reason=relative-heading-rank; base=h2

\subsection{至于我们对\OriginalTerm{《论原理》}{\GreekText{Περὶ} \GreekText{᾿Αρχῶν}}的译文，你的译文造成伤害，而我的译文是出于自卫的必要。你应该高兴异端被揭露。}

34. \Quote{但为什么，}你问，\Quote{你接受了我那被篡改的手稿？并且，当我翻译了\OriginalTerm{论原初原理}{\GreekText{Περὶ} \GreekText{᾿Αρχῶν}}之后，你竟敢也动笔做同一工作？如果我犯了错误（这是人之常情），难道你不该用一封私人信件召我来回答，并和颜悦色地对我说话，如同我现在在这封信里和颜悦色地说话一样吗？}我的全部过错在于，当有人用伪善的赞扬来指控我时，我试图为自己辩解，而这并没有恶意地提及你的名字。我希望把仅由你提出的指控归于许多人，不是为了将异端的指控回掷给你，而是为了从我自己身上驳斥它。我怎能知道，如果我写文章反对异端者，你会生气呢？你曾说过，你已经从\Person{奥力振}的著作中删除了异端段落。因此，我把攻击转向异端者而非你，因为我不相信你是异端的袒护者。请原谅我，如果我做得过于激烈。我以为这会令你高兴。你说，是由于替我办事者的欺诈伎俩，你的手稿才被公之于众，而它们原本秘密地藏在你的房间里，或者只归那个要求翻译的人所有。但这如何与你先前所说要么无人要么极少人拥有它们一致呢？如果它们秘密地藏在你的房间里，又怎能归那个要求翻译的人所有？如果那唯一被委托写手稿的人为了隐藏而拿到了它们，那么它们就不是秘密地藏在你的房间里，也不在你现在宣称拥有它们的少数人手中。你指控我们偷走了它们；然后你又责备我们用一大笔钱和巨额贿赂买下了它们。同一件事，同一封短信中，何等杂乱矛盾的谎言！你完全有指控的自由，而我却没有辩护的自由。当你提出指控时，你根本不考虑友谊。当我开始回答时，你满脑子都是友谊的权利。让我问你：你写这些手稿是为了隐藏还是为了发表？如果是为了隐藏，为何要写？如果是为了发表，为何要隐藏它们？\par

35. 但你会说，我的过错在于没有阻止那些是我的朋友的你的控告者。为什么？我忙于回答他们对我的指控；因为他们指控我虚伪——我可以用出示他们的信来证明——因为我明知你是异端者却保持沉默；并且因为我不谨慎地与你维持和平，从而助长了教会的内战。你称他们为我的门徒；他们怀疑我是你的同道门徒；并且，因为我在拒绝你的赞扬时有些保留，他们认为我和你们一起被引入了异端的奥秘。这就是你的序言为我所做的好事；你以朋友的面貌出现比以敌人的面貌出现更能伤害我。他们已经一劳永逸地说服了自己（不论对错，那是他们的事）你是异端者。如果我决定为你辩护，只会使自己和他们一起被你指责。他们用你对我的赞美来责备我，他们认为这些赞美不是出于诡计而是出于真诚；他们猛烈地指责我那些你一贯赞扬我的地方。我该怎么做？为了你的缘故，把我的门徒变成我的控告者？用我的头去接那些投向朋友的武器？\par

36. 关于\OriginalTerm{论原初原理}{\GreekText{Περὶ} \GreekText{᾿Αρχῶν}}这些书，我甚至对你有感激之情。你说你删除了任何冒犯性的部分，并用更好的内容代替。我完全按照希腊原文的内容来呈现。通过这种方式，两件事都显现出来了：你的信德和你所翻译之人的异端。罗马的主要基督徒们写信给我：回答你的控告者；如果你保持沉默，你就会被视为同意他的指控。他们全体一致要求我揭露\Person{奥力振}的微妙错误，并将异端者的毒害告知罗马人的耳朵，以使他们警惕。这怎么会对你造成伤害呢？难道你对这些书的翻译有垄断权吗？难道没有其他人参与这项翻译吗？当你翻译\WorkTitle{七十贤士译本}的部分时，你难道要禁止所有人在你的译本出版后再翻译吗？为什么？我也从希腊文翻译了许多书。你完全有权力随意对它们进行第二次翻译；因为其中的好坏都必须归咎于其作者。这在你身上也适用，如果你没有说过你已经删除了异端部分，只翻译了绝对好的内容。这是你为自己制造的一个困难，除非你承认自己犯了人常有的错误，并谴责你之前的意见，否则无法解决。\par

% structure: p0077 source=h2 target=subsection reason=relative-heading-rank; base=h2

\subsection{你为\Person{奥力振}的辩护并没有拯救他，反而使你卷入了异端。}

37. 然而，你为你所写的\Person{奥力振}著作辩护——更确切地说，你为\Person{优西比乌}的著作辩护，尽管你作了大量修改，并将一位异端者的作品冠以殉道者的名义翻译出来——但你却写下了更多与教会信仰不相容的内容。你和我一样将拉丁文著作译成希腊文；\Quote{你岂能禁止我将外国人的作品呈现给我本族的人呢？}如果我针对你的某部未曾攻击我的著作作出回应，那么人们或许会认为，在翻译你已经译过的作品时，我是出于敌意对待你，并希望证明你不够准确或不可信赖。但这是一类新的抱怨：当你因我回应你指责我的问题而心生不满时，情况便不同了。据说全罗马因你的译本而骚动；每个人都向我索求补救；这并非因为我有什么了不起，而是那些询问者如此认为。你说你是我的朋友，是你翻译了那作品。但你希望我怎么做呢？我们应当服从天主还是服从人？是守护我们主人的财产，还是隐瞒同仆的偷窃？难道除非我与你一同犯下招致谴责的行为，否则我就不能与你平安相处吗？如果你没有提到我的名字，如果你没有用你的奉承来装扮我，我或许还有逃避的途径，并能为不翻译已经译过的作品找出许多借口。但是，我的朋友，你迫使我在这一工作上浪费了不少时日，并将本应沉入\Place{卡律布狄斯}的东西公之于众；然而，尽管我受到了伤害，我仍遵守友谊的法则，并尽可能在不指控你的情况下为自己辩护。你所表现出的是一种过于多疑和抱怨的性情，当你将对异端者所说的话视为对自己的责备时。如果除非我是异端者的朋友，否则就不能成为你的朋友，那么我宁愿忍受你的敌意，也不愿忍受他们的友谊。\par

% structure: p0079 source=h2 target=subsection reason=relative-heading-rank; base=h2

\subsection{我的友好信件是为了防止不和：另一封是为了粉碎虚假见解。}

38. 你想象我编造了另一个谎言，即我以我自己的名义给你写了一封信，假装这信是许久以前写的，在信中我显得友善和谦恭；而你却从未收到过。真相很容易查证。过去三年来，罗马的许多人都有这封信的副本；但他们拒绝将它寄给你，因为他们知道你一直在散布诋毁我名誉的暗示，并编造最可耻、与我们的基督信仰不相称的故事。我在对此一无所知的情况下写信给你，如同写给一位朋友；但他们不愿将这封信转交给一位敌人——他们知道你就是如此——这样既免除了我错误的后果，也免除了你良心的谴责。接着，你提出论证，表明如果我写了这样一封信，我就没有权利再写另一封包含许多对你责备的信。但这是你说的所有话中常犯的错误，我有权对此表示不满：凡是我对异端者所说的话，你都想象是对你说的。\Quote{怎么！难道我因为给异端者一块石头砸碎他们的脑袋，就是在拒绝给你面包吗？}但是，为了证明你不相信我的信，你不得不认为教宗\Person{阿纳斯塔西乌斯}的信也基于类似的欺诈。关于这一点，我先前已经回答过你。如果你真的怀疑那不是他的手笔，你有办法证明我伪造。但如果那是他的手笔——正如他今年写的反对你的信也证明的那样——你妄图用虚假的推理来证明我的信是伪造的，那是徒劳的，因为我可以用他真实的信来证明我的信也是真实的。\par

% structure: p0081 source=h2 target=subsection reason=relative-heading-rank; base=h2

\subsection{我正确引用了\Person{毕达哥拉斯}。我提出他的一些言论。}

39. 为了反驳虚假的指控，你竟变得吹毛求疵。你没有被要求交出你所提及的奥力振六千卷书；但你却指望我熟悉毕达哥拉斯的所有记录。你鼓着胀起的双颊脱口而出的那些夸夸其谈，有什么真理可言？你声称通过引入你在奥力振其他书中读到的词句，已经校订了\OriginalTerm{论首要原理}{\GreekText{Περὶ} \GreekText{᾿Αρχῶν}}，因而你并未添加他人的话，而是恢复了他自己的话。在他著作的整片森林中，你连一根枝或一棵蘖都拿不出来。你指责我扬起烟雾和迷雾。这确实是烟雾和迷雾。你知道我已经驱散并消除了它们；但即使你的脖子断了，你也不肯低头，反而以超越你无知的厚颜无耻说我在否认显而易见的事，以便为自己开脱：你曾承诺金山，却连一个皮革般的铜板都拿不出。我承认你对我怀有敌意是有充分理由的，你的愤怒和激情是真实的；因为，若不是我坚持要求那些不存在的东西，你就会被认为拥有你所没有的东西。你向我索要毕达哥拉斯的书。但谁告诉你他的书还存在呢？确实，在你所批评的那封信中有这样的话：\Quote{假设我年轻时犯了错误，受过世俗文学的训练，在我基督徒生涯之初没有足够教义知识，并且我将我在毕达哥拉斯、柏拉图或恩培多克勒那里读到的内容归于宗徒；}然而我谈的是他们的学说，而不是他们的书，我能够通过西塞罗、布鲁图斯和塞内加了解这些学说。请读一下为瓦提尼乌斯作的简短演讲以及其他提及秘密社团的演讲。翻阅西塞罗的对话录。搜遍意大利海岸（从前称为大希腊），你会发现那里的公共纪念碑上铭刻着毕达哥拉斯的各种学说，刻在铜器上。那些\WorkTitle{金言}是谁的？是毕达哥拉斯的；其中以概要形式包含了他所有的原则。扬布里柯为之写了注释，至少部分遵循了莫德拉图斯（一位极擅辞令的人）以及毕达哥拉斯门徒阿基普斯和吕西德斯的做法。这些人中，阿基普斯和吕西德斯在希腊（即底比斯）执教；他们如此完整地保留了老师的训诫，以至于他们用记忆代替了书籍。其中一条训诫是：\Quote{我们必须用任何方法抛弃，并用火、剑和各种方法从身体切除疾病，从灵魂切除无知，从肚腹切除奢侈，从城邦切除叛乱，从家庭切除不和，从一切事物切除过度。}毕达哥拉斯还有其他训诫，如下：\Quote{朋友之间一切共有。}\Quote{朋友是另一个自我。}\Quote{有两个时刻尤需注意：早晨和晚上——即我们将要做的事和我们已经做的事。}\Quote{在天主之后，我们必须崇拜真理，因为只有真理使人与天主相似。}还有亚里士多德在其著作中勤奋收集的谜语：\Quote{永远不要越过天平杆，}即\Quote{不要违背正义的准则；}\Quote{绝不要用剑拨火，}即\Quote{当一个人愤怒激动时，不要用尖刻的话激怒他。}\Quote{不可触摸冠冕，}即\Quote{我们必须维护国家的法律。}\Quote{不要吞噬你的心，}即\Quote{从心中消除忧愁。}\Quote{一旦出发，就不要返回，}即\Quote{死后不要留恋此生。}\Quote{不要在大路上行走，}即\Quote{不要跟随群众的错误。}\Quote{绝不要让燕子进屋，}即\Quote{不要与饶舌多言之人同住一屋檐下。}\Quote{把新的担子加在已负重的人身上；不要加在放下担子的人身上；}即\Quote{当人们走向德行的道路时，用新的训诫来操练他们；当他们沉溺于怠惰时，就任凭他们。}我说我读过毕达哥拉斯学派的学说。让我告诉你，毕达哥拉斯是第一个发现灵魂不朽及其从一个身体转到另一个身体的人。维吉尔在\WorkTitle{埃涅阿斯记}第六卷中赞同这一观点，其言如下：\par

\Quote{这些灵魂，当千年轮转已满，\newline{}天主召叫，长长的一群，只愿在迷川中\newline{}淹没往昔，渴望再次看到\newline{}头上苍穹，重新回到肉身。}\par

40. 因此，\Person{毕达哥拉斯}曾教导，他本人原本是\Person{欧福尔布斯}，然后是\Person{卡利德斯}，第三是\Person{赫尔莫提摩斯}，第四是\Person{皮洛斯}，最后是\Person{毕达哥拉斯}；并且那些曾经存在的事物，在经过一定周期的循环后，会重新出现；因此世间不应认为有任何新事物。他说真正的哲学是对死亡的默想；其每日的奋斗是将灵魂从身体的牢笼中释放到自由之中；我们的学习即是回忆，还有许多其他内容，\Person{柏拉图}在其对话录中，特别是在\WorkTitle{斐多}和\WorkTitle{蒂迈欧}中详细阐述。因为\Person{柏拉图}在创立了学园并赢得了无数门徒之后，感到他的哲学在许多方面有欠缺，于是前往\Place{大希腊}，在那里从\Person{塔伦图姆的阿尔基塔斯}和\Person{洛克里的蒂迈欧}那里学习了\Person{毕达哥拉斯}的学说；并将这一体系以他从\Person{苏格拉底}学来的优雅形式和文风表现出来。我们可以证明，\Person{奥力振}将所有这些都搬进了他的著作\OriginalTerm{论原理}{\GreekText{Περὶ} \GreekText{᾿Αρχῶν}}，只是更换了名称。那么，当我曾说我在年轻时曾将我在\Person{毕达哥拉斯}、\Person{柏拉图}和\Person{恩培多克勒}那里发现的思想归之于宗徒们时，我犯了什么错误呢？我并非如你恶意推测的那样，说我曾读过\Person{毕达哥拉斯}、\Person{柏拉图}和\Person{恩培多克勒}的著作，而是说我曾读到过那些存在于他们著作中的东西，也就是说，其他人的著作向我表明这些东西曾存在于他们之中。这种表达方式很常见。例如，我可以说：\Quote{我在\Person{苏格拉底}那里读到的意见，我以为是真实的}，意思是我在\Person{柏拉图}和\Person{苏格拉底}学派的其他人的著作中读到作为他意见的东西，尽管\Person{苏格拉底}本人并未著书。同样，我可以说，我希望效仿我在\Person{亚历山大}和\Person{西庇阿}那里读到的功业，并非指他们描述了自己的功业，而是说我在他人的作品中读到了我钦佩的他们所做的功业。因此，即使我无法告知你任何\Person{毕达哥拉斯}本人尚存的记录，并有他儿子或女儿或其他门徒的证明，你也不能认为我犯了虚假之罪，因为我说的不是我读过他的书，而是他的学说。如果你以为这可以成为你谎言的掩护，并断言因为我不能提供任何\Person{毕达哥拉斯}所写的书，你就有权声称\Person{奥力振}的六千本书已经遗失，那你就大错特错了。\par

% structure: p0085 source=h2 target=subsection reason=relative-heading-rank; base=h2

\subsection{你以毁灭威胁我。我不会以同样的方式回应。论及信仰的争论中，应当避免人身攻击。}

41. 现在我来谈你的结语（即你对我倾倒的辱骂），在其中你劝我悔改，并以毁灭威胁我，除非我改变，也就是说，除非我在你的指控下保持沉默。你说，此丑事将落在我的头上，因为是我以答辩激起了你写作的疯狂，而你本是一个极其温和、具有\Person{梅瑟}般温良的人。你宣称你知道一些罪行，是我在你作为我最亲密的友人时单独向你承认的，并且你将把这些公之于众；我将被描绘成我本来的颜色；我应当记得我伏在你脚下，否则你可能会用你口中的剑砍下我的头。在诸如此类许多事情之后，你像个疯子一样翻来覆去，然后你挺直身板说，你希望和平，但仍然暗示我将来要保持安静，即我不得再写作反对异端者，也不得回应你提出的任何指控；如果我这样做，我将成为你的好兄弟和同僚，最雄辩的人，你的朋友和同伴；而且，更甚者，你将宣布我所有从\Person{奥力振}翻译的作品都是正统的。但是，如果我说一个字或移动一步，我立刻就会成为不健全和异端者，不配与你有任何联系。你就是这样吹嘘对我的赞美，你就是这样劝我和平。你连让我在悲伤中呻吟或流泪的自由都不给。\par

42. 我同样也可能将你描绘成你本来的颜色，以同样的狂怒回应你的疯狂；说出我所知道的，并加上我所不知道的；以你那样的放纵，或者说是暴怒和疯狂，真假混杂，直到我羞于开口，你羞于听闻；并以那种方式责备你，使得被指控者或指控者都难逃谴责；仅凭厚颜无耻将自己强加于读者，让他相信我所写的无耻之词乃是真实。但基督徒的实践远非如此，他们献出生命而非寻求他人的生命，他们不是用刀剑而是用意志成为杀人者。这或许符合你的温柔与天真；因为你可以从你胸中的粪堆里同时发出玫瑰的芳香和尸体的恶臭；并且，违背先知的诫命，称那曾经赞美为甘甜的为苦。但在处理基督徒话题时，我们无需抛出那些本应提交法庭的指控。你从我这里听到的不会比这句俗语更多：\Quote{当你说了你喜欢的话时，你将承受你不喜欢的事。}或者，如果这粗俗的谚语对你来说太过低俗，而作为有文化的人，你更喜欢哲学家或诗人的语句，那么请接受我引用\Person{荷马}的话。\par

\Quote{你说出什么话，你也会听到同样的话。}\par

有一件事，我倒想向一位如此杰出的圣德和苛求的人请教，（他的圣德大到连你们的手帕和围裙都能令魔鬼喊叫）；你以谁为你在写作中的榜样？在天主教作家中，有谁在意见争论时曾指控与他辩论之人的道德过失？你的师傅们教了你这样做吗？这就是你所受的训练制度，当你无法回答一个人时，就应该砍掉他的头？当你无法制止一个人的舌头时，就应该把它割掉？你没有什么值得夸耀的，因为你所做的不过是蝎子和斑蝥所做的。这就是\Person{富尔维亚}对\Person{西塞罗}所做的，以及\Person{黑落狄雅}对\Person{若翰}所做的。他们不能忍受听真理，因此用分开头发的簪子刺穿了说真理的舌头。狗的职责是在为主人服务时吠叫；我为什么不能在为基督服务时吠叫？许多人曾写作反对\Person{马西昂}或\Person{瓦伦提努斯}、\Person{亚流}或\Person{欧诺米乌斯}。其中有谁提出了不道德行为的指控？他们不是每次都全力以赴驳斥异端吗？这是异端者的阴谋，也就是你的师傅们的阴谋，当他们被定罪背叛信仰时，便求助于恶言毁谤。同样，\Place{安提约基雅}主教\Person{欧斯塔修斯}在不知不觉中成了父亲。同样，\Place{亚历山大}主教\Person{亚大纳削}切断了\Person{阿尔塞尼乌斯}的第三只手；因为当他被认为已死后又活着出现时，被发现只有两只手。如今同样的事也被诬告给同一个教会的主教，真正的信仰被金钱攻击，而金钱构成你和你的朋友们的势力。但我无需谈论与异端者的争论，他们虽然实质上是外人，却仍自称为基督徒。我们多少作家曾与那些最不虔敬的人，即\Person{切尔苏斯}和\Person{波斐利}争辩！但其中有谁放下了他所从事的事业，忙着将罪行归咎于他的对手，这类本应记载在法官的档案中而非教会著作中？如果你确定了某人的罪行却在辩论中失败，你得到了什么好处？你在提出控告时，根本不必冒上自己的脑袋。如果你的目标是报复，你可以雇一个刽子手，满足你的欲望。你假装害怕丑闻，却准备杀死一个曾经是你的兄弟、如今被你控告、并且一直被你当作敌人的人。然而我惊叹，像你这样一个知道自己在做什么的人，怎么会如此被疯狂蒙蔽，竟然想通过将我的灵魂从监狱中拉出来而给我好处，却不允许它与你一同留在这个世界的黑暗中。\par

% structure: p0090 source=h2 target=subsection reason=relative-heading-rank; base=h2

\subsection{和平之路是通过\BibleBook{}{箴言篇}所教导的智慧，以及通过信仰中的合一。}

43. 如果你希望我保持沉默，就停止控告我。放下你的剑，我便扔掉我的盾。只有一件事我不能同意，那就是宽恕异端者，而不为我的正统信仰辩护。如果那是我们之间不和谐的原因，我可以忍受死亡，但不能忍受沉默。原本应该通读整部\WorkTitle{圣经}作为对你狂言乱语的回答，并且像\Person{达味}弹奏竖琴一样，用神圣的话语来平复你愤怒的胸怀。但我只满足于从一本书中引用几句话；我将以智慧对抗愚昧；因为我希望，如果你轻视人的言语，你不会轻视天主的圣言。那么，请听智者\Person{撒罗满}关于你和一切喜好毁谤和辱骂的人所说的话：\par

\Quote{愚昧人渴望伤害，成为不虔敬的人，憎恶智慧。不要图谋恶事反对你的朋友。不要无故向人发怒。不虔敬的人高抬侮辱。从你身上除去邪恶的口，远离邪恶的嘴唇，说恶言的眼睛，不义者的舌头，流义人血的手，图谋恶念的心，急速行恶的脚。依靠谎言的人喂养风，追随飞翔的鸟。因为他离开了自己葡萄园的道路，使耕作的轮子误入歧途。他穿过干旱和荒芜的地方，用手收集荒芜。乖戾的口接近毁灭，出恶言的人是愚人的首领。每一个纯朴的人是有福的灵魂；但强暴的人是可耻的。因嘴唇的过失，罪人陷入罗网。愚人的一切道路在自己眼中看为正。愚人当日就显明他的怒气。说谎的嘴唇为主所憎恶。保守嘴唇的保全自己的灵魂；嘴唇急躁的使自己惊惶。恶人因强暴行恶，愚人铺张愚昧。在恶人中寻找智慧，你必找不到。急躁的人必吃自己道路的果实。智慧人因谨慎避开恶；愚人却自信，依附于恶。恒久忍耐的人在智慧上刚强；心小的人极其愚昧。欺压穷人的辱骂造他的主。智慧人的舌头知道善事，愚人的口却说恶。好争竞的人喜爱纷争，凡高抬自己心的在天主面前是不洁的。手连手行不义，他们必不能免受刑罚。爱惜生命的必须约束自己的口。傲慢在伤害之前，恶念在跌倒之前。闭着眼睛的说出悖逆的话，用嘴唇惹动一切恶。愚人的嘴唇引他入恶，冒失的言语招致死亡。恶谋的人必受大亏损。穷乏的义人胜过撒谎的富人。远离恶言是人的荣耀；愚昧的人却把自己束缚在其中。不要喜爱毁谤，免得你被除灭。谎言之饼对人甘甜，此后他的口却必充满沙砾。用说谎的舌头得财宝的，是追逐虚妄，必陷入死亡的网罗。不要在愚人耳边说话，免得智慧人嗤笑你的言语。棍棒、刀剑和箭都是伤害人的；同样，作假见证陷害朋友的人也是。如同鸟和麻雀飞去，咒诅也必落空，不得临到他。不要照愚昧人的愚昧回答他，免得你像他一样；却要照愚昧人的愚昧回答他，免得他自以为有智慧。埋伏等待朋友的人被揭露时说：我不过是开玩笑。木柴给炭火，木材给火焰，恶言的人给纷争的喧嚷。如果你的敌人向你求什么，吝啬却大声说话，你不要同意他，因为他心里有七级的邪恶。石头沉重，沙土难担；愚人的愤怒比这两样更重；愤恨残忍，怒气锋利，嫉妒急躁。不虔敬的人说穷人的坏话；心里狂妄的人极其愚昧。愚昧人尽发怒气，智慧人却部分地处理。恶子——他的牙齿是刀剑，他的臼齿如犁耙，要吞灭地上的软弱人和世人中的穷人。}\par

这些是我所受的教导，因此我不愿以牙还牙，以报复的方式攻击你；我认为更好的做法是驱除一个疯狂者的癫狂，并将单本书籍的解毒剂注入他中毒的胸膛。但我担心我将毫无成效，我将被迫唱\Person{达味}的诗歌，并以他的话为我唯一的安慰：\par

\BibleQuote{恶人一出母胎就与天主疏远，一离母腹便走迷路说谎话。他们的忿怒好像蛇的忿怒，像聋的虺蛇塞住耳朵，不听念咒者的声音和行邪术者巧妙的法术。天主必在牠们口中打碎他们的牙齿，主必打碎壮狮的大牙。他们必归无有，如水流去；祂拉紧弓弦，直到他们仆倒。他们如蜡熔化，被带走；火落在他们身上，他们不见太阳。}\BibleRef{咏 58}\par

又说：\par

\BibleQuote{义人见恶人遭报就欢喜，他要在罪人的血中洗手。人必说：果然有义人的报酬；果然有审判地上人的天主。}\BibleRef{咏 58}\par

44. 在信末你说：\Quote{我希望你喜爱和平。}对此我简答几句：你若渴求和平，就放下武器。我能与友善之人和平相处；我不畏惧恐吓我的人。让我们在信德上合而为一，和平便会立即相随。\par

\SourceRecord{Translated by W.H. Fremantle.；Nicene and Post-Nicene Fathers, Second Series；Vol. 3.；Edited by Philip Schaff and Henry Wace.；Buffalo, NY: Christian Literature Publishing Co.,；1892.；<http://www.newadvent.org/fathers/27103.htm>.}
